\documentclass[8pt,a4paper,onesided]{article}
\usepackage{multicol}
\usepackage[utf8]{inputenc}
\usepackage[english]{babel}
\usepackage{listings}
\usepackage[usenames,dvipsnames]{color}
\usepackage{verbatim}
\usepackage{hyperref}
\usepackage{color}
\usepackage{geometry}
\usepackage{minted}
\usepackage{amsmath, amssymb, amsfonts}
\usepackage{enumitem}
\usepackage[compact]{titlesec}
\usepackage[top=0.9cm, bottom=0.5cm, left=0.5cm, right=0.5cm, includehead]{geometry}
\usepackage{enumitem}
\setlist{nosep}
\makeatletter
\renewcommand{\subsection}{\@startsection{subsection}{2}{0mm}
    {-1ex plus -0.5ex minus -0.2ex} % Before-skip (Negative to pull it up)
    {0.5ex plus .2ex}                % After-skip (Small positive to separate text)
    {\normalfont\large\bfseries}}    % Font formatting
\makeatother
\setlength{\headsep}{3pt}
\geometry{verbose,portrait,a4paper,tmargin=0.9cm,bmargin=0.5cm,lmargin=0.5cm,rmargin=0.5cm}
\usepackage{listings}
\usepackage{color}
\setminted{ breaklines=true,
  breakatwhitespace=true,
  showstringspaces=false,
  columns=flexible,
  fontsize=\footnotesize,
  showtabs=false,
  style=default,
  breakanywhere=true,
  tabsize=1
}
\titlespacing{\subsection}{0pt}{0pt}{0pt}
\raggedbottom
\setlength{\parskip}{0pt}
\setlength{\columnsep}{0.1in}
\setlength{\columnseprule}{1px}
\usepackage{fancyhdr}
\pagestyle{fancy}
\fancyhf{}
\renewcommand{\headrulewidth}{0pt}
\fancyhead[R]{\thepage}
\fancyhead[L]{Shahjalal University of Science and Technology - GodFathers Inc.}
\setlist[enumerate]{leftmargin=17pt, itemsep=2pt, parsep=0pt}
\begin{document}
\title{\vspace{-1.5cm}GodFathers Inc. \\ \large \textit{sys.set\_int\_max\_str\_digits(10000); from heapq import *; heap = []; heappush(heap, ...); heapop(heap); } \\ \large \textit{list(map(int, input().strip().split()))} }
\author{Shahjalal University of Science and Technology}
\date{}
\maketitle
\begin{multicols}{3}
\tableofcontents
\end{multicols}
\vspace{0.5cm}
\begin{multicols}{3}
\lstloadlanguages{C++,Java}
\section{DP}
\subsection{Slope Trick}
\begin{minted}{c++}
//Maximize profit over N days by choosing to buy, sell, or hold one share daily based on known future prices. You must begin and end with zero inventory and are prohibited from selling shares you do not currently own.
priority_queue<int> q; q.push(a[1]); int ans = 0;
for (int i = 2; i <= n; i++) { q.push(a[i]);
if(q.top() > a[i]) { ans += q.top() - a[i]; q.pop(); q.push(a[i]);}}
\end{minted}
\subsection{Aliens Trick}
\begin{minted}{c++}
//choose at most k disjoint subarrays from array 'a' (-1e9 <= a_i <= 1e9) such that the sum of elements included in a subarray is maximized.

/* returns the maximum sum along with the number of subarrays used, if creating a subarray penalizes the sum by "lmb" and there is no limit to the number of subarrays you can create*/
auto solve_lambda = [&](ll lmb) {
 pair<ll, ll> dp[n][2];
 dp[0][0] = {0, 0}; dp[0][1] = {a[0] - lmb, 1};
 for (int i = 1; i < n; i++) {
  dp[i][0] = max(dp[i - 1][0], dp[i - 1][1]);
  dp[i][1] = max(make_pair(dp[i - 1][0].first + a[i] - lmb, dp[i - 1][0].second + 1),
    make_pair(dp[i - 1][1].first + a[i], dp[i - 1][1].second));}
 return max(dp[n - 1][0], dp[n - 1][1]);};
ll lo = 0; ll hi = 1e18;
while (lo < hi) {
 ll mid = (lo + hi + 1) / 2;
 solve_lambda(mid).second >= k ? lo = mid : hi = mid - 1;}
cout << solve_lambda(lo).first + lo * k << endl;
\end{minted}
\subsection{DNC}
\begin{minted}{c++}
// Divide 1,2,3...n people in k consecutive parts so that sum of cost of each individual part is minimum
int a[N][N], c[N][N], dp[810][N]; // dp[i][j]=minimum cost for dividing [1...j] into i parts
int cost(int i, int j) {
 // assert(i <= j);
 return c[j][j] - c[i - 1][j] - c[j][i - 1] + c[i - 1][i - 1];}
void yo(int i, int l, int r, int optl, int optr) {
 if(l > r) return;
 int mid = (l + r) / 2;
 dp[i][mid] = inf; // for maximum cost change it to 0
 int opt = -1;
 for(int k = optl; k <= min(mid - 1, optr); k++) {
  int c = dp[i - 1][k] + cost(k + 1, mid);
  if(c < dp[i][mid]) { // for maximum cost just change < to > only and rest of the algo should not be changed
   dp[i][mid] = c;
   opt = k;}}
 // for opt[1..j] <= opt[1...j+1]
 if (opt == -1) {
  // if we can't divide into k parts, then go right
  yo(i, mid + 1, r, optl, optr);
  return;}
 yo(i, l, mid - 1, optl, opt);
 yo(i, mid + 1, r, opt, optr);
 // for opt[1...j] >= opt[1...j+1]
 // yo(i, l, mid-1, opt, optr);
 // yo(i, mid+1, r, optl, opt);}
int main() {
 int i, j, k, n, m;
 n = sc(); k = sc();
 for(i = 1; i <= n; i++) for(j = 1; j <= n; j++) a[i][j] = sc();
 for(i = 1; i <= n; i++) {
  for(j = 1; j <= n; j++) {
   c[i][j] = a[i][j] + c[i - 1][j] + c[i][j - 1] - c[i - 1][j - 1];}}
  for(i = 1; i <= k; i++) {
    for (j = 1; j <= n; j++) {
      dp[i][j] = inf;}}
 for(i = 1; i <= n; i++) dp[1][i] = cost(1, i);
 for(i = 2; i <= k; i++) yo(i, 1, n, 1, n);
 cout << dp[k][n] / 2 << endl;
 return 0;}\end{minted}
\subsection{CHT}
\begin{minted}{c++}
struct CHT {
vector<ll> m, b; int ptr = 0;
bool bad(int l1, int l2, int l3) {
 return 1.0 * (b[l3] - b[l1]) * (m[l1] - m[l2]) <= 1.0 * (b[l2] - b[l1]) * (m[l1] - m[l3]); //(slope dec+query min),(slope inc+query max)
 return 1.0 * (b[l3] - b[l1]) * (m[l1] - m[l2]) > 1.0 * (b[l2] - b[l1]) * (m[l1] - m[l3]); //(slope dec+query max), (slope inc+query min)}
void add(ll _m, ll _b) {
 m.push_back(_m); b.push_back(_b);
 int s = m.size();
 while(s >= 3 && bad(s - 3, s - 2, s - 1)) {
  s--; m.erase(m.end() - 2);
  b.erase(b.end() - 2);}}
ll f(int i, ll x) {
 return m[i] * x + b[i];}
//(slope dec+query min), (slope inc+query max) -> x increasing
//(slope dec+query max), (slope inc+query min) -> x decreasing
ll query(ll x) {
 if(ptr >= m.size()) ptr = m.size() - 1;
 while(ptr < m.size() - 1 && f(ptr + 1, x) < f(ptr, x)) ptr++;
 return f(ptr, x);}
ll bs(int l, int r, ll x) {
 int mid = (l + r) / 2;
 if(mid + 1 < m.size() && f(mid + 1, x) < f(mid, x)) return bs(mid + 1, r, x); // > for max
 if(mid - 1 >= 0 && f(mid - 1, x) < f(mid, x)) return bs(l, mid - 1, x); // > for max
 return f(mid, x);}};
\end{minted}
\subsection{Dynamic CHT}
\begin{minted}{c++}
//add lines with -m and -b and return -ans to
//make this code work for minimums.(not -x)
const ll inf = -(1LL << 62);
struct line { ll m, b;
 mutable function<const line*() > succ;
 bool operator < (const line& rhs) const {
  if (rhs.b != inf) return m < rhs.m;
  const line* s = succ();
  if (!s) return 0;
  ll x = rhs.m;
  return b - s->b < (s->m - m) * x;}};
struct CHT : public multiset<line> {
 bool bad(iterator y) {
  auto z = next(y);
  if (y == begin()) {
   if (z == end()) return 0; return y -> m == z -> m && y -> b <= z -> b;}
  auto x = prev(y);
  if (z == end()) return y -> m == x -> m && y -> b <= x -> b;
  return 1.0 * (x -> b - y -> b) * (z -> m - y -> m) >= 1.0 * (y -> b - z -> b) * (y -> m - x -> m);}
 void add(ll m, ll b) {
  auto y = insert({ m, b });
  y->succ = [=, this] { return next(y) == end() ? 0 : &*next(y); };
  if (bad(y)) {erase(y); return;}
  while (next(y) != end() && bad(next(y))) erase(next(y));
  while (y != begin() && bad(prev(y))) erase(prev(y));}
 ll query(ll x) {
  assert(!empty());
  auto l = *lower_bound((line) {x, inf});
  return l.m * x + l.b;}};
CHT* cht;
\end{minted}
\section{DS and Misc}
\subsection{BIT}
\begin{minted}{c++}
int query(int i) {
int ans = 0;
for (; i >= 1; i -= (i & -i)) ans += t[i];
return ans;}
void upd(int i, int val) {
if (i <= 0) return;
for (; i <= n; i += (i & -i)) t[i] += val;}
\end{minted}
\subsection{Link Cut Tree}
\begin{minted}{c++}
struct Node {
 int x; Node *l = 0;
 Node *r = 0; Node *p = 0;
 bool rev = false; Node() = default;
 Node(int v) { x = v; }
 void push() {
  if (rev) {
   rev = false; swap(l, r);
   if (l) l->rev ^= true;
   if (r) r->rev ^= true;}}
 bool is_root() { return p == 0 || (p->l != this && this != p->r); }};
struct LCT {
 vector<Node> a; LCT(int n) {
  a.resize(n + 1);
  for (int i = 1; i <= n; ++i) a[i].x = i;}
 void rot(Node *c) {
  auto p = c->p;
  auto g = p->p;  if (!p->is_root()) (g->r == p ? g->r : g->l) = c;  p->push();
  c->push();  if (p->l == c) { // rtr
   p->l = c->r; c->r = p;
   if (p->l) p->l->p = p;} else { // rtl
   p->r = c->l; c->l = p;
   if (p->r) p->r->p = p;}  p->p = c;
  c->p = g;} void splay(Node *c) {
  while (!c->is_root()) {
   auto p = c->p; auto g = p->p;
   if (!p->is_root()) rot((g->r == p) == (p->r == c) ? p : c); rot(c);}
  c->push();} Node *access(int v) {
  Node *last = 0; Node *c = &a[v];
  for (Node *p = c; p; p = p->p) {
   splay(p); p->r = last; last = p;}
  splay(c); return last;} void make_root(int v) {
  access(v); auto *c = &a[v];
  if (c->l) c->l->rev ^= true, c->l = 0;} void link(int u, int v) {
  make_root(v); Node *c = &a[v];
  c->p = &a[u];} void cut(int u, int v) {
  make_root(u); access(v);
  if (a[v].l) {a[v].l->p = 0; a[v].l = 0;}}
 bool connected(int u, int v) {
  access(u); access(v);
  return a[u].p;}};
\end{minted}
\subsection{Min Plus Convolution}
\begin{minted}{c++}
//convex and arbitary
using poly = std::vector<int>;
const int inf = int(2e9) + 5;
poly solve(poly a,poly b) {
 int n = a.size(), m = b.size();
 int z = n + m - 1;
 poly c(z, inf), pos(z + 1);
 c[0] = a[0] + b[0];
 pos.back() = m - 1;
 int d = 1; while(d < z) d <<= 1;
 for(int len = d >> 1; len; len >>= 1) {
  for(int i = len; i < z; i += (len << 1)) {
   int l = i - len, r = std::min(z, i + len);
   pos[i] = pos[l]; for(int t = pos[l]; t <= pos[r]; ++t) {
    if(t <= i && i - t < n && c[i] > b[t] + a[i - t]) {
     c[i] = b[t] + a[i - t], pos[i] = t;}}}}
\end{minted}
\subsection{MOs with Update}
\begin{minted}{c++}
const int N = 2e5 + 9, B = 2500;
struct query {
 int l, r, t, id;
 bool operator < (const query &x) const {
 if(l / B == x.l / B) {
  if(r / B == x.r / B) return t < x.t;
  return r / B < x.r / B;}
 return l / B < x.l / B;}} Q[N];
struct upd {
 int pos, old, cur;} U[N];
int a[N];
int cnt[N], f[N], ans[N], l, r, t;
inline void add(int x) {
 f[cnt[x]]--, ++cnt[x], f[cnt[x]]++;}
inline void del(int x) {
 f[cnt[x]]--, --cnt[x], f[cnt[x]]++;}
inline void update(int pos, int x) {
 if (l <= pos && pos <= r) { add(x); del(a[pos]);} a[pos] = x;}
map <int, int> mp;
int nxt = 0;int get(int x) {return mp.count(x) ? mp[x] : mp[x] = ++nxt;}
int main() {
 int n, q; cin >> n >> q;
 for (int i = 1; i <= n; i++) {
 cin >> a[i]; a[i] = get(a[i]);}
 int nq = 0, nu = 0;
 for (int i = 1; i <= q; i++) {
 int ty, l, r; cin >> ty >> l >> r;
 if (ty == 1) ++nq, Q[nq] = {l, r, nu, nq};
 else ++nu, U[nu].pos = l, U[nu].old = a[l], a[l] = get(r), U[nu].cur = a[l];}
 sort(Q + 1, Q + nq + 1);
 t = nu, l = 1, r = 0;
 for (int i = 1; i <= nq; i++) {
 int L = Q[i].l, R = Q[i].r, T = Q[i].t;
 while(t < T) t++, update(U[t].pos, U[t].cur);
while(t > T) update(U[t].pos, U[t].old), t--;
 if(R < l) {
  while(l > L) add(a[--l]);
  while(l < L) del(a[l++]);
  while(r < R) add(a[++r]);
  while(r > R) del(a[r--]);
 } else {
  while(r < R) add(a[++r]);
  while(r > R) del(a[r--]);
  while(l > L) add(a[--l]);
  while(l < L) del(a[l++]);}
 int cur = 1; while(f[cur] > 0) cur++;
 ans[Q[i].id] = cur;}
 for (int i = 1; i <= nq; i++) cout << ans[i] << '\n';}
\end{minted}
\subsection{PBDS}
\begin{minted}{c++}
struct chash {
  const uint64_t C = 2713661325311;
  const uint32_t RANDOM = rng();
  size_t operator()(uint64_t x) const {
      return __builtin_bswap64((x ^ RANDOM) * C); }
  size_t operator()(pair<uint64_t, uint64_t> x) const { 
    return x.first * 31 + x.second;}};
#pragma comment(linker, "/stack:200000000")
#pragma GCC optimize("unroll-loops,-O3,Ofast,fast-math,no-stack-protector")
#pragma GCC target("avx2,bmi2,lzcnt,ssse3,sse4,popcnt,abm,mmx,tune=native")
#include <ext/pb_ds/assoc_container.hpp>
#include <ext/pb_ds/tree_policy.hpp>
using namespace __gnu_pbds;
template <typename T> using o_set = tree<T, null_type, less_equal<T>, rb_tree_tag, tree_order_statistics_node_update>;
mt19937 rng(chrono::steady_clock::now().time_since_epoch().count());
\end{minted}
\subsection{Pers Segment Tree}
\begin{minted}{c++}
#define lc t[cur].l
#define rc t[cur].r
struct node {
  int l = 0, r = 0, val = 0;
} t[20 * N];
int T = 0;
int build(int b, int e) {
  int cur = ++T; if(b == e) return cur;
  int mid = b + e >> 1;
  lc = build(b, mid); rc = build(mid + 1, e);
  t[cur].val = t[lc].val + t[rc].val;
  return cur;
}
int upd(int pre, int b, int e, int i, int v) {
  int cur = ++T; t[cur] = t[pre];
  if(b == e) {t[cur].val += v; return cur;}
  int mid = b + e >> 1;
  if(i <= mid) lc = upd(t[pre].l, b, mid, i, v);
  else rc = upd(t[pre].r, mid + 1, e, i, v);
  t[cur].val = t[lc].val + t[rc].val; return cur;
}
  int query(int cur, int b, int e, int i, int j) {
  if(b > j || e < i) return 0;
  if(i <= b && e <= j) {return t[cur].val;}
  int mid = b + e >> 1;
  return query(lc, b, mid, i, j) + query(rc, mid + 1, e, i, j);
}
\end{minted}
\subsection{SOS}
\begin{minted}{c++}
 // sum over subsets
 for (int i = 0; i < B; i++) {
  for (int mask = 0; mask < (1 << B); mask++) {
   if ((mask & (1 << i)) != 0) {
    f[mask] += f[mask ^ (1 << i)];}}}
 // sum over supersets
 for (int i = 0; i < B; i++) {
  for (int mask = (1 << B) - 1 ; mask >= 0 ; mask--) {
   if ((mask & (1 << i)) == 0) g[mask] += g[mask ^ (1 << i)] ;}}
\end{minted}
\subsection{Sparse Table}
\begin{minted}{c++}
#define LVL 20 //ceil(log_2(MAXN))+1
int arr[MAX], dst[MAX][LVL];
int merge(int x, int y) { return __gcd(x,y); }
//0 based
void build(int n){
  for(int i=0; i<n; ++i) dst[i][0] = arr[i];
  for(int k=1; k<LVL; ++k)
    for(int i=0; (i+(1<<k)-1)<n; ++i)
      dst[i][k]=merge(dst[i][k-1], dst[i+(1<<(k-1))][k-1]);}
int query(int l, int r){
  int k=31-__builtin_clz(r-l+1);
  return merge(dst[l][k], dst[r-(1<<k)+1][k]);  }
vector<vector<vector<int>>> st;
int merge(int a,int b,int c,int d) {
    return min({a, b, c, d}); }
//O(n*m*Log(max(n,m)) ///0-based
void build2DSparseTable(int n,int m){
  int LOG = log2(max(n, m))+1;
  vector<vector<vector<int>>>().swap(st);
  st.resize(n, vector<vector<int>>(m, vector<int>(LOG)));
  for(int i = 0; i < n; i++)
    for(int j = 0; j < m; j++)
      st[i][j][0] = arr[i][j];
  for(int k = 1; k < LOG; k++) {
  for(int i = 0; i + (1<<k) <= n; i++) {
  for(int j = 0; j + (1<<k) <= m; j++) {
    st[i][j][k] = merge( st[i][j][k-1],
      st[i+(1<<(k-1))][j][k-1],
      st[i][j+(1<<(k-1))][k-1],
      st[i+(1<<(k-1))][j+(1<<(k-1))][k-1]);}}}}
int query(int x1,int y1,int x2,int y2) {
  int k = log2(min(x2-x1+1, y2-y1+1))+1;
  return merge( st[x1][y1][k-1],
    st[x2-(1<<(k-1))+1][y1][k-1],
    st[x1][y2-(1<<(k-1))+1][k-1],
    st[x2-(1<<(k-1))+1][y2-(1<<(k-1))+1][k-1]);}
\end{minted}
\subsection{VS Code Build}
\begin{minted}{c++}
{"tasks": [{"type": "process","label": "CP build","command": "bash",
 "args": ["-c",
 "ulimit -s262144 -f102400 -t5; g++ -O2 -g -fsanitize=address,undefined -fno-omit-frame-pointer -std=c++17 \"${file}\" -o /tmp/a.out && /usr/bin/time -f \"%esec | %MKB\" /tmp/a.out <in.txt >out.txt"],
 "problemMatcher":
 {"owner": "cpbuild", "fileLocation": "absolute",
 "pattern": [{
 "regexp": "^(.*):([0-9]+):([0-9]+):\\s+(.*error):\\s+(.*)$",
 "file": 1, "line": 2, "column": 3, "severity": 4, "message": 5}]},
 "options": {"cwd": "${fileDirname}"},
 "group": {"kind": "build", "isDefault": true},
 "presentation": {"showReuseMessage": false, "clear": true, "revealProblems": "onProblem"}}
 ], "version": "2.0.0"}
\end{minted}
\subsection{Submasks 3^n}
\begin{minted}{c++}
int total_masks = (1 << n);
for (int mask = 0; mask < (1 << n); ++mask) {
  for (int submask = mask; submask > 0; submask = (submask - 1) & mask) {}}
\end{minted}
\subsection{Treap set}
\begin{minted}{c++}
template <typename T> struct treap {
  struct node {
    node *l, *r; int prior, sz;
    T key; node(T id) { l = r = nullptr; key = id; prior = rnd(); sz = 1;}};
  node *root;
  treap() {
    root = nullptr;}
  int get_sz (node *x) {if (x == nullptr) return 0;
    return x->sz;}
  void resize(node *x) { x->sz = get_sz(x->l) + get_sz(x->r) + 1;}
  void split(node *t, T pos, node *&l, node *&r) {
    if (t == nullptr) {
      l = r = nullptr; return;}
    if (t->key <= pos) {
      split(t->r, pos, l, r);
      t->r = l; l = t;}
    else { split(t->l, pos, l, r);
      t->l = r; r = t;}
    resize(t);}
  node* merge(node *l, node *r) {
    if (!l || !r) return l ? l : r;
    if (l->prior < r->prior) {
      l->r = merge(l->r, r);
      resize(l);return l;}
    r->l = merge(l, r->l);
    resize(r); return r;}
  node* merge_treap(node *l, node *r) {
    if(!l) return r; if(!r) return l;
    if(l->prior < r->prior) swap(l, r);
    node *L, *R; split(r, l->key, L, R);
    l->r = merge_treap(l->r, R);
    l->l = merge_treap(L, l->l);
    resize(l); return l;}
  void insert(T id) {
    node *l, *r; split(root, id, l, r);
    root = merge(merge(l, new node(id)), r);}
  node* erase(T L, T R) {
    node *l, *r, *mid, *mr;
    split(root, L - 1, l, r);
    split(r, R, mid, mr);
    root = merge(l, mr); return mid;}
  void combine(node *x) { //combine with another treap
    root = merge_treap(root, x);}
  vector<T> ans;
  void dfs(node * cur) {
    if (!cur) return;
    ans.push_back(cur -> key);
    dfs(cur -> l); dfs(cur -> r);}
  vector<T> get() {
    ans.clear(); dfs(root); return ans;}
  int calc_sz (node * t, T x) { if (t == nullptr) return 0;
    if (t->key < x) return get_sz(t->l) + 1 + calc_sz(t->r, x);
    return calc_sz(t->l, x);}
  node* fnd (node* t, int ord) { if (get_sz(t) <= 1) return t;
    if (get_sz(t->l) >= ord) return fnd(t->l, ord);
    if (get_sz(t->l) + 1 == ord) return t;
    return fnd(t->r, ord - get_sz(t->l) - 1);}
  int order_of_key (T x) { return calc_sz(root, x);}
  node* find_by_order (int ord) {if (ord >= get_sz(root) || ord < 0) return nullptr;
    return fnd(root, ord+1);}}
\end{minted}
\subsection{Treap}
\begin{minted}{c++}
//split by size, insert at any point, binary search using lower_bound, reverse using lazy
template <typename T> struct treap {
  struct node {
    node *l, *r, *par; int prior, sz; T key; bool laz;
    node(T id) { l = r = par = nullptr; key = id; prior = rnd(); sz = 1, laz = 0;}};
  node *root;
  treap() {root = nullptr;}
    int get_sz (node *x) {
    if (x == nullptr) return 0; return x->sz;}
  void propagate (node *t) {
    if (t->laz == 0) return;
    if (t->l) {t->l->laz ^= t->laz;}
    if (t->r) {t->r->laz ^= t->laz;}
    t->laz = 0;
    swap(t->l, t->r);  }
    void resize(node *x) { if (x == nullptr) return;
    x->sz = get_sz(x->l) + get_sz(x->r) + 1;
    if (x->l != nullptr) x->l->par = x;
    if (x->r != nullptr) x->r->par = x;
    x->par = nullptr;}
  void split(node *t, int pos, node *&l, node *&r) {
    if (t == nullptr) {
      l = r = nullptr; return;}
    propagate(t);
    if (get_sz(t->l) < pos) {
      split(t->r, pos - get_sz(t->l) - 1, l, r);
      t->r = l; l = t;}
    else {
      split(t->l, pos, l, r);
      t->l = r; r = t;}
    resize(t);}
    node* merge(node *l, node *r) {
    if (!l || !r) (l ? resize(l) : resize(r));
    if (!l || !r) return l? l : r;
    propagate(l); propagate(r);
    if (l->prior < r->prior) {
      l->r = merge(l->r, r);
      resize(l); return l;}
    r->l = merge(l, r->l);
    resize(r); return r;}
    node* insert(int pos, T id) {
    if (pos < 0 || pos > get_sz(root)) return nullptr;
    node *l, *r; split(root, pos, l, r); node* newnode = new node(id);
    root = merge(merge(l, newnode), r);
    return newnode;}  node* insert(int pos, node* id) {
    if (pos < 0 || pos > get_sz(root)) return nullptr;
    node *l, *r; split(root, pos, l, r);
    node* newnode = id; root = merge(merge(l, newnode), r);
    return newnode;}
  node* push_back(T id) {
    node* newnode = new node(id); root = merge(root, newnode);
    return newnode;}
  node* erase(int L, int R) {
    node *l, *r, *mid, *mr;
    split(root, L, l, r); split(r, R-L+1, mid, mr); root = merge(l, mr); return mid;} 
  vector<T> ans;
  void dfs(node * cur) {
    if (!cur) return; dfs(cur -> l); ans.push_back(cur -> key); dfs(cur -> r);}
  vector<T> get() {
    ans.clear(); dfs(root); return ans;}
  node* fnd (node* t, int ord) {
    if (get_sz(t) <= 1) return t;
    propagate(t);
    if (get_sz(t->l) >= ord) return fnd(t->l, ord);
    if (get_sz(t->l) + 1 == ord) return t;
    return fnd(t->r, ord - get_sz(t->l) - 1);}
  node* at (int ord) {
    if (ord >= get_sz(root) || ord < 0) return nullptr;
    return fnd(root, ord+1);}
  int order (node* x) {
    if (x == nullptr) return get_sz(root);
    int ans = get_sz(x->l);
    while (x != nullptr) { node* nx = x->par; if (nx && nx->r == x) ans += get_sz(nx->l) + 1; x = nx;}
    return ans;}
  node* lower_bound (const auto &x, auto cmp) {
    node* v = root, *res = nullptr;
    while (v != nullptr) {
      if (cmp(v->key, x)) { v = v->r;}
      else { res = v; v = v->l;}} return res;}
  void update (node *t, int L, int R) {
    node *l, *r, *mid, *mr;
    split(root, L, l, r); split(r, R-L+1, mid, mr);
    mid->laz ^= 1; root = merge(merge(l, mid), mr);}};
\end{minted}
\subsection{WavLet Tree}
\begin{minted}{c++}
#define pii pair<int,int>
#define x first #define y second
int n;
struct wavelet_tree {
  static const int unfound = INT_MIN;
  static const int maxn = 1e5 + 5;
  static const int mlog = 20;
  typedef int Int; typedef long long Long;
  Int s[maxn], tree[mlog][maxn];
  int L[mlog][maxn];  Long ls[mlog][maxn], sl;
  Int & operator[](int x) {return tree[0][x];}
  void build(int l = 1, int r = n, int d = 0) {
    if (l == r)return;
    int m=(l + r)>>1, cnt=0, lc=l, rc=m+1, ec=0;
    for(int i=l; i<=r; i++) if(tree[d][i]<s[m])cnt++;
    for (int i = l; i <= r; i++) {
      if ( (tree[d][i] < s[m]) || (tree[d][i] == s[m] && ec < (m - l + 1 - cnt) ) ) {
        tree[d + 1][lc++] = tree[d][i];
        ls[d + 1][i] = ls[d+1][i-1]+tree[d][i];
        if (tree[d][i] == s[m]) ec++;} else {
        tree[d + 1][rc++] = tree[d][i];
        ls[d + 1][i] = ls[d + 1][i - 1];} L[d][i] = L[d][l - 1] + lc - l;} build(l, m, d + 1); build(m + 1, r, d + 1); }
  void init(Int *arr, int n) {
    for(int i=1; i<=n; i++) tree[0][i] = arr[i]; init(n); }
  void init(int n) {
    for(int i=1; i<=n; i++)s[i]=tree[0][i], ls[0][i]=ls[0][i-1]+s[i];
    sort(s + 1, s + 1 + n); build(); }
  Long sum(pii a,int d=0){return ls[d][a.y]-ls[d][a.x-1];}
  int cn(pii a, int d) {return L[d][a.y] - L[d][a.x - 1];}
  pii left(pii a, int d, int l) {return {l + cn({l, a.x-1}, d), l-1+cn({l, a.y}, d)};}
  pii right(pii a, int d, int r) {return {a.x + cn({a.x, r}, d), a.y + cn({a.y+1, r}, d)};}
  Int kth(int x,int y,int k,int l=1, int r=n, int d=0) {
  if (y - x + 1 < k || x > y)return unfound;
    if (l == r) { sl += tree[d][l]; return tree[d][l];} int cnt = cn({x, y}, d), m = (l + r) >> 1, nx, ny;
    if (cnt >= k) { tie(nx, ny) = left({x, y}, d, l);
      return kth(nx, ny, k, l, m, d + 1);} else {
      sl += sum({x, y}, d+1); tie(nx, ny)=right({x,y},d,r);
      return kth(nx, ny, k - cnt, m + 1, r, d + 1); } }
  int leq(int x, int y, Int k, int l=1, int r=n, int d=0){
    if (x > y)return 0; if (l == r) {
      if (l > y || l < x)return 0; // is it important?
      sl += tree[d][l] * (tree[d][l] <= k);
      return tree[d][l] <= k;} int cnt = cn({x, y}, d), m = (l + r) >> 1, nx, ny;
    if (s[m] <= k) {
      sl += sum({x,y}, d+1); tie(nx,ny) = right({x,y},d,r);
      return cnt + leq(nx, ny, k, m + 1, r, d + 1);} else {  tie(nx, ny) = left({x, y}, d, l);
      return leq(nx, ny, k, l, m, d + 1); } }
  Int rmin(int x, int y, int l=1, int r=n, int d=0) {return kth(x, y, 1, l, r, d);}
  Int rmax(int x, int y, int l=1, int r=n, int d=0) {return kth(x, y, y-x+1, l, r, d);}
  Int floor(int x, int y, Int k, int l = 1, int r = n, int d = 0) {
    if (x > y)return INT_MIN;
    if (l == r) {
      if (l > y || l < x)return INT_MIN;
      return (tree[d][l] <= k ? tree[d][l] : INT_MIN);} Int ans = INT_MIN; int m = (l + r) >> 1, nx, ny;
    if (s[m] <= k) { tie(nx, ny) = right({x, y}, d, r);
      ans = max(ans, floor(nx, ny, k, m + 1, r, d + 1));
      if (ans == INT_MIN) {
        tie(nx, ny) = left({x, y}, d, l);
        auto an = rmax(nx, ny, l, m, d + 1);
        if (an != unfound)ans = max(ans, an);}} else { tie(nx, ny) = left({x, y}, d, l);
        ans = max(ans, floor(nx, ny, k, l, m, d + 1));} return ans; }
  Int ceil(int x, int y, Int k, int l=1, int r=n, int d=0){
    if (l == r) {
      if (l > y || l < x)return INT_MAX;
      return (tree[d][l] >= k ? tree[d][l] : INT_MAX);} Int ans = INT_MAX; int m = (l + r) >> 1, nx, ny;
    if (s[m] >= k) { tie(nx, ny) = left({x, y}, d, l);
      ans = min(ans, ceil(nx, ny, k, l, m, d + 1));
      if (ans == INT_MAX) {
        tie(nx, ny) = right({x, y}, d, r);
        auto an = rmin(nx, ny, m + 1, r, d + 1);
        if (an != unfound)ans = min(ans, an);}} else { tie(nx, ny) = right({x, y}, d, r);
        ans = min(ans, ceil(nx, ny, k, m + 1, r, d + 1));} return ans;  } } wvt;
\end{minted}
\subsection{Xor Basis}
\begin{minted}{c++}
struct XorBasis{
  vector<ll> basis; ll N=0,tmp=0;
  void add(ll x){ N++; tmp|=x;
    for(auto &i : basis) x=min(x,x^i); if(!x) return;
    for(auto &i : basis) if((i^x)<i) i^=x;
    basis.push_back(x); sort(basis.begin(),basis.end());}
  ll size(){ return (ll)basis.size(); }
  void clear(){ N=0;tmp=0; basis.clear(); }
  bool possible(ll x){
    for(auto &i: basis) x=min(x, x^i); return !x; }
  ll maxxor(ll x=0){
    for(auto &i: basis) x=max(x, x^i); return x; }
  ll minxor(ll x=0){
    for(auto &i: basis) x=min(x, x^i); return x; }
  ll cntxor(ll x){
    if(!possible(x)) return 0; return (1LL<<(N-size()));}
  ll sumOfAll(){
    ll ans=tmp*(1LL<<(N-1)); return ans; }
  ll kth(ll k){ ll sz=size(); if(k > (1LL<<sz)) return -1;
    k--; ll ans=0;
    for(ll i=0; i<sz; i++) if(k>>i & 1) ans^=basis[i];
    return ans;}};
\end{minted}
\subsection{Josephus}
\begin{minted}{c++}
// n = total person. will kill every kth person, if k = 2, 2,4,6.. returns the mth killed person 
// O(k log n) 
int josephus(int n, int k, int m) {
 m = n - m; int i = m;
 if (k <= 1) return n - m;
 while (i < n) {
  int r = (i - m + k - 2) / (k - 1);
  if ((i + r) > n) r = n - i;
  else if (!r) r = 1;
  i += r; m = (m + (r * k)) % i;}
 return m + 1;}
\end{minted}
\subsection{Power Tower Mod}
\begin{minted}{c++}
int mod(int x){if(x<m) return x; return x%m + m;}
\end{minted}
\section{Geometry}
\subsection{2D Full Template}
\begin{minted}{c++}
typedef long double ld; //typedef long long T;
typedef long double T; #define PI acos(-1.0)
#define eps 1e-9 #define inf (1e100)
int sign(T x) {return (x>eps)-(x<-eps);}
struct PT{
 T x,y; PT(){x=0,y=0;}
 PT(T x,T y) : x(x),y(y){}
 PT(const PT& p) : x(p.x), y(p.y) {}
 PT operator + (const PT &a) const { return PT(x + a.x, y + a.y); }
 PT operator - (const PT &a) const { return PT(x - a.x, y - a.y); }
 PT operator * (const T a) const { return PT(x * a, y * a); }
 friend PT operator * (const T &a, const PT &b) { return PT(a * b.x, a * b.y); }
 PT operator / (const T a) const { return PT(x / a, y / a); }
 PT operator*(PT p) {return {x*p.x-y*p.y, x*p.y+y*p.x};}
 PT operator/(PT p) {return {((*this)*p.conj())/(p*p.conj()).x};}
 bool operator == (PT a) const { return sign(a.x - x) == 0 && sign(a.y - y) == 0; }
 bool operator != (PT a) const { return !(*this == a); }
 bool operator < (PT a) const { return sign(a.x - x) == 0 ? y < a.y : x < a.x; }
 bool operator > (PT a) const { return sign(a.x - x) == 0 ? y > a.y : x > a.x; }
 ld norm() { return sqrt(x * x + y * y); }
 T norm2() { return x * x + y * y; }//squared dis
 PT perp() { return PT(-y, x); }
 PT unit() {return (*this)/norm();}
 PT conj() {return {x,-y};}
 ld arg() { return atan2(y, x); }
 //a vector with norm r and having same direction
 PT truncate(T r) { 
  ld k = norm(); if(!sign(k)) return *this;
  r /= k; return PT(x * r, y * r); }};
typedef vector<PT> VPT;
T dot(PT a, PT b) { return a.x * b.x + a.y * b.y; }
T dist2(PT a, PT b) { return dot(a - b, a - b); }
ld dist(PT a, PT b) { return sqrt(dot(a - b, a - b)); }
T cross(PT a, PT b) { return a.x * b.y - a.y * b.x; }
T cross2(PT a, PT b, PT c) { return cross(b - a, c - a); }
int orientation(PT a, PT b, PT c) { return sign(cross(b - a, c - a)); }
PT perp(PT a) { return PT(-a.y, a.x); }
PT rotate_ccw90(PT a) { return PT(-a.y, a.x); }
PT rotate_cw90(PT a) { return PT(a.y, -a.x); }
PT rotate_ccw(PT a, T t) { return PT(a.x * cos(t) - a.y * sin(t), a.x * sin(t) + a.y * cos(t)); }
PT rotate_cw(PT a, T t) { return PT(a.x * cos(t) + a.y * sin(t), -a.x * sin(t) + a.y * cos(t)); }
//O er sapekke A ke radian kune CCW
PT rotate(PT O, PT A, T radian){ return O + rotate_ccw(A - O, radian); }
T SQ(T x) { return x * x; }
ld rad_to_deg(T r) { return (r * 180.0 / PI); }
ld deg_to_rad(T d) { return (d * PI / 180.0); }
T get_angle(PT a, PT b) {
  T costheta = dot(a, b) / a.norm() / b.norm();
  return acos(max((T)-1.0, min((T)1.0, costheta))); }
// does point p lie in angle <bac
bool is_point_in_angle(PT b, PT a, PT c, PT p) {
  assert(orientation(a, b, c) != 0);
  if (orientation(a, c, b) < 0) swap(b, c);
  return orientation(a, c, p) >= 0 && orientation(a, b, p) <= 0; }
bool half(PT p) {
  return p.y > 0.0 || (p.y == 0.0 && p.x < 0.0); }
// sort points in counterclockwise
void polar_sort(VPT &v) {
  sort(v.begin(), v.end(), [](PT a,PT b) {
    return make_tuple(half(a), (T)0.0, a.norm2()) < make_tuple(half(b), cross(a, b), b.norm2()); }); } 
struct line {
 PT a, b;// goes through points a and b
 PT v; T c;//line form: direction vec [cross](x,y)=c 
 line() {}
 //direction vector v and offset c
 line(PT v, T c) : v(v), c(c) {
  auto p = get_points(); a = p.first; b = p.second; }
 // equation ax + by + c = 0
 line(T _a, T _b, T _c) : v({_b, -_a}), c(-_c) {
   auto p = get_points();
   a = p.first; b = p.second; }
 //goes through points p and q
 line(PT p, PT q) : v(q - p), c(cross(v, p)), a(p), b(q) {}
  pair<PT, PT> get_points() {
  PT p, q; T a = -v.y, b = v.x; // ax + by = c
  if(sign(a) == 0) {
   p = PT(0, c / b); q = PT(1, c / b); }
  else if(sign(b) == 0) {
   p = PT(c / a, 0); q = PT(c / a, 1); }
  else {
   p = PT(0, c / b); q = PT(1, (c - a) / b); }
  return {p, q}; }
 array<T, 3> get_abc() {
  T a = -v.y, b = v.x; return {a, b, c}; }
 //1 if on the left,-1 if on the right,0 if on the line
 int side(PT p) { return sign(cross(v, p) - c); }
 line perpendicular_through(PT p){return {p,p+perp(v)};}
 line translate(PT t) { return {v, c + cross(v, t)}; }
 bool cmp_by_projection(PT p, PT q) { return dot(v, p) < dot(v, q); }
 line shift_left(T d) {
  PT z = v.perp().truncate(d);
  return line(a + z, b + z); }}
;
//find a point from A through B with distance d
PT point_along_line(PT a, PT b, T d) {
 return a + (((b - a) / (b - a).norm()) * d); }
PT project_from_point_to_line(PT a, PT b, PT c) {
 return a+(b-a)*dot(c-a,b-a)/(b - a).norm2(); }
PT reflection_from_point_to_line(PT a, PT b, PT c) {
 PT p = project_from_point_to_line(a,b,c);
 return point_along_line(c, p, (T)2 * dist(c, p)); }
ld dist_from_point_to_line(PT a, PT b, PT c) {
 return fabs(cross(b - a, c - a) / (b - a).norm()); }
bool is_point_on_seg(PT a, PT b, PT p) {
 if (abs(cross(p - b, a - b)) < eps) {
  if(p.x<min(a.x,b.x) || p.x>max(a.x,b.x)) return 0;
  if(p.y<min(a.y,b.y) || p.y>max(a.y,b.y)) return 0;
  return 1; } return 0; }
//minimum distance point from point c to segment ab that lies on segment ab
PT project_from_point_to_seg(PT a, PT b, PT c) {
 T r = dist2(a, b); if(abs(r) < eps) return a;
 r = dot(c - a, b - a) / r; if (r < 0) return a;
 if (r > 1) return b; return a + (b - a) * r; }
ld dist_from_point_to_seg(PT a, PT b, PT c) {
 return dist(c, project_from_point_to_seg(a, b, c));}
//0 if not parallel, 1 if parallel, 2 if collinear
int is_parallel(PT a, PT b, PT c, PT d) {
 T k = abs(cross(b - a, d - c));
 if (k < eps){
  if (abs(cross(a - b, a - c)) < eps && abs(cross(c - d, c - a)) < eps) return 2;
  else return 1;} else return 0; }
bool are_lines_same(PT a, PT b, PT c, PT d) {
 if(abs(cross(a - c, c - d)) < eps && abs(cross(b - c, c - d)) < eps) return true; return false; }
//bisector vector of <abc
PT angle_bisector(PT &a, PT &b, PT &c){
 PT p=a-b,q=c-b; return p+q*sqrt(dot(p,p)/dot(q,q)); }
//1 if point is ccw to the line, 2 if point is cw to the line, 3 if point is on the line
int point_line_relation(PT a, PT b, PT p) {
 int c = sign(cross(p - a, b - a)); if (c < 0) return 1;
 if (c > 0) return 2; return 3;}
bool line_line_intersection(PT a,PT b,PT c,PT d,PT &ans){
 T a1 = a.y - b.y, b1 = b.x - a.x, c1 = cross(a, b);
 T a2 = c.y - d.y, b2 = d.x - c.x, c2 = cross(c, d);
 T det=a1*b2-a2*b1; if(det == 0) return 0;
 ans=PT((b1*c2-b2*c1)/det,(c1*a2-a1*c2)/det); return 1;}
bool seg_seg_intersection(PT a,PT b,PT c,PT d,PT &ans){
 T oa = cross2(c, d, a), ob = cross2(c, d, b);
 T oc = cross2(a, b, c), od = cross2(a, b, d);
 if (oa * ob < 0 && oc * od < 0){
  ans=(a*ob-b*oa)/(ob-oa); return 1; } else return 0;}
//se.size()==0 means no intersection
//se.size()==1 means one intersection
//se.size()==2 means range intersection
set<PT> seg_seg_intersection_inside(PT a,PT b,PT c,PT d){
 PT ans; set<PT> se;
 if(seg_seg_intersection(a,b,c,d,ans)) return {ans};
 if (is_point_on_seg(c, d, a)) se.insert(a);
 if (is_point_on_seg(c, d, b)) se.insert(b);
 if (is_point_on_seg(a, b, c)) se.insert(c);
 if(is_point_on_seg(a,b,d)) se.insert(d); return se;}
//[ab],(cd). 0->not, 1->proper,2->segment intersect
int seg_line_relation(PT a, PT b, PT c, PT d) {
 T p = cross2(c, d, a); T q = cross2(c, d, b);
 if (sign(p) == 0 && sign(q) == 0) return 2;
 else if (p * q < 0) return 1; else return 0; }
bool seg_line_intersection(PT a,PT b,PT c,PT d,PT &ans){
 bool k=seg_line_relation(a, b, c, d); assert(k != 2);
 if(k) line_line_intersection(a,b,c,d,ans); return k; }
ld dist_from_seg_to_seg(PT a, PT b, PT c, PT d) {
 PT dummy;
 if(seg_seg_intersection(a,b,c,d,dummy)) return (T)0.0;
 else return min({dist_from_point_to_seg(a,b,c),
  dist_from_point_to_seg(a, b, d), 
  dist_from_point_to_seg(c, d, a),
  dist_from_point_to_seg(c, d, b)});}
//c to [a,b)
ld dist_from_point_to_ray(PT a, PT b, PT c) {
 b=a+b; T r=dot(c-a,b-a); if(r<0.0) return dist(c,a);
 return dist_from_point_to_line(a, b, c);}
//starting point as and direction vector ad
bool ray_ray_intersection(PT as, PT ad, PT bs, PT bd){
 T dx = bs.x - as.x, dy = bs.y - as.y;
 T det = bd.x * ad.y - bd.y * ad.x;
 if(abs(det)<eps) return 0;
 T u = (dy * bd.x - dx * bd.y) / det;
 T v = (dy * ad.x - dx * ad.y) / det;
 if(sign(u) >= 0 && sign(v) >= 0) return 1;
 else return 0; }
ld ray_ray_distance(PT as,PT ad,PT bs,PT bd){
 if(ray_ray_intersection(as,ad,bs,bd)) return 0.0;
 ld ans = dist_from_point_to_ray(as, ad, bs);
 ans = min(ans, dist_from_point_to_ray(bs,bd,as));
 return ans;}
struct circle {
 PT p; T r; circle() {}
 circle(PT _p, T _r): p(_p), r(_r) {};
 circle(T x, T y, T _r): p(PT(x, y)), r(_r) {};
 //circumcircle of a triangle
 circle(PT a, PT b, PT c) {
  b=(a+b)*0.5; c=(a+c)*0.5;
  line_line_intersection(b, b + rotate_cw90(a - b), c, c + rotate_cw90(a - c), p);
  r = dist(a, p); }
 //inscribed circle of a triangle
 circle(PT a, PT b, PT c, bool t) {
  line u, v; ld m = atan2(b.y-a.y, b.x-a.x);
  ld n = atan2(c.y-a.y, c.x-a.x); u.a = a;
  u.b=u.a+(PT(cos((n+m)/2.0),sin((n+m)/2.0))); v.a=b;
  m=atan2(a.y-b.y,a.x-b.x); n=atan2(c.y-b.y,c.x-b.x);
  v.b = v.a+(PT(cos((n + m)/2.0), sin((n + m)/2.0)));
  line_line_intersection(u.a, u.b, v.a, v.b, p);
  r = dist_from_point_to_seg(a, b, p); }
 bool operator==(circle v){ return p == v.p && sign(r - v.r) == 0; }
 ld area() { return PI * r * r; }
 ld circumference() { return 2.0 * PI * r; } };
//0->outside,1->on circumference, 2->inside circle
int circle_point_relation(PT p, T r, PT b) {
 ld d = dist(p,b); if(sign(d - r) < 0) return 2;
 if(sign(d - r) == 0) return 1; return 0; }
int circle_line_relation(PT p, T r, PT a, PT b){
 ld d = dist_from_point_to_line(a, b, p);
 if(sign(d - r) < 0) return 2;
 if(sign(d - r) == 0) return 1; return 0; }
VPT circle_line_intersection(PT c,T r,PT a,PT b){
 VPT ret; b = b - a; a = a - c;
 T A = dot(b, b), B = dot(a, b);
 T C=dot(a,a)-r*r, D=B*B-A*C; if(D < -eps) return ret;
 ret.push_back(c + a + b * (-B + sqrt(D + eps)) / A);
 if(D > eps) ret.push_back(c+a+b*(-B - sqrt(D)) / A);
 return ret;  }
//5 - outside and do not intersect
//4 - intersect outside in one point
//3 - intersect in 2 points
//2 - intersect inside in one point
//1 - inside and do not intersect
int circle_circle_relation(PT a, T r, PT b, T R) {
 ld d = dist(a, b); if(sign(d - r - R) > 0) return 5;
 if(sign(d - r - R) == 0) return 4; ld l = fabs(r - R);
 if(sign(d - r - R) < 0 && sign(d - l) > 0) return 3;
 if(sign(d-l) == 0) return 2; if(sign(d-l)<0) return 1;
 assert(0); return -1; }
VPT circle_circle_intersection(PT a,T r,PT b,T R){
 if(a == b && sign(r-R) == 0) return {PT(1e18, 1e18)};
 VPT ret; ld d = sqrt(dist2(a, b));
 if(d > r + R || d + min(r,R) < max(r,R)) return ret;
 ld x = (d * d - R * R + r * r) / (2 * d);
 ld y = sqrt(r * r - x * x); PT v = (b - a) / d;
 ret.push_back(a + v * x + rotate_ccw90(v) * y);
 if(y>0) ret.push_back(a + v * x - rotate_ccw90(v) * y);
 return ret; }
//circle through points a, b and of radius r, return cnt
int get_circle(PT a, PT b, T r, circle &c1, circle &c2){
 VPT v = circle_circle_intersection(a, r, b, r);
 int t = v.size(); if(!t) return 0; c1.p=v[0], c1.r=r;
 if(t == 2) c2.p = v[1], c2.r = r; return t; }
//returns two circle c1, c2 which is tangent to line u
//goes through point q and has radius r1; returns cnt
int get_circle(line u,PT q,T r1,circle &c1,circle &c2){
 T d = dist_from_point_to_line(u.a, u.b, q);
 if(sign(d - r1 * 2.0) > 0) return 0;
 if(sign(d) == 0) {
  c1.p = q + rotate_ccw90(u.v).truncate(r1);
  c2.p = q + rotate_cw90(u.v).truncate(r1);
  c1.r = c2.r = r1; return 2; }
 line u1 = line(u.a + rotate_ccw90(u.v).truncate(r1), u.b + rotate_ccw90(u.v).truncate(r1));
 line u2 = line(u.a + rotate_cw90(u.v).truncate(r1), u.b + rotate_cw90(u.v).truncate(r1));
 circle cc = circle(q, r1); PT p1, p2; VPT v;
 v = circle_line_intersection(q, r1, u1.a, u1.b);
if(!v.size()) v=circle_line_intersection(q,r1,u2.a,u2.b);
 v.push_back(v[0]); p1 = v[0], p2 = v[1];
 c1 = circle(p1, r1);
 if(p1 == p2) { c2 = c1; return 1; }
 c2 = circle(p2, r1); return 2; }
//area of intersection between two circles
ld circle_circle_area(PT a, T r1, PT b, T r2){
 ld d=(a-b).norm(); if(r1+r2 < d+eps) return 0;
 if(r1 + d < r2 + eps) return PI * r1 * r1;
 if(r2 + d < r1 + eps) return PI * r2 * r2;
 double theta_1 = acos((r1 * r1 + d * d - r2 * r2) / (2 * r1 * d)), theta_2 = acos((r2 * r2 + d * d - r1 * r1)/(2 * r2 * d));
 return r1 * r1 * (theta_1 - sin(2 * theta_1)/2.) + r2 * r2 * (theta_2 - sin(2 * theta_2)/2.); }
//tangent lines from point q to the circle
int tangent_lines_from_point(PT p, T r, PT q, line &u, line &v) { int x = sign(dist2(p, q) - r * r);
 if(x<0) return 0;
 if(x==0){u=line(q,q+rotate_ccw90(q-p));v=u; return 1;}
 ld d=dist(p,q); ld l=r*r/d; ld h=sqrt(r*r - l*l);
 u = line(q, p + ((q - p).truncate(l) + (rotate_ccw90(q - p).truncate(h))));
 v = line(q, p + ((q - p).truncate(l) + (rotate_cw90(q - p).truncate(h)))); return 2; }
// returns outer tangents line of two circles
// if inner == 1 it returns inner tangent lines
int tangents_lines_from_circle(PT c1, T r1, PT c2, T r2, bool inner, line &u, line &v) {
 if(inner) r2 = -r2; PT d = c2 - c1;
 T dr = r1 - r2, d2 = d.norm(), h2 = d2 - dr * dr;
 if(d2 == 0 || h2 < 0) { return 0;}
 vector<pair<PT, PT>>out;
 for(int tmp: {- 1, 1}) {
  PT v = (d*dr+rotate_ccw90(d)*sqrt(h2)*tmp)/d2;
  out.push_back({c1 + v * r1, c2 + v * r2}); }
 u = line(out[0].first, out[0].second);
 if(out.size()==2) v=line(out[1].first,out[1].second);
 return 1 + (h2 > 0); }
//O(n^2 log n)
struct CircleUnion {
 int n,covered[2020]; T x[2020],y[2020],r[2020];
 vector<pair<T, T> > seg, cover; T arc, pol;
 int sign(T x) {return x < -eps ? -1 : x > eps;}
 int sign(T x, T y) {return sign(x - y);}
 T SQ(const T x) {return x * x;}
 ld dist(T x1, T y1, T x2, T y2) {return sqrt(SQ(x1 - x2) + SQ(y1 - y2));}
 ld angle(T A, T B, T C) {
  ld val = (SQ(A) + SQ(B) - SQ(C)) / (2 * A * B);
  if(val < -1) val = -1; if(val > +1) val = +1;
  return acos(val); }
 CircleUnion(){
  n = 0; seg.clear(), cover.clear(); arc = pol = 0; }
 void init(){
  n = 0; seg.clear(), cover.clear(); arc = pol = 0; }
 void add(T xx, T yy, T rr) {
  x[n]=xx, y[n]=yy, r[n]=rr, covered[n]=0, n++; }
 void getarea(int i, T lef, T rig){
  arc += 0.5*r[i]*r[i]*(rig - lef - sin(rig - lef));
  ld x1 = x[i]+r[i]*cos(lef),y1 = y[i]+r[i]*sin(lef);
  ld x2 = x[i]+r[i]*cos(rig),y2 = y[i]+r[i]*sin(rig);
  pol += x1 * y2 - x2 * y1; }
 ld solve() {
 for(int i = 0; i < n; i++) {
  for(int j = 0; j < i; j++) {
   if(!sign(x[i] - x[j]) && !sign(y[i] - y[j]) && !sign(r[i] - r[j])) { r[i] = 0.0; break; }}}
 for(int i = 0; i < n; i++) {
  for(int j = 0; j < n; j++) {
   if(i != j && sign(r[j] - r[i]) >= 0 && sign(dist(x[i], y[i], x[j], y[j]) - (r[j] - r[i])) <= 0) { covered[i] = 1; break; }}}
 for(int i = 0; i < n; i++) {
  if(sign(r[i]) && !covered[i]) {
   seg.clear();
   for(int j = 0; j < n; j++) {
    if(i != j) {
     ld d = dist(x[i], y[i], x[j], y[j]);
     if(sign(d - (r[j] + r[i])) >= 0 || sign(d - abs(r[j] - r[i])) <= 0) { continue; }
     ld alpha = atan2(y[j] - y[i], x[j] - x[i]);
     ld beta = angle(r[i], d, r[j]);
     pair<ld,ld> tmp(alpha - beta, alpha + beta);
     if(sign(tmp.first)<=0 && sign(tmp.second)<=0){ seg.push_back(pair<ld, ld>(2 * PI + tmp.first, 2 * PI + tmp.second)); }
     else if (sign(tmp.first) < 0){
      seg.push_back(pair<ld, ld>(2 * PI + tmp.first, 2 * PI)); seg.push_back(pair<ld, ld>(0, tmp.second)); }
     else { seg.push_back(tmp); } }}
     sort(seg.begin(), seg.end()); ld rig = 0;
     for(auto it=seg.begin(); it!=seg.end();it++){
      if(sign(rig - it->first) >= 0) {
       rig = max(rig, it->second); }
      else { getarea(i, rig, it->first);
       rig = it->second; } }
     if(!sign(rig)) { arc += r[i] * r[i] * PI;}
     else { getarea(i, rig, 2 * PI); }} }
  return pol / 2.0 + arc; }   } CU;
ld area_of_triangle(PT a, PT b, PT c) {
 return fabs(cross(b - a, c - a) * 0.5); }
//-1->inside, 0->on the polygon, 1->outside
int is_point_in_triangle(PT a, PT b, PT c, PT p) {
 if(sign(cross(b - a,c - a)) < 0) swap(b, c);
 int c1 = sign(cross(b - a,p - a));
 int c2 = sign(cross(c - b,p - b));
 int c3 = sign(cross(a - c,p - c));
 if(c1<0 || c2<0 || c3 < 0) return 1;
 if(c1 + c2 + c3 != 3) return 0; return -1; }
T perimeter(VPT &p) {
 T ans=0; int n = p.size();
 for(int i=0;i<n;i++) ans+=dist(p[i],p[(i+1) % n]);
 return ans;  }
ld area(VPT &p) { ld ans=0; int n=p.size();
 for(int i=0;i<n;i++) ans+=cross(p[i],p[(i+1)%n]);
 return fabs(ans) * 0.5;  }
//centroid of a (possibly non-convex) polygon, assuming that the coordinates are listed in a CW or CCW fashion
PT centroid(VPT &p) {
 int n = p.size(); PT c(0, 0); ld sum = 0;
 for(int i=0;i<n;i++) sum+=cross(p[i],p[(i+1)%n]);
 ld scale = 3.0 * sum;
 for(int i = 0; i < n; i++) {
  int j=(i+1)%n; c=c+(p[i]+p[j])*cross(p[i],p[j]); }
 return c / scale;  }
// 0 if cw, 1 if ccw
bool get_direction(VPT &p) {
 ld ans = 0; int n = p.size();
 for(int i=0;i<n;i++) ans += cross(p[i],p[(i+1)%n]);
 if(sign(ans) > 0) return 1; return 0;  }
//returns a point such that the sum of distances
//from that point to all points in p is minimum
//O(n log^2 MX)
PT geometric_median(VPT p) {
 auto tot_dist = [&](PT z) {
  T res = 0;
  for(int i=0; i<p.size(); i++) res += dist(p[i], z);
  return res; };
 auto findY = [&](T x) {
  T yl = -1e5, yr = 1e5;
  for(int i = 0; i < 60; i++) {
   T ym1 = yl + (yr - yl) / 3;
   T ym2 = yr - (yr - yl) / 3;
   T d1 = tot_dist(PT(x, ym1));
   T d2 = tot_dist(PT(x, ym2));
   if (d1 < d2) yr = ym2; else yl = ym1; }
  return pair<T, T> (yl, tot_dist(PT(x, yl)));};
 T xl = -1e5, xr = 1e5;
 for(int i = 0; i < 60; i++) {
  T xm1 = xl + (xr - xl) / 3;
  T xm2 = xr - (xr - xl) / 3;
  T y1, d1, y2, d2;
  auto z=findY(xm1); y1=z.first; d1=z.second;
  z=findY(xm2); y2 = z.first; d2 = z.second;
  if (d1 < d2) xr = xm2; else xl = xm1; }
 return {xl, findY(xl).first };  }
VPT convex_hull(VPT &p) {
 if(p.size() <= 1) return p; VPT v = p;
 sort(v.begin(), v.end()); VPT up, dn;
 for(auto& p : v) {
  while (up.size() > 1 && orientation(up[up.size() - 2], up.back(), p) >= 0) { up.pop_back(); }
  while (dn.size() > 1 && orientation(dn[dn.size() - 2], dn.back(), p) <= 0) { dn.pop_back(); }
  up.push_back(p); dn.push_back(p); } v = dn;
 if(v.size() > 1) v.pop_back();
 reverse(up.begin(), up.end());
 up.pop_back(); for(auto& p:up){v.push_back(p);}
 if(v.size() == 2 && v[0] == v[1]) v.pop_back();
 return v; }
bool is_convex(VPT &p) {
 bool s[3]; s[0]=s[1]=s[2]=0; int n = p.size();
 for(int i = 0; i < n; i++) {
  int j = (i+1)%n; int k = (j + 1) % n;
  s[sign(cross(p[j] - p[i], p[k] - p[i])) + 1] = 1;
  if(s[0] && s[2]) return 0; } return 1; }
//-1->inside,0->on boundary,1->outside;// O(log n)
int is_point_in_convex(VPT &p, const PT& x) {
 int n = p.size(); int a = orientation(p[0], p[1], x);
 int b=orientation(p[0],p[n-1],x);if(a<0||b>0) return 1;
 int l = 1, r = n - 1;
 while(l + 1 < r) { int mid = l + r >> 1;
  if(orientation(p[0], p[mid], x) >= 0) l = mid;
  else r = mid; }
 int k=orientation(p[l],p[r],x); if(k<=0) return -k;
 if(l == 1 && a == 0) return 0; 
 if(r==n-1 && b==0) return 0; return -1; }
bool is_point_on_polygon(VPT &p, const PT& z) {
 int n = p.size();
 for(int i = 0; i < n; i++)
  if(is_point_on_seg(p[i], p[(i+1)%n], z)) return 1;
 return 0; }
int winding_number(VPT &p, const PT& z){
 if (is_point_on_polygon(p, z)) return 1e9;
 int n = p.size(), ans = 0;
 for(int i = 0; i < n; ++i) {
  int j=(i+1)%n; bool below = p[i].y < z.y;
  if(below != (p[j].y < z.y)) {
   auto orient = orientation(z, p[j], p[i]);
   if(orient == 0) return 0;
   if(below == (orient > 0)) ans += below ? 1 : -1;} } return ans; }
int is_point_in_polygon(VPT &p, const PT& z){
 int k = winding_number(p, z);
 return k == 1e9 ? 0 : k == 0 ? 1 : -1; }
//id of the vertex having maximum dot product with z
//top - upper right vertex
//for minimum dot prouct negate z and return -dot(z,p[id])
int extreme_vertex(VPT &p,const PT &z,int top){
 int n = p.size(); if (n == 1) return 0;
 T ans = dot(p[0], z); int id = 0;
 if(dot(p[top], z) > ans) ans=dot(p[top], z), id=top;
 int l = 1, r = top - 1;
 while(l < r){  int mid = l + r >> 1;
  if(dot(p[mid+1], z) >= dot(p[mid], z)) l = mid+1;
  else r = mid; }
 if(dot(p[l], z) > ans) ans = dot(p[l], z), id = l;
 l = top + 1, r = n - 1;
 while(l < r){  int mid = l + r >> 1;
  if(dot(p[(mid+1)%n],z)>=dot(p[mid],z)) l = mid+1;
  else r = mid; } l %= n;
 if(dot(p[l],z)>ans) ans=dot(p[l],z),id=l; return id; }
ld diameter(VPT &p){ int n = (int)p.size();
 if(n==1) return 0; if(n==2) return dist(p[0],p[1]);
 T ans = 0; int i = 0, j = 1;
 while(i < n){
  while(cross(p[(i+1)%n]-p[i],p[(j+1)%n]-p[j])>=0){
   ans=max(ans, dist2(p[i], p[j])); j = (j+1)%n; }
  ans = max(ans, dist2(p[i], p[j])); i++;} return sqrt(ans); }
ld width(VPT &p){
 int n = p.size(); if(n <= 2) return 0;
 ld ans = inf; int i = 0, j = 1;
 while(i < n){
  while(cross(p[(i + 1) % n] - p[i], p[(j + 1) % n] - p[j]) >= 0) j = (j + 1) % n;
  ans = min(ans, dist_from_point_to_line(p[i], p[(i + 1) % n], p[j])); i++; } return ans; }
ld minimum_enclosing_rectangle(VPT &p){
 int n = p.size(); if (n <= 2) return perimeter(p);
 int mndot = 0; T tmp = dot(p[1] - p[0], p[0]);
 for(int i = 1; i < n; i++) {
  if(dot(p[1] - p[0], p[i]) <= tmp) {
   tmp = dot(p[1] - p[0], p[i]); mndot = i; } }
 ld ans = inf; int i = 0, j = 1, mxdot = 1;
 while(i < n){ PT cur = p[(i + 1) % n] - p[i];
  while(cross(cur,p[(j+1)%n]-p[j])>=0) j=(j+1)%n;
  while (dot(p[(mxdot + 1) % n], cur) >= dot(p[mxdot], cur)) mxdot = (mxdot + 1) % n;
  while (dot(p[(mndot + 1) % n], cur) <= dot(p[mndot], cur)) mndot = (mndot + 1) % n;
  ans = min(ans, 2.0 * ((dot(p[mxdot], cur) / cur.norm() - dot(p[mndot], cur) / cur.norm()) + dist_from_point_to_line(p[i], p[(i + 1) % n], p[j]))); i++; } return ans; }
//expected O(n), 1st call convex_hull() for faster
circle minimum_enclosing_circle(VPT &p){
 random_shuffle(p.begin(), p.end());
 int n=p.size(); circle c(p[0], 0);
 for(int i = 1; i < n; i++) {
  if(sign(dist(c.p, p[i]) - c.r) > 0) {
   c = circle(p[i], 0);
   for(int j = 0; j < i; j++) {
    if(sign(dist(c.p, p[j]) - c.r) > 0) {
     c=circle((p[i]+p[j])/2,dist(p[i],p[j])/2);
     for(int k = 0; k < j; k++) {
      if(sign(dist(c.p, p[k]) - c.r) > 0)
        c = circle(p[i], p[j], p[k]);} } } } } return c; }
//returns a vector with the vertices of a polygon with everything to the left of the line going from a to b cut away
VPT cut(VPT &p, PT a, PT b) {
 VPT ans; int n = (int)p.size();
 for(int i = 0; i < n; i++) {
  T c1=cross(b-a,p[i]-a); T c2=cross(b-a,p[(i+1)%n]-a);
  if(sign(c1) >= 0) ans.push_back(p[i]);
  if(sign(c1 * c2) < 0) {
   if(!is_parallel(p[i], p[(i + 1) % n], a, b)){
    PT tmp; line_line_intersection(p[i],p[(i+1)%n],a,b,tmp); ans.push_back(tmp); }} } return ans; }
//not necessarily convex, boundary is included in the intersection, returns total intersected length
ld polygon_line_intersection(VPT &p, PT a, PT b){
 int n=p.size(); p.push_back(p[0]); line l=line(a, b);
 ld ans = 0.0; vector< pair<ld, int> > vec;
 for(int i = 0; i < n; i++) {
  int s1 = sign(cross(b - a, p[i] - a));
  int s2 = sign(cross(b - a, p[i+1] - a));
  if(s1 == s2) continue;
  line t = line(p[i], p[i + 1]);
  PT inter = (t.v * l.c - l.v * t.c) / cross(l.v, t.v);
  ld tmp = dot(inter, l.v); int f;
  if(s1 > s2) f = s1 && s2 ? 2 : 1;
  else f = s1 && s2 ? -2 : -1;
  vec.push_back(make_pair(tmp, f)); }
 sort(vec.begin(), vec.end());
 for(int i=0,j=0; i+1 < (int)vec.size(); i++){
  j += vec[i].second;
  if (j) ans += vec[i + 1].first - vec[i].first; }
 ans = ans / sqrt(dot(l.v, l.v)); p.pop_back();
 return ans; }
pair<PT,int> point_poly_tangent(VPT &p,PT Q,int dir,int l,int r){
 while(r - l > 1) {
  int mid = (l + r) >> 1;
  bool pvs = orientation(Q, p[mid], p[mid - 1]) != -dir;
  bool nxt = orientation(Q, p[mid], p[mid + 1]) != -dir;
  if(pvs && nxt) return {p[mid], mid};
  if(!(pvs || nxt)) {
   auto p1=point_poly_tangent(p, Q, dir, mid + 1, r);
   auto p2=point_poly_tangent(p, Q, dir, l, mid - 1);
return orientation(Q,p1.first,p2.first) == dir ? p1:p2;}
  if(!pvs) {
   if(orientation(Q,p[mid],p[l]) == dir) r=mid-1;
   else if(orientation(Q,p[l],p[r]) == dir) r=mid-1;
   else l = mid + 1; }
  if(!nxt) {
   if(orientation(Q,p[mid],p[l]) == dir) l=mid+1;
   else if(orientation(Q,p[l],p[r]) == dir) r=mid-1;
   else l = mid + 1; } }
 pair<PT, int> ret = {p[l], l};
 for(int i = l + 1; i <= r; i++) ret = orientation(Q, ret.first, p[i]) != dir ? make_pair(p[i], i) : ret; return ret; }
//(cw, ccw) tangents from a point that is outside this convex polygon returns indexes of the points
pair<int,int> tangents_from_point_to_polygon(VPT &p,PT Q){
 int cw=point_poly_tangent(p,Q,1,0,p.size()-1).second;
 int ccw=point_poly_tangent(p,Q,-1,0,p.size()-1).second;
 return make_pair(cw, ccw); }
// minimum distance from a point to a convex polygon
// it assumes point does not lie strictly inside the polygon
ld dist_from_point_to_polygon(vector<PT> &p, PT z) {
 ld ans = inf; int n = p.size();
 if (n <= 3) {
  for(int i = 0; i < n; i++) ans = min(ans, dist_from_point_to_seg(p[i], p[(i + 1) % n], z));
  return ans; }
 auto [r, l] = tangents_from_point_to_polygon(p, z);
 if(l > r) r += n;
 while (l < r) {
  int mid = (l + r) >> 1;
  ld left = dist2(p[mid % n], z), right= dist2(p[(mid + 1) % n], z);
  ans = min({ans, left, right});
  if(left < right) r = mid; else l = mid + 1; }
 ans = sqrt(ans);
 ans = min(ans, dist_from_point_to_seg(p[l % n], p[(l + 1) % n], z));
 ans = min(ans, dist_from_point_to_seg(p[l % n], p[(l - 1 + n) % n], z));
 return ans; }
ld dist_from_polygon_to_line(VPT &p,PT a,PT b,int top){ PT orth = (b - a).perp();
 if(orientation(a, b, p[0]) > 0) orth = (a - b).perp();
 int id = extreme_vertex(p, orth, top);
 if(dot(p[id] - a, orth) > 0) return 0.0;
 return dist_from_point_to_line(a, b, p[id]); }
ld dist_from_polygon_to_polygon(VPT &p1, VPT &p2){
 ld ans = inf; //NLogN
 for(int i = 0; i < p1.size(); i++)
  ans=min(ans, dist_from_point_to_polygon(p2, p1[i]));
 for(int i = 0; i < p2.size(); i++)
  ans=min(ans, dist_from_point_to_polygon(p1, p2[i]));
 return ans; }
ld maximum_dist_from_polygon_to_polygon(VPT &u, VPT &v){
 int n=u.size(), m=v.size(); ld ans = 0; //0(n)
 if(n < 3 || m < 3) {
  for(int i = 0; i < n; i++) {
  for(int j=0; j<m; j++) ans=max(ans,dist2(u[i],v[j]));} return sqrt(ans); }
 if(u[0].x > v[0].x) swap(n, m), swap(u, v);
 int i = 0, j = 0, step = n + m + 10;
 while (j + 1 < m && v[j].x < v[j + 1].x) j++ ;
 while (step--) {
 if(cross(u[(i+1)%n]-u[i],v[(j+1)%m]-v[j])>=0) j=(j+1)%m;
 else i = (i + 1) % n; ans = max(ans, dist2(u[i], v[j]));} return sqrt(ans); }
//n polygons(not necessarily convex). points within each polygon must be given in CCW order. O(N^2), N total points
ld rat(PT a, PT b, PT p) {
 return !sign(a.x - b.x) ? (p.y - a.y) / (b.y - a.y) : (p.x - a.x) / (b.x - a.x); };
ld polygon_union(vector<VPT> &p) {
 int n = p.size(); ld ans=0;
 for(int i = 0; i < n; ++i) {
  for(int v = 0; v < (int)p[i].size(); ++v){
   PT a = p[i][v], b = p[i][(v + 1) % p[i].size()];
   vector<pair<ld, int>> segs;
   segs.emplace_back(0,0), segs.emplace_back(1,0);
   for(int j = 0; j < n; ++j) if(i != j){
    for(size_t u = 0; u < p[j].size(); ++u) {
      PT c=p[j][u], d=p[j][(u+1)%p[j].size()];
      int sc = sign(cross(b - a, c - a)), sd = sign(cross(b - a, d - a));
      if(!sc && !sd) {
      if(sign(dot(b - a, d - c)) > 0 && i > j)
        segs.emplace_back(rat(a, b, c), 1), segs.emplace_back(rat(a, b, d), -1); }
      else {
      ld sa=cross(d-c,a-c), sb=cross(d-c,b-c);
      if(sc >= 0 && sd < 0) segs.emplace_back(sa / (sa - sb), 1);
      else if(sc < 0 && sd >= 0) segs.emplace_back(sa / (sa - sb), -1);} } }
   sort(segs.begin(), segs.end());
   ld pre=min(max(segs[0].first,(ld)0.0),(ld)1.0),now, sum = 0; int cnt = segs[0].second;
   for(int j = 1; j < segs.size(); ++j) {
     now=min(max(segs[j].first, (ld)0.0), (ld)1.0);
     if(!cnt) sum += now-pre; cnt += segs[j].second;
     pre = now; }
   ans += cross(a, b) * sum; } } return ans * 0.5; }
struct HP {
 PT a, b; HP() {}
 HP(PT a, PT b) : a(a), b(b) {}
 HP(const HP& rhs) : a(rhs.a), b(rhs.b) {}
 int operator < (const HP& rhs) const {
  PT p = b - a; PT q = rhs.b - rhs.a;
  int fp = (p.y < 0 || (p.y == 0 && p.x < 0));
  int fq = (q.y < 0 || (q.y == 0 && q.x < 0));
  if (fp != fq) return fp == 0;
  if (cross(p, q)) return cross(p, q) > 0;
  return cross(p, rhs.b - a) < 0; }
 PT line_line_intersection(PT a, PT b, PT c, PT d){
  b = b - a; d = c - d; c = c - a;
  return a + b * cross(c, d) / cross(b, d); }
 PT intersection(const HP &v) {
  return line_line_intersection(a,b,v.a,v.b); } };
int check(HP a, HP b, HP c) {
 return cross(a.b - a.a, b.intersection(c) - a.a) > -eps;}
//consider half-plane of counter-clockwise side of each line. returns a convex polygon, a point can occur multiple times though. O(n log(n))
VPT half_plane_intersection(vector<HP> h) {
 sort(h.begin(), h.end()); vector<HP> tmp;
 for(int i = 0; i < h.size(); i++) {
  if(!i || cross(h[i].b-h[i].a, h[i-1].b-h[i-1].a))
   tmp.push_back(h[i]); }
 h=tmp; vector<HP> q(h.size()+10); int qh=0,qe=0;
 for(int i = 0; i < h.size(); i++) {
  while(qe-qh>1 && !check(h[i],q[qe-2],q[qe-1])) qe--;
  while(qe-qh > 1 && !check(h[i],q[qh],q[qh+1])) qh++;
  q[qe++] = h[i]; }
 while(qe-qh>2 && !check(q[qh],q[qe-2],q[qe-1])) qe--;
 while(qe-qh>2 && !check(q[qe-1],q[qh],q[qh+1])) qh++;
 vector<HP> res; VPT hull; int n = res.size();
 for(int i = qh; i < qe; i++) res.push_back(q[i]);
 if(n > 2) {
  for(int i = 0; i < n; i++)
   hull.push_back(res[i].intersection(res[(i+1)%n]));} return hull; }
//a and b -> strictly convex polygons of DISTINCT points
VPT minkowski_sum(VPT &a, VPT &b) {
 int n = a.size(), m = b.size(),i = 0, j = 0; VPT c;
 c.push_back(a[i] + b[j]);
 while(1){
  PT p1=a[i]+b[(j+1)% m], p2=a[(i+1)%n]+b[j];
  int t = orientation(c.back(), p1, p2);
  if(t>=0) j=(j+1)%m; if(t<=0) i=(i+1)%n, p1=p2;
  if(t == 0) p1 = a[i] + b[j]; if(p1 == c[0]) break;
  c.push_back(p1); } return c;  }
ld triangle_circle_intersection(PT c,T r,PT a,PT b){
 ld sd1 = dist2(c, a), sd2 = dist2(c, b);
 if(sd1 > sd2) swap(a, b), swap(sd1, sd2);
 ld sd = dist2(a, b);
 ld d1 = sqrtl(sd1), d2 = sqrtl(sd2), d = sqrt(sd);
 ld x=abs(sd2-sd-sd1)/(2*d); ld h=sqrtl(sd1-x*x);
 if(r >= d2) return h * d / 2; ld area = 0;
 if(sd + sd1 < sd2) {
  if(r < d1) area=r*r*(acos(h/d2)-acos(h/d1))/2;
  else{
   area=r*r*(acos(h/d2)-acos(h/r))/2;
   ld y=sqrtl(r*r-h*h); area+=h*(y-x)/2; } }
 else {
  if(r<h) area=r*r*(acos(h/d2) + acos(h/d1)) / 2;
  else {
   area += r * r * (acos(h / d2) - acos(h / r)) / 2;
   ld y=sqrtl(r*r-h*h); area += h * y / 2;
   if(r < d1) {
    area += r*r*(acos(h / d1) - acos(h / r)) / 2;
    area+=h*y/2; } 
   else area += h*x/2; } } return area; }
//intersection between a simple polygon and a circle
ld polygon_circle_intersection(VPT &v, PT p, T r) {
 int n=v.size(); ld ans=0.00; PT org = {0, 0};
 for(int i = 0; i < n; i++) {
  int x = orientation(p, v[i], v[(i + 1) % n]);
  if(x == 0) continue;
  ld area = triangle_circle_intersection(org, r, v[i] - p, v[(i + 1) % n] - p); if (x < 0) ans -= area;
  else ans += area; } return ans * 0.5; }
//circle of radius r that contains as many points as possible. O(n^2 log n);
int maximum_circle_cover(VPT p, T r, circle &c) {
 int n = p.size(),ans = 0,id = 0; ld th = 0;
 for(int i = 0; i < n; ++i) {
  vector<pair<ld, int>> events = {{-PI, +1},{PI, -1}};
  for(int j = 0; j < n; ++j) {
   if(j == i) continue;
   ld d = dist(p[i], p[j]); if(d > r * 2) continue;
   ld dir = (p[j]-p[i]).arg(); ld ang = acos(d/2/r);
   ld st=dir-ang, ed=dir+ang; if(st>PI) st -= PI*2;
   if(st <= -PI) st += PI*2; if(ed > PI) ed -= PI*2;
   if(ed<=-PI) ed += PI*2;events.push_back({st, +1});
   events.push_back({ed, -1});
   if(st > ed) {
events.push_back({-PI,1});events.push_back({PI,-1});} }
  sort(events.begin(), events.end()); int cnt = 0;
  for(auto &&e: events) {
   cnt += e.second;
   if(cnt>ans){ ans=cnt; id=i; th=e.first; } } }
 PT w=PT(p[id].x+r*cos(th), p[id].y + r * sin(th));
 c = circle(w, r); return ans; }
ld maximum_inscribed_circle(VPT p){ //returns radius
 int n=p.size(); if(n<3) return 0; ld l=0, r=20000;
 while(r - l > eps){ vector<HP> h;
  ld mid = (l + r) * 0.5; const int L = 1e9;
  h.push_back(HP(PT(-L, -L), PT(L, -L)));
  h.push_back(HP(PT(L, -L), PT(L, L)));
  h.push_back(HP(PT(L, L), PT(-L, L)));
  h.push_back(HP(PT(-L, L), PT(-L, -L)));
  for(int i = 0; i < n; i++) {
   PT z=(p[(i+1)%n]-p[i]).perp(); z=z.truncate(mid);
   PT y = p[i] + z, q = p[(i + 1) % n] + z;
   h.push_back(HP(p[i] + z, p[(i + 1) % n] + z)); }
  VPT nw = half_plane_intersection(h);
  if(!nw.empty()) l = mid; else r=mid; } return l;}
\end{minted}
\subsection{Closest Pair}
\begin{minted}{c++}
#define x first
#define y second
long long dist2(pair<int, int> a, pair<int, int> b) {
 return 1LL * (a.x - b.x) * (a.x - b.x) + 1LL * (a.y - b.y) * (a.y - b.y);}pair<int, int> closest_pair(vector<pair<int, int>> a) {
 int n = a.size();
 assert(n >= 2);
 vector<pair<pair<int, int>, int>> p(n);
 for (int i = 0; i < n; i++) p[i] = {a[i], i};
 sort(p.begin(), p.end());
 int l = 0, r = 2;
 long long ans = dist2(p[0].x, p[1].x);
 pair<int, int> ret = {p[0].y, p[1].y};
 while (r < n) {
  while (l < r && 1LL * (p[r].x.x - p[l].x.x) * (p[r].x.x - p[l].x.x) >= ans) l++;
  for (int i = l; i < r; i++) {
   long long nw = dist2(p[i].x, p[r].x);
   if (nw < ans) {
    ans = nw;
    ret = {p[i].y, p[r].y};}}r++;}
 return ret;}
\end{minted}
\subsection{Pick Theorem}
\begin{minted}{c++}
theorem:area=In+(on_boundary/2)-1;in=strictly inside polygon
T point(PT a,PT b){  T x=abs(a.x-b.x);
  T y=abs(a.y-b.y); return __gcd(x,y)-1; }
main(){ ll on_boundary=n,area=0LL;
 for(i=0; i<n; i++){ area+=cross(p[i],p[(i+1)%n]);
   on_boundary+=point(p[i],p[(i+1)%n]); } }
\end{minted}
\subsection{Transform}
\begin{minted}{c++}
void translate(double dx, double dy, double dz)
{array <array <double, 4>, 4> mat2;
 mat2[0] = {1, 0, 0, dx};
 mat2[1] = {0, 1, 0, dy};
 mat2[2] = {0, 0, 1, dz};
 mat2[3] = {0, 0, 0, 1};
 multiply(mat2);}
void scale(double dx, double dy, double dz)
{array <array <double, 4>, 4> mat2;
 mat2[0] = {dx, 0, 0, 0};
 mat2[1] = {0, dy, 0, 0};
 mat2[2] = {0, 0, dz, 0};
 mat2[3] = {0, 0, 0, 1};
 multiply(mat2);}
void rotate (double x, double y, double z, double d)
{d = d / 180.0 * acosl(-1.0);
 double norm = sqrtl(x * x + y * y + z * z);
 x /= norm; y /= norm; z /= norm;
 array <array <double, 4>, 4> mat2;
 mat2[0] = {(1-cosl(d))*x*x+cosl(d), (1-cosl(d))*x*y-sinl(d)*z, (1-cosl(d))*x*z+sinl(d)*y, 0};
 mat2[1] = {(1-cosl(d))*y*x+sinl(d)*z, (1-cosl(d))*y*y+cosl(d), (1-cosl(d))*y*z-sinl(d)*x, 0};
 mat2[2] = {(1-cosl(d))*z*x-sinl(d)*y, (1-cosl(d))*z*y+sinl(d)*x, (1-cosl(d))*z*z+cosl(d), 0};
 mat2[3] = {0, 0, 0, 1};
 multiply(mat2);}
\end{minted}
\subsection{dynamic Hull}
\begin{minted}{c++}
typedef set<pair<int, int> >::iterator iter;
struct PT {
 int x, y;
 PT() : x(0), y(0) {}
 PT(int x, int y) : x(x), y(y) {}
 PT(const PT& rhs) : x(rhs.x), y(rhs.y) {}
 PT(const iter& p) : x(p->first), y(p->second) {}
 PT operator - (const PT& rhs) const {
  return PT(x - rhs.x, y - rhs.y);}};
long long cross(PT a, PT b) {
 return (long long) a.x * b.y - (long long) a.y * b.x;}
inline int inside(set<pair<int, int>>& hull, const PT& p) { //border inclusive
 int x = p.x, y = p.y;
 iter p1 = hull.lower_bound(make_pair(x, y));
 if (p1 == hull.end()) return 0;
 if (p1->first == x) return p1 != hull.begin() && y <= p1->second;
 if (p1 == hull.begin()) return 0;
 iter p2(p1--);
 return cross(p - PT(p1), PT(p2) - p) >= 0;}
inline void del(set<pair<int, int>>& hull, iter it, long long& scross) {
 if (hull.size() == 1) {
  hull.erase(it); return;}
 if (it == hull.begin()) {
  iter p1 = it++;
  scross -= cross(p1, it);
  hull.erase(p1); return;}
 iter p1 = --it, p2 = ++it;
 if (++it == hull.end()) {
  scross -= cross(p1, p2);
  hull.erase(p2); return;}
 scross -= cross(p1, p2) + cross(p2, it) - cross(p1, it);
 hull.erase(p2);}
inline void add(set<pair<int, int>>& hull, iter it, long long& scross) {
 if (hull.size() == 1) return;
 if (it == hull.begin()) { iter p1 = it++; scross += cross(p1, it); return;}
 iter p1 = --it, p2 = ++it;
 if (++it == hull.end()) { scross += cross(p1, p2); return;}
 scross += cross(p1, p2) + cross(p2, it) - cross(p1, it);}
inline void add(set<pair<int, int>>& hull, const PT& p, long long& scross) { //no collinear PTs
 if (inside(hull, p)) return;
 int x = p.x, y = p.y;
 iter pnt = hull.insert(make_pair(x, y)).first, p1, p2;
 add(hull, pnt, scross);
 for ( ; ; del(hull, p2, scross)) {
  p1 = pnt;
  if (++p1 == hull.end()) break;
  p2 = p1;
  if (++p1 == hull.end()) break;
  if (cross(PT(p2) - p, PT(p1) - p) < 0) break;}
 for ( ; ; del(hull, p2, scross)) {
  if ((p1 = pnt) == hull.begin()) break;
  if (--p1 == hull.begin()) break;
  p2 = p1--;
  if (cross(PT(p2) - p, PT(p1) - p) > 0) break;}}
int main() {
 long long ucross = 0, dcross = 0;
 set<pair<int, int>> uhull, dhull;
 PT p[] = {PT(0, 0), PT(3, 0), PT(3, 3), PT(0, 3), PT(0, 1), PT(0, 2), PT(3, 1), PT(3, 2)};
 for (int i = 0; i < 5; i++) {
  add(uhull, PT(+p[i].x, +p[i].y), ucross); add(dhull, PT(-p[i].x, -p[i].y), dcross);}
 cout << fixed << setprecision(1) << "Area: " << fabs(ucross + dcross) / 2.0 << "\n";}
\end{minted}
\section{Graph and Tree}
\subsection{2 SAT}
\begin{minted}{c++}
/*zero Indexed
we have vars variables
F=(x_0 XXX y_0) and (x_1 XXX y_1) and ... (x_{vars-1} XXX y_{vars-1})
here {x_i,y_i} are variables
and XXX belongs to {OR,XOR}
is there any assignment of variables such that F=true*/
struct twosat {
 int n; // total size combining +, -. must be even.
 vector< vector<int> > g, gt;
 vector<bool> vis, res;
 vector<int> comp;
 stack<int> ts;
 twosat(int vars = 0) {
  n = vars << 1;
  g.resize(n);
  gt.resize(n);}
//zero indexed, be careful
 //if you want to force variable a to be true in OR or XOR combination
 //add addOR (a,1,a,1);
 //if you want to force variable a to be false in OR or XOR combination
 //add addOR (a,0,a,0);
 //(x_a or (not x_b))-> af=1,bf=0
 void addOR(int a, bool af, int b, bool bf) {
  a += a + (af ^ 1);
  b += b + (bf ^ 1);
  g[a ^ 1].push_back(b); // !a => b
  g[b ^ 1].push_back(a); // !b => a
  gt[b].push_back(a ^ 1);
  gt[a].push_back(b ^ 1);}
 //(!x_a xor !x_b)-> af=0, bf=0
 void addXOR(int a, bool af, int b, bool bf) {
  addOR(a, af, b, bf);
  addOR(a, !af, b, !bf);}
 void _add(int a,bool af,int b,bool bf) {
  a += a + (af ^ 1);
  b += b + (bf ^ 1);
  g[a].push_back(b);
  gt[b].push_back(a);}
 //add this type of condition->
 //add(a,af,b,bf) means if a is af then b must need to be bf
 void add(int a,bool af,int b,bool bf) {
  _add(a, af, b, bf);
  _add(b, !bf, a, !af);}
 void dfs1(int u) {
  vis[u] = true;
  for(int v : g[u]) if(!vis[v]) dfs1(v);
  ts.push(u);}
 void dfs2(int u, int c) {
  comp[u] = c;
  for(int v : gt[u]) if(comp[v] == -1) dfs2(v, c);}
 bool ok() {
  vis.resize(n, false);
  for(int i = 0; i < n; ++i) if(!vis[i]) dfs1(i);
  int scc = 0;
  comp.resize(n, -1);
  while(!ts.empty()) {
   int u = ts.top();
   ts.pop();
   if(comp[u] == -1) dfs2(u, scc++);}
  res.resize(n / 2);
  for(int i = 0; i < n; i += 2) {
   if(comp[i] == comp[i + 1]) return false;
   res[i / 2] = (comp[i] > comp[i + 1]);}
  return true;}};
\end{minted}
\subsection{Articulation Bridges}
\begin{minted}{c++}
struct TECC { // 0 indexed
 int n, k;
 vector<vector<int>> g, t;
 vector<bool> used;
 vector<int> comp, ord, low;
 using edge = pair<int, int>;
 vector<edge> br;
 void dfs(int x, int prv, int &c) {
  used[x] = 1; ord[x] = c++; low[x] = n;
  bool mul = 0;
  for (auto y:g[x]) {
   if (used[y]) {
    if (y != prv || mul) low[x] = min(low[x], ord[y]);
    else mul = 1;
    continue;}
   dfs(y, x, c);
   low[x] = min(low[x], low[y]);}}
 void dfs2(int x, int num) {
  comp[x] = num;
  for (auto y: g[x]) {
   if (comp[y] != -1) continue;
   if (ord[x] < low[y]) {
    br.push_back({x, y});
    k++;
    dfs2(y, k);} else dfs2(y, num);}}
 TECC(const vector<vector<int>> &g): g(g), n(g.size()), used(n), comp(n, -1), ord(n), low(n), k(0) {
  int c = 0;
  for (int i = 0; i < n; i++) {
   if (used[i]) continue;
   dfs(i, -1, c);
   dfs2(i, k);
   k++;}}
 void build_tree() {
  t.resize(k);
  for (auto e: br) {
   int x = comp[e.first], y = comp[e.second];
   t[x].push_back(y);
   t[y].push_back(x);}}}
;\end{minted}
\subsection{Block Cut Tree}
\begin{minted}{c++}
const int N = 4e5 + 9;
int T, low[N], dis[N], art[N], sz;
vector<int> g[N], bcc[N], st;
void dfs(int u, int pre = 0) {
 low[u] = dis[u] = ++T;
 st.push_back(u);
 for(auto v: g[u]) {
  if(!dis[v]) {
   dfs(v, u);
   low[u] = min(low[u], low[v]);
   if(low[v] >= dis[u]) {
    sz ++;
    int x;
    do {
     x = st.back();
     st.pop_back();
     bcc[x].push_back(sz);} while(x ^ v);
    bcc[u].push_back(sz);}} else if(v != pre) low[u] = min(low[u], dis[v]);}}
int dep[N], par[N][20], cnt[N], id[N];
vector<int> bt[N];
void dfs1(int u, int pre = 0) {
 dep[u] = dep[pre] + 1;
 cnt[u] = cnt[pre] + art[u];
 par[u][0] = pre;
 for(int k = 1; k <= 18; k++) par[u][k] = par[par[u][k - 1]][k - 1];
 for(auto v: bt[u]) if(v != pre) dfs1(v, u);}
int lca(int u, int v) {
 if(dep[u] < dep[v]) swap(u, v);
 for(int k = 18; k >= 0; k--) if(dep[par[u][k]] >= dep[v]) u = par[u][k];
 if(u == v) return u;
 for(int k = 18; k >= 0; k--) if(par[u][k] != par[v][k]) u = par[u][k], v = par[v][k];
 return par[u][0];}
int dist(int u, int v) {
 int lc = lca(u, v);
 return cnt[u] + cnt[v] - 2 * cnt[lc] + art[lc];}
int32_t main() {
 while(m--) {
  int u, v;
  cin >> u >> v;
  g[u].push_back(v);
  g[v].push_back(u);}
 dfs(1);
 for(int u = 1; u <= n; u++) {
  if(bcc[u].size() > 1) { //AP
   id[u] = ++sz;
   art[id[u]] = 1; //if id in BCT is an AP on real graph or not
   for(auto v: bcc[u]) {
    bt[id[u]].push_back(v);
    bt[v].push_back(id[u]);}} else if(bcc[u].size() == 1) id[u] = bcc[u][0];}
 dfs1(1);
 int q; cin >> q;
 while(q--) {
  int u, v; cin >> u >> v; int ans;
  if(u == v) ans = 0;
  else ans = dist(id[u], id[v]) - art[id[u]] - art[id[v]];
  cout << ans << '\n';;//number of articulation points in the path from u to v except u and v
  //u and v are in the same bcc if ans == 0}
 return 0;}
\end{minted}
\subsection{Blossom}
\begin{minted}{c++}
//O(v^2 * E)
const int N = 2e3 + 9;
mt19937 rnd(chrono::steady_clock::now().time_since_epoch().count());
struct Blossom {
 int vis[N], par[N], orig[N], match[N], aux[N], t;
 int n; bool ad[N]; vector<int> g[N];queue<int> Q;
 Blossom(int _n) {
  n = _n; t = 0;
  for (int i = 0; i <= _n; ++i) {
   g[i].clear(); match[i] = aux[i] = par[i] = vis[i] = aux[i] = ad[i] = orig[i] = 0;}}
 void add_edge(int u, int v) {
  g[u].push_back(v); g[v].push_back(u);}
 void augment(int u, int v) {
  int pv = v, nv;
  do {
   pv = par[v]; nv = match[pv]; match[v] = pv; match[pv] = v;
   v = nv;} while (u != pv);}
 int lca(int v, int w) {
  ++t;
  while (true) {
   if (v) {
    if (aux[v] == t) return v;
    aux[v] = t; v = orig[par[match[v]]];}
   swap(v, w);}}
 void blossom(int v, int w, int a) {
  while (orig[v] != a) {
   par[v] = w; w = match[v]; ad[v] = true;
   if (vis[w] == 1) Q.push(w), vis[w] = 0;
   orig[v] = orig[w] = a; v = par[w];}}
 //it finds an augmented path starting from u
 bool bfs(int u) {
  fill(vis + 1, vis + n + 1, -1);
  iota(orig + 1, orig + n + 1, 1);
  Q = queue<int> ();
  Q.push(u); vis[u] = 0;
  while (!Q.empty()) {
   int v = Q.front(); Q.pop(); ad[v] = true;
   for (int x : g[v]) {
    if (vis[x] == -1) {
     par[x] = v; vis[x] = 1;
     if (!match[x]) return augment(u, x), true;
     Q.push(match[x]); vis[match[x]] = 0;}
     else if (vis[x] == 0 && orig[v] != orig[x]) {
     int a = lca(orig[v], orig[x]);
     blossom(x, v, a); blossom(v, x, a);}}}
  return false;}
 int maximum_matching() {
  int ans = 0;
  vector<int> p(n - 1);
  iota(p.begin(), p.end(), 1);
  shuffle(p.begin(), p.end(), rnd);
  for (int i = 1; i <= n; i++) shuffle(g[i].begin(), g[i].end(), rnd);
  for (auto u : p) {
   if (!match[u]) { for(auto v : g[u]) { if (!match[v]) {
      match[u] = v, match[v] = u; ++ans; break;}}}}
  for(int i = 1; i <= n; ++i) if (!match[i] && bfs(i)) ++ans;
  return ans;}}
 M;
\end{minted}
\subsection{Centroid Decomposition}
\begin{minted}{c++}
vector<int> g[N];
int sz[N];
int tot, done[N], cenpar[N];
void calc_sz(int u, int p) {
 tot ++;
 sz[u] = 1;
 for (auto v : g[u]) {
  if(v == p || done[v]) continue;
  calc_sz(v, u);
  sz[u] += sz[v];}}
int find_cen(int u, int p) {
 for (auto v : g[u]) {
  if(v == p || done[v]) continue;
  else if(sz[v] > tot / 2) return find_cen(v, u);}return u;}
void decompose(int u, int pre) {
 tot = 0;
 calc_sz(u, pre);
 int cen = find_cen(u, pre);
 cenpar[cen] = pre;
 done[cen] = 1;
 for(auto v : g[cen]) {
  if(v == pre || done[v]) continue;
  decompose(v, cen);}}
int main() {
 decompose(1, 0);
 for(int i = 1; i <= n; i++) g[i].clear(); int root;
 for(int i = 1; i <= n; i++) {
  g[cenpar[i]].push_back(i);
  g[i].push_back(cenpar[i]);
  if (cenpar[i] == 0) root = i;}}
\end{minted}
\subsection{Dinics Algorithm}
\begin{minted}{c++}
//General graphs O(V^2E)
//Unweighted bipartite graphs O(E * root(V))
//Dense graphs O(V^3)
struct Dinic {
 struct edge {
  int to, rev;
  long long flow, w;
  int id;};
 int n, s, t, mxid;
 vector<int> d, flow_through;
 vector<int> done;
 vector<vector<edge>> g;
 Dinic() {}
 Dinic(int _n) {
  n = _n + 10;
  mxid = 0;
  g.resize(n);}
 void add_edge(int u, int v, long long w, int id = -1) {
  edge a = {v, (int)g[v].size(), 0, w, id};
  edge b = {u, (int)g[u].size(), 0, 0, -2};//for bidirectional edges cap(b) = w
  g[u].emplace_back(a);
  g[v].emplace_back(b);
  mxid = max(mxid, id);}
 bool bfs() {
  d.assign(n, -1);
  d[s] = 0;
  queue<int> q;
  q.push(s);
  while (!q.empty()) {
   int u = q.front();
   q.pop();
   for (auto &e : g[u]) {
    int v = e.to;
    if (d[v] == -1 && e.flow < e.w) d[v] = d[u] + 1, q.push(v);}}
  return d[t] != -1;}
 long long dfs(int u, long long flow) {
  if (u == t) return flow;
  for (int &i = done[u]; i < (int)g[u].size(); i++) {
   edge &e = g[u][i];
   if (e.w <= e.flow) continue;
   int v = e.to;
   if (d[v] == d[u] + 1) {
    long long nw = dfs(v, min(flow, e.w - e.flow));
    if (nw > 0) {
     e.flow += nw;
     g[v][e.rev].flow -= nw;
     return nw;}}}
  return 0;}
 long long max_flow(int _s, int _t) {
  s = _s;
  t = _t;
  long long flow = 0;
  while (bfs()) {
   done.assign(n, 0);
   while (long long nw = dfs(s, inf)) flow += nw;}
  flow_through.assign(mxid + 10, 0);
  for(int i = 0; i < n; i++) for(auto e : g[i]) if(e.id >= 0) flow_through[e.id] = e.flow;
  return flow;}};
\end{minted}
\subsection{Dominator Tree}
\begin{minted}{c++}
const int N = 2e5 + 9;vector<int> g[N];
vector<int> t[N], rg[N], bucket[N]; //t = dominator tree of the nodes reachable from root
int sdom[N], par[N], idom[N], dsu[N], label[N];
int id[N], rev[N], T;
int find_(int u, int x = 0) {
 if(u == dsu[u]) return x ? -1 : u;
 int v = find_(dsu[u], x+1);
 if(v < 0)return u;
 if(sdom[label[dsu[u]]] < sdom[label[u]]) label[u] = label[dsu[u]];
 dsu[u] = v;
 return x ? v : label[u];}void dfs(int u) {
 T++; id[u] = T;
 rev[T] = u; label[T] = T;
 sdom[T] = T; dsu[T] = T;
 for(int i = 0; i < g[u].size(); i++) {
  int w = g[u][i];
  if(!id[w]) dfs(w), par[id[w]] = id[u];
  rg[id[w]].push_back(id[u]);}}
void build(int r, int n) {
 dfs(r);
 n = T;
 for(int i = n; i >= 1; i--) {
  for(int j = 0; j < rg[i].size(); j++) sdom[i] = min(sdom[i], sdom[find_(rg[i][j])]);
  if(i > 1) bucket[sdom[i]].push_back(i);
  for(int j = 0; j < bucket[i].size(); j++) {
   int w = bucket[i][j];
   int v = find_(w);
   if(sdom[v] == sdom[w]) idom[w] = sdom[w];
   else idom[w] = v;}
  if(i > 1) dsu[i] = par[i];}
 for(int i = 2; i <= n; i++) {
  if(idom[i] != sdom[i]) idom[i]=idom[idom[i]];
  t[rev[i]].push_back(rev[idom[i]]);
  t[rev[idom[i]]].push_back(rev[i]);}}
int st[N], en[N];
void yo(int u, int pre = 0) {
 st[u] = ++T;
 for(auto v: t[u]) {
  if(v == pre) continue;
  yo(v, u);}
 en[u] = T;}int main() {
 int n, m; while(cin >> n >> m) {
  vector<pair<int, int>> ed;
  for(int i = 0; i < m; i++) {
   int u, v; cin >> u >> v;
   g[u].push_back(v);
   ed.push_back({u, v});}
  build(1, n); T = 0; yo(1);
  vector<int> ans;
  for(int i = 0; i < m; i++) {
   int u = ed[i].first, v = ed[i].second;
   if(st[u] && !(st[v] <= st[u] && en[u] <= en[v])) ans.push_back(i);} yo(1);
  cout << ans.size() << '\n';
  for(auto x: ans) cout << x + 1 << ' ';
  cout << '\n';
  T = 0; for(int i = 0; i <= n; i++) {
   t[i].clear(), g[i].clear(), rg[i].clear(), bucket[i].clear();
   sdom[i] = par[i] = idom[i] = dsu[i] = label[i] = id[i] = rev[i] = st[i] = en[i] = 0;}}
 return 0;}
\end{minted}
\subsection{Dyn Connectivity}
\begin{minted}{c++}
struct persistent_dsu {
 struct state {
  int u, v, rnku, rnkv;
  state() {
   u = -1; v = -1;
   rnkv = -1; rnku = -1;}
  state(int _u, int _rnku, int _v, int _rnkv) {
   u = _u; rnku = _rnku; v = _v;
   rnkv = _rnkv;}}
; stack<state> st;
 int par[N], depth[N]; int comp;
 persistent_dsu() {
  comp = 0;
  memset(par, -1, sizeof(par));
  memset(depth, 0, sizeof(depth));} int root(int x) {
  if(x == par[x]) return x;
  return root(par[x]);} void init(int n) {
  comp = n;
  for(int i = 0; i <= n; i++) {
   par[i] = i; depth[i] = 1;}}
 bool connected(int x, int y) {
  return root(x) == root(y);} void unite(int x, int y) {
  int rx = root(x), ry = root(y);
  if(rx == ry) {
   st.push(state()); return;}
  if(depth[rx] < depth[ry])
   par[rx] = ry;
  else if(depth[ry] < depth[rx])
   par[ry] = rx;
  else {
   par[rx] = ry; depth[rx]++;}
  comp--; st.push(state(rx, depth[rx], ry, depth[ry]));}
 void backtrack(int c) {
  while(!st.empty() && c) {
   if(st.top().u == -1) {
    st.pop(); c--; continue;}
   par[st.top().u] = st.top().u;
   par[st.top().v] = st.top().v;
   depth[st.top().u] = st.top().rnku;
   depth[st.top().v] = st.top().rnkv;
   st.pop(); c--; comp++;}}};
persistent_dsu d;
vector<pair<int, int>> alive[4 * N];
void upd(int n, int b, int e, int i, int j, pair<int, int> &p) {
 if(b > j || e < i) return;
 if(b >= i && e <= j) {
  alive[n].push_back(p);/*this edge was alive in this time range */ return;}
 int l = 2 * n, r = l + 1, mid = b + e >> 1;
 upd(l, b, mid, i, j, p);
 upd(r, mid + 1, e, i, j, p);}
int ans[N];
void query(int n, int b, int e) {
 if(b > e) return;
 int prevsz = d.st.size();
 ///add edges which were alive in this range
 for(auto p : alive[n]) d.unite(p.first, p.second);
 if(b == e) {
  ans[b] = d.comp; d.backtrack(d.st.size() - prevsz); return;}
 int l = 2 * n, r = l + 1, mid = b + e >> 1;
 query(l, b, mid); query(r, mid + 1, e);
 d.backtrack(d.st.size() - prevsz);}
\end{minted}
\subsection{Dynamic MST Offline}
\begin{minted}{c++}
struct disj{
 int pa[MAXN], rk[MAXN];
 void init(int n){
  iota(pa, pa + n + 1, 0);
  memset(rk, 0, sizeof(rk));}
 int find(int x){
  return pa[x] == x ? x : find(pa[x]);}
 bool uni(int p, int q, vector<pair<int, int>> &snapshot){
  p = find(p);
  q = find(q);
  if(p == q) return 0;
  if(rk[p] < rk[q]) swap(p, q);
  snapshot.push_back({q, pa[q]});
  pa[q] = p;
  if(rk[p] == rk[q]){
   snapshot.push_back({p, -1});
   rk[p]++;
  } return 1;}
 void revert(vector<pair<int, int>> &snapshot){
  reverse(snapshot.begin(), snapshot.end());
  for(auto &x : snapshot){
   if(x.second < 0) rk[x.first]--;
   else pa[x.first] = x.second;
  } snapshot.clear();
 }}disj;
int n, m, q, st[MAXN], ed[MAXN], cost[MAXN], chk[MAXN];
pair<int, int> qr[MAXN];
bool cmp(int &a, int &b){ return pair<int, int>(cost[a], a) < pair<int, int>(cost[b], b); }
void contract(int s, int e, vector<int> v, vector<int> &must_mst, vector<int> &maybe_mst){
 sort(v.begin(), v.end(), cmp);
 vector<pair<int, int>> snapshot;
 for(int i=s; i<=e; i++) disj.uni(st[qr[i].first], ed[qr[i].first], snapshot);
 for(auto &i : v) if(disj.uni(st[i], ed[i], snapshot)) must_mst.push_back(i);
 disj.revert(snapshot);
 for(auto &i : must_mst) disj.uni(st[i], ed[i], snapshot);
 for(auto &i : v) if(disj.uni(st[i], ed[i], snapshot)) maybe_mst.push_back(i);
 disj.revert(snapshot);
}
void solve(int s, int e, vector<int> v, long long int cv){
 if(s == e){
  cost[qr[s].first] = qr[s].second;
  if(st[qr[s].first] == ed[qr[s].first]){
   printf("%lld\n", cv); return;}
  int minv = qr[s].second;
  for(auto &i : v) minv = min(minv, cost[i]);
  printf("%lld\n",minv + cv); return;}
 int m = (s+e)/2;
 vector<int> lv = v, rv = v;
 vector<int> must_mst, maybe_mst;
 for(int i=m+1; i<=e; i++){
  chk[qr[i].first]--;
  if(chk[qr[i].first] == 0) lv.push_back(qr[i].first);}
 vector<pair<int, int>> snapshot;
 contract(s, m, lv, must_mst, maybe_mst);
 long long int lcv = cv;
 for(auto &i : must_mst) lcv += cost[i], disj.uni(st[i], ed[i], snapshot);
 solve(s, m, maybe_mst, lcv);
 disj.revert(snapshot);
 must_mst.clear(); maybe_mst.clear();
 for(int i=m+1; i<=e; i++) chk[qr[i].first]++;
 for(int i=s; i<=m; i++){
  chk[qr[i].first]--;
  if(chk[qr[i].first] == 0) rv.push_back(qr[i].first);}
 long long int rcv = cv;
 contract(m+1, e, rv, must_mst, maybe_mst);
 for(auto &i : must_mst) rcv += cost[i], disj.uni(st[i], ed[i], snapshot);
 solve(m+1, e, maybe_mst, rcv);
 disj.revert(snapshot);
 for(int i=s; i<=m; i++) chk[qr[i].first]++;}
int main(){
 scanf("%d %d",&n,&m); vector<int> ve;
 for(int i=0; i<m; i++){
  scanf("%d %d %d",&st[i],&ed[i],&cost[i]);}
 scanf("%d",&q);
 for(int i=0; i<q; i++){
  scanf("%d %d",&qr[i].first,&qr[i].second);
  qr[i].first--; chk[qr[i].first]++;}
 disj.init(n);
 for(int i=0; i<m; i++) if(!chk[i]) ve.push_back(i);
 solve(0, q-1, ve, 0);}
\end{minted}
\subsection{Euler Path Directed}
\begin{minted}{c++}
const int N = 4e5 + 9;
/*all the edges should be in the same connected component
#directed graph: euler path: for all -> indeg = outdeg or nodes having indeg > outdeg = outdeg > indeg = 1 and for others in = out
#directed graph: euler circuit: for all -> indeg = outdeg*/
//euler path in a directed graph
//it also finds circuit if it exists
vector<int> g[N], ans;
int done[N];
void dfs(int u) {
 while (done[u] < g[u].size()) dfs(g[u][done[u]++]);
 ans.push_back(u);}
int solve(int n) {
 int edges = 0;
 vector<int> in(n + 1, 0), out(n + 1, 0);
 for (int u = 1; u <= n; u++) {
  for (auto v : g[u]) in[v]++, out[u]++, edges++;}
 int ok = 1, cnt1 = 0, cnt2 = 0, root = 0;
 for (int i = 1; i <= n; i++) {
  if (in[i] - out[i] == 1) cnt1++;
  if (out[i] - in[i] == 1) cnt2++, root = i;
  if (abs(in[i] - out[i]) > 1) ok = 0;}
 if (cnt1 > 1 || cnt2 > 1) ok = 0;
 if (!ok) return 0;
 if (root == 0) {
  for (int i = 1; i <= n; i++) if (out[i]) root = i;}
 if (root == 0) return 1; //empty graph
 dfs(root);
 if (ans.size() != edges + 1) return 0; //connectivity
 reverse(ans.begin(), ans.end());
 return 1;}
\end{minted}
\subsection{Euler Path Undirected}
\begin{minted}{c++}
/*all the edges should be in the same connected component
#undirected graph: euler path: all degrees are even or exactly two of them are odd.
#undirected graph: euler circuit: all degrees are even*/
//euler path in an undirected graph
//it also finds circuit if it exists
vector<pair<int, int>> g[N];
vector<int> ans;
int done[N];
int vis[N * N]; //number of edges
void dfs(int u) {
 while (done[u] < g[u].size()) {
  auto e = g[u][done[u]++];
  if (vis[e.second]) continue;
  vis[e.second] = 1;
  dfs(e.first);}
 ans.push_back(u);}
int solve(int n) {
 int edges = 0; ans.clear();
 memset(done, 0, sizeof done);
 memset(vis, 0, sizeof vis);
 vector<int> deg(n + 1, 0);
 for (int u = 1; u <= n; u++) {
  for (auto e : g[u]) {
   deg[e.first]++, deg[u]++, edges++;}}
 int odd = 0, root = 0;
 for (int i = 1; i <= n; i++) {
  if (deg[i] & 1) odd++, root = i;}
 if (odd > 2) return 0;
 if (root == 0) {
  for (int i = 1; i <= n; i++) if (deg[i]) root = i;}
 if (root == 0) return 1; //empty graph
 dfs(root);
 if (ans.size() != edges / 2 + 1) return 0;
 reverse(ans.begin(), ans.end());
 return 1;}int32_t main() {
  int n, m; cin >> n >> m;
  vector<int> deg(n + 1, 0);
  for (int i = 1; i <= m; i++) {
    int u, v; cin >> u >> v;
    g[u].push_back({v, i});
    g[v].push_back({u, i});
    deg[u]++, deg[v]++;}
  int sz = m;
  for (int i = 1; i <= n; i++) {
    if (deg[i] & 1) { ++sz;
    g[n + 1].push_back({i, sz});
    g[i].push_back({n + 1, sz});}}
  int ok = solve(n + 1);
  assert(ok);
  vector<int> in(n + 2, 0), out(n + 2, 0);
  for (int i = 0; i + 1 < ans.size(); i++) {
    if (ans[i] != n + 1 && ans[i + 1] != n + 1) {
    in[ans[i + 1]]++;
    out[ans[i]]++;}}
  int res = 0;
  for (int i = 1; i <= n; i++) res += in[i] == out[i];
  cout << res << '\n';
  for (int i = 0; i + 1 < ans.size(); i++) if (ans[i] != n + 1 && ans[i + 1] != n + 1) cout << ans[i] << ' ' << ans[i + 1] << '\n';
  return 0;}
\end{minted}
\subsection{Ford-Fulkerson}
\begin{minted}{c++}
//Ford-Fulkerson Edmonds-Karp : Max Flow
/*Directed: cap[x][y]=c; cap[y][x]=0;
  Undirected: cap[x][y]=c; cap[y][x]=c;
  Multiple Edge: cap[x][y]+=c; */
int bfs(int s, int t){
  memset(par,-1,sizeof par);
  par[s] = s; queue<pair<int, int>> q; q.push({s, 1e9});
  while (!q.empty()) {
    int cur=q.front().first; int flow=q.front().second;
    q.pop(); if (cur == t) return flow;
    for (int v : g[cur]) {
      if(par[v] == -1 && capacity[cur][v]) {
        par[v] = cur; 
        int new_flow = min(flow, capacity[cur][v]);
        q.push({v, new_flow}); } } } return 0; }
//O(V*E*E)
int maxflow(int s, int t) {
  int flow = 0,new_flow; while (1){
    new_flow = bfs(s, t); if(new_flow<=0) break;
    flow += new_flow; int cur = t; while (cur != s) {
      int prev=parent[cur]; capacity[prev][cur]-=new_flow;
      capacity[cur][prev] += new_flow; cur = prev;} } return flow; }
struct FlowEdge{
  int v,u; ll cap,flow=0LL;
  FlowEdge(int vv, int uu, ll c) {
    v=vv; u=uu; cap=c; /*Edge v to u*/ } };
\end{minted}
\subsection{HLD Path Query}
\begin{minted}{c++}
vector<int> g[N];
int par[N][LG + 1], dep[N], sz[N];
void dfs(int u, int p = 0) {
 par[u][0] = p;
 dep[u] = dep[p] + 1;
 sz[u] = 1;
 for (int i = 1; i <= LG; i++) par[u][i] = par[par[u][i - 1]][i - 1];
 if (p) g[u].erase(find(g[u].begin(), g[u].end(), p));
 for (auto &v : g[u]) if (v != p) {
   dfs(v, u);
   sz[u] += sz[v];
   if(sz[v] > sz[g[u][0]]) swap(v, g[u][0]);}}
int lca(int u, int v) {
 if (dep[u] < dep[v]) swap(u, v);
 for (int k = LG; k >= 0; k--) if (dep[par[u][k]] >= dep[v]) u = par[u][k];
 if (u == v) return u;
 for (int k = LG; k >= 0; k--) if (par[u][k] != par[v][k]) u = par[u][k], v = par[v][k];
 return par[u][0];}
int kth(int u, int k) {
 assert(k >= 0);
 for (int i = 0; i <= LG; i++) if (k & (1 << i)) u = par[u][i];
 return u;}
int T, head[N], st[N], en[N];
void dfs_hld(int u) {
 st[u] = ++T;
 for (auto v : g[u]) {
  head[v] = (v == g[u][0] ? head[u] : v);
  dfs_hld(v);}
 en[u] = T;}
int n;
int query_up(int u, int v) {
 int ans = -inf;
 while(head[u] != head[v]) {
  ans = max(ans, t.query(1, 1, n, st[head[u]], st[u]));
  u = par[head[u]][0];}
 ans = max(ans, t.query(1, 1, n, st[v], st[u]));
 return ans;}
\end{minted}
\subsection{HopCroft Karp}
\begin{minted}{c++}
//O(root(v) * E)
struct HopcroftKarp {
 static const int inf = 1e9;
 int n;
 vector<int> l, r, d;
 vector<vector<int>> g;
 HopcroftKarp(int _n, int _m) {
  n = _n;
  int p = _n + _m + 1;
  g.resize(p);
  l.resize(p, 0);
  r.resize(p, 0);
  d.resize(p, 0);}
 void add_edge(int u, int v) {
  g[u].push_back(v + n); //right id is increased by n, so is l[u]}
 bool bfs() {
  queue<int> q;
  for (int u = 1; u <= n; u++) {
   if (!l[u]) d[u] = 0, q.push(u);
   else d[u] = inf;}
  d[0] = inf;
  while (!q.empty()) {
   int u = q.front();
   q.pop();
   for (auto v : g[u]) {
    if (d[r[v]] == inf) {
     d[r[v]] = d[u] + 1;
     q.push(r[v]);}}}
  return d[0] != inf;}
 bool dfs(int u) {
  if (!u) return true;
  for (auto v : g[u]) {
   if(d[r[v]] == d[u] + 1 && dfs(r[v])) {
    l[u] = v; r[v] = u;
    return true;}}
  d[u] = inf;
  return false;}
 int maximum_matching() {
  int ans = 0;
  while (bfs()) {
   for(int u = 1; u <= n; u++) if (!l[u] && dfs(u)) ans++;}
  return ans;}}
\end{minted}
\subsection{Hungarian}
\begin{minted}{c++}
const int N = 509;
/* Complexity: O(n^3) but optimized
It finds minimum cost maximum matching.
For finding maximum cost maximum matching
add -cost and return -matching()
1-indexed */
struct Hungarian {
 long long c[N][N], fx[N], fy[N], d[N];
 int l[N], r[N], arg[N], trace[N];
 queue<int> q;
 int start, finish, n;
 const long long inf = 1e18;
 Hungarian() {}
 Hungarian(int n1, int n2): n(max(n1, n2)) {
  for (int i = 1; i <= n; ++i) {
   fy[i] = l[i] = r[i] = 0;
   for (int j = 1; j <= n; ++j) c[i][j] = inf;}}
 void add_edge(int u, int v, long long cost) { c[u][v] = min(c[u][v], cost); }
 inline long long getC(int u, int v) { return c[u][v] - fx[u] - fy[v]; }
 void initBFS() {
  while (!q.empty()) q.pop(); q.push(start);
  for (int i = 0; i <= n; ++i) trace[i] = 0;
  for (int v = 1; v <= n; ++v) { d[v] = getC(start, v); arg[v] = start; }
  finish = 0;}
 void findAugPath() {
  while (!q.empty()) {
   int u = q.front(); q.pop();
   for (int v = 1; v <= n; ++v) if (!trace[v]) {
     long long w = getC(u, v);
     if (!w) {
      trace[v] = u;
      if (!r[v]) { finish = v; return; }
      q.push(r[v]);}
     if (d[v] > w) { d[v] = w; arg[v] = u; }}}}
 void subX_addY() {
  long long delta = inf;
  for (int v = 1; v <= n; ++v) if (trace[v] == 0 && d[v] < delta) delta = d[v];
  fx[start] += delta;
  for (int v = 1; v <= n; ++v) if (trace[v]) {
    int u = r[v]; fy[v] -= delta; fx[u] += delta;} else d[v] -= delta;
  for (int v = 1; v <= n; ++v) if (!trace[v] && !d[v]) {
    trace[v] = arg[v];
    if (!r[v]) { finish = v; return; }
    q.push(r[v]);}}
 void Enlarge() {
  do {
   int u = trace[finish]; int nxt = l[u];
   l[u] = finish; r[finish] = u; finish = nxt;} while (finish);}
 long long maximum_matching() {
  for (int u = 1; u <= n; ++u) {
   fx[u] = c[u][1];
   for (int v = 1; v <= n; ++v) fx[u] = min(fx[u], c[u][v]);}
  for (int v = 1; v <= n; ++v) {
   fy[v] = c[1][v] - fx[1];
   for (int u = 1; u <= n; ++u) fy[v] = min(fy[v], c[u][v] - fx[u]);}
  for (int u = 1; u <= n; ++u) {
   start = u; initBFS();
   while (!finish) { findAugPath(); if (!finish) subX_addY(); }
   Enlarge();}
  long long ans = 0;
  for (int i = 1; i <= n; ++i) {
   if (c[i][l[i]] != inf) ans += c[i][l[i]];
   else l[i] = 0;}return ans;}};
\end{minted}
\subsection{Kuhn}
\begin{minted}{c++}
//O(n * m) 
struct Kuhn {
 int n; vector<vector<int>> g;
 vector<int> l, r; vector<bool> vis;
 Kuhn(int _n, int _m) {
  n = _n; g.resize(n + 1);
  vis.resize(n + 1, false);
  l.resize(n + 1, -1);
  r.resize(_m + 1, -1);}
 void add_edge(int a, int b) {
  g[a].push_back(b);}
 bool yo(int u) {
  if (vis[u]) return false;
  vis[u] = true;
  for (auto v : g[u]) {
   if (r[v] == -1 || yo(r[v])) {
    l[u] = v; r[v] = u;
    return true;}}
  return false;}
 int maximum_matching() {
  for (int i = 1; i <= n; i++) shuffle(g[i].begin(), g[i].end(), rnd);
  vector<int> p(n);
  iota(p.begin(), p.end(), 1);
  shuffle(p.begin(), p.end(), rnd);
  bool ok = true;
  while (ok) {
   ok = false; vis.assign(n + 1, false);
   for (auto &i : p) if(l[i] == -1) ok |= yo(i);}
  int ans = 0;
  for (int i = 1; i <= n; i++) ans += l[i] != -1;
  return ans;}};
bool vis[N];
vector<int> g[N], r[N], G[N], vec; //G is the condensed graph
void dfs1(int u) {
 vis[u] = 1;
 for(auto v : g[u]) if(!vis[v]) dfs1(v);
 vec.push_back(u);}vector<int> comp;
void dfs2(int u) {
 comp.push_back(u);
 vis[u] = 1;
 for(auto v : r[u]) if(!vis[v]) dfs2(v);}
int idx[N];// Given a graph having perfect matching, find the number of edges which are not in all Maximum Matchings in a Bipartite Graph
int32_t main() {
 int n, m; cin >> n >> m; Kuhn M(n, n);
 multiset<pair<int, int>> edges;
 for (int i = 1; i <= m; i++) {
  int u, v; cin >> u >> v; //u -> left side, v -> right side
  M.add_edge(u, v); edges.insert({u, v});}
 int cnt = M.maximum_matching();
 assert(cnt == n);
 for (int i = 1; i <= n; i++) {
  int u = i, v = M.l[i];
  edges.erase(edges.find({u, v})); //u -> v for matched edges
  g[u].push_back(v + n);
  r[v + n].push_back(u);}
 for (auto e : edges) {
  g[e.second + n].push_back(e.first); //v -> u for unmatched edges
  r[e.first].push_back(e.second + n);}
 for(int i = 1; i <= 2 * n; i++) if(!vis[i]) dfs1(i);
 reverse(vec.begin(), vec.end());
 memset(vis, 0, sizeof vis);
 int scc = 0; for(auto u : vec) {
  if(!vis[u]) {
   comp.clear(); dfs2(u); scc++;
   for(auto x : comp) idx[x] = scc;}}
 int ans = m - n;
 for (int i = 1; i <= n; i++) {
  int u = i, v = M.l[i];
  if (idx[u] == idx[v + n]) ans++; //if nodes of a matched edge are in the same SCC then we can replace it}
 cout << ans << '\n';
 return 0;}
\end{minted}
\subsection{LR flow}
\begin{minted}{c++}
struct LR_Flow {
 Dinic F;
 int n, s, t; struct edge {int u, v, l, r, id;};
 vector<edge> edges;
 LR_Flow() {}
 LR_Flow(int _n) {
 n = _n + 10; s = n - 2, t = n - 1;; edges.clear();}
 void add_edge(int u, int v, int l, int r, int id = -1) {
 assert(0 <= l && l <= r);
 edges.push_back({u, v, l, r, id});}
 bool feasible(int _s = -1, int _t = -1, int L = -1, int R = -1) {
 if (L != -1) edges.push_back({_t, _s, L, R, -1});
 F = Dinic(n); long long target = 0;
 for (auto e : edges) {
  int u = e.u, v = e.v, l = e.l, r = e.r, id = e.id;
  if (l != 0) {
  F.add_edge(s, v, l); F.add_edge(u, t, l); target += l;}
  F.add_edge(u, v, r - l, id);}
 auto ans = F.max_flow(s, t);
 if (L != -1) edges.pop_back();
 if (ans < target) return 0; //not feasible
 return 1;}
 int max_flow(int _s, int _t) { //-1 means flow is not feasible
 int mx = 1e5 + 9;
 if (!feasible(_s, _t, 0, mx)) return -1;
 return F.max_flow(_s, _t);}
 int min_flow(int _s, int _t) { //-1 means flow is not feasible
 int mx = 1e9; int ans = -1, l = 0, r = mx;
 while (l <= r) {
  int mid = l + r >> 1;
  if (feasible(_s, _t, 0, mid)) ans = mid, r = mid - 1;
  else l = mid + 1;} return ans;}};
int get_id(map<int, int> &mp, int k) {
 if (mp.find(k) == mp.end()) mp[k], mp[k] = mp.size();
 return mp[k];}
int Lx[N], Rx[N], Ly[N], Ry[N], degx[N], degy[N];
int32_t main() {
 cin >> n >> m; LR_Flow F(2 * n + 10);
 int r, b; cin >> r >> b; int sp = r > b;
 if (sp) swap(r, b); map<int, int> mx, my;
 for (int i = 1; i <= n; i++) {
 int x, y; cin >> x >> y; if (sp) swap(x, y);
 F.add_edge(get_id(mx, x), get_id(my, y) + n, 0, 1, i);
 degx[mx[x]]++;
 degy[my[y]]++;}
 for (int i = 1; i <= mx.size(); i++) Lx[i] = 0, Rx[i] = degx[i];
 for (int i = 1; i <= my.size(); i++) Ly[i] = 0, Ry[i] = degy[i];
 while (m--) {
 int ty, x, d; cin >> ty >> x >> d; ty--; ty ^= sp;
 if (ty == 0) {
  if (mx.find(x) != mx.end()) {
  int i = mx[x]; int p = degx[i];
  int l = (p - d + 1) / 2, r = (p + d) / 2;
  l = max(0, l); r = min(p, r);
  Lx[i] = max(Lx[i], l); Rx[i] = min(Rx[i], r);
  }} else {
  if (my.find(x) != my.end()) {
  int i = my[x]; int p = degy[i];
  int l = (p - d + 1) / 2, r = (p + d) / 2;
  l = max(0, l); r = min(p, r);
  Ly[i] = max(Ly[i], l); Ry[i] = min(Ry[i], r);}}}
 int s = 2 * n + 2, t = s + 1;
 for (int i = 1; i <= mx.size(); i++) {
 if (Lx[i] > Rx[i]) return cout << -1 << '\n', 0;
 F.add_edge(s, i, Lx[i], Rx[i]);}
 for (int i = 1; i <= my.size(); i++) {
 if (Ly[i] > Ry[i]) return cout << -1 << '\n', 0;
 F.add_edge(i + n, t, Ly[i], Ry[i]);}
 int c = F.max_flow(s, t);
 if (c == -1) return cout << -1 << '\n', 0;
 long long ans = 1LL * c * r + 1LL * (n - c) * b;
 cout << ans << '\n';
 for (int i = 1; i <= n; i++) {
 cout << "br"[F.F.flow_through[i] ^ sp];}
 cout << '\n';
 return 0;
}\end{minted}
\subsection{MIS}
\begin{minted}{c++}
/* Given a graph as a symmetric bitset matrix (without any self edges)
 * Finds the maximum clique
 * Can be used to find the maximum independent set by finding a clique of the complement graph.
 * Runs in about 1s for n=155, and faster for sparse graphs
 * 0 indexed*/
const int N = 105;
typedef vector<bitset<N>> graph;
#define sz(v) (int)v.size()
struct Maxclique {
 double limit = 0.025, pk = 0;
 struct Vertex {
  int i, d = 0;};
 typedef vector<Vertex> vv;
 graph e; vv V;
 vector<vector<int>> C;
 vector<int> qmax, q, S, old;
 void init(vv& r) {
  for (auto& v : r) v.d = 0;
  for (auto& v : r) for (auto j : r) v.d += e[v.i][j.i];
  sort(r.begin(), r.end(), [](auto a, auto b) {
   return a.d > b.d;});
  int mxD = r[0].d;
  for (int i = 0; i < sz(r); i++) r[i].d = min(i, mxD) + 1;}
 void expand(vv& R, int lev = 1) {
  S[lev] += S[lev - 1] - old[lev];
  old[lev] = S[lev - 1];
  while (sz(R)) {
   if (sz(q) + R.back().d <= sz(qmax)) return;
   q.push_back(R.back().i);
   vv T;
   for(auto v : R) if (e[R.back().i][v.i]) T.push_back({v.i});
   if (sz(T)) {
    if (S[lev]++ / ++pk < limit) init(T);
    int j = 0, mxk = 1, mnk = max(sz(qmax) - sz(q) + 1, 1);
    C[1].clear(), C[2].clear();
    for (auto v : T) {
     int k = 1;
     auto f = [&](int i) {
      return e[v.i][i];
     };
     while (any_of(C[k].begin(), C[k].end(), f)) k++;
     if (k > mxk) mxk = k, C[mxk + 1].clear();
     if (k < mnk) T[j++].i = v.i;
     C[k].push_back(v.i);}
    if (j > 0) T[j - 1].d = 0;
    for (int k = mnk; k <= mxk; k++) for (int i : C[k])
     T[j].i = i, T[j++].d = k;
    expand(T, lev + 1);
   } else if (sz(q) > sz(qmax)) qmax = q;
   q.pop_back(), R.pop_back();}}
 Maxclique(graph g) : e(g), C(sz(e) + 1), S(sz(C)), old(S) {
  for (int i = 0; i < sz(e); i++) V.push_back({i});}
 vector<int> solve() { // returns the clique
  init(V), expand(V); return qmax;}};
\end{minted}
\subsection{Min Cost Max Flow}
\begin{minted}{c++}
const int N = 3e5 + 9;
//Works for both directed, undirected and with negative cost too
//doesn't work for negative cycles
//for undirected edges just make the directed flag false
//Complexity: O(min(E^2 *V log V, E logV * flow))
using T = long long;
const T inf = 1LL << 61;
struct MCMF {
 struct edge {
  int u, v;
  T cap, cost;
  int id;
  edge(int _u, int _v, T _cap, T _cost, int _id) {
   u = _u; v = _v; cap = _cap;
   cost = _cost; id = _id;}};
 int n, s, t, mxid;
 T flow, cost;
 vector<vector<int>> g; vector<edge> e;
 vector<T> d, potential, flow_through;
 vector<int> par; bool neg;
 MCMF() {}
 MCMF(int _n) { // 0-based indexing
  n = _n + 10; g.assign(n, vector<int> ());
  neg = false; mxid = 0;}
 void add_edge(int u, int v, T cap, T cost, int id = -1, bool directed = true) {
  if(cost < 0) neg = true;
  g[u].push_back(e.size());
  e.push_back(edge(u, v, cap, cost, id));
  g[v].push_back(e.size());
  e.push_back(edge(v, u, 0, -cost, -1));
  mxid = max(mxid, id);
  if(!directed) add_edge(v, u, cap, cost, -1, true);}
 bool dijkstra() {
  par.assign(n, -1);
  d.assign(n, inf);
  priority_queue<pair<T, T>, vector<pair<T, T>>, greater<pair<T, T>> > q;
  d[s] = 0; q.push(pair<T, T>(0, s));
  while (!q.empty()) {
   int u = q.top().second;
   T nw = q.top().first; q.pop();
   if(nw != d[u]) continue;
   for (int i = 0; i < (int)g[u].size(); i++) {
    int id = g[u][i];
    int v = e[id].v;
    T cap = e[id].cap;
    T w = e[id].cost + potential[u] - potential[v];
    if (d[u] + w < d[v] && cap > 0) {
     d[v] = d[u] + w;
     par[v] = id;
     q.push(pair<T, T>(d[v], v));}}}
  for (int i = 0; i < n; i++) {
   if (d[i] < inf) d[i] += (potential[i] - potential[s]);}
  for (int i = 0; i < n; i++) {
   if (d[i] < inf) potential[i] = d[i];}
  return d[t] != inf; // for max flow min cost
  // return d[t] <= 0; // for min cost flow}
 T send_flow(int v, T cur) {
  if(par[v] == -1) return cur;
  int id = par[v]; int u = e[id].u;
  T w = e[id].cost; T f = send_flow(u, min(cur, e[id].cap));
  cost += f * w; e[id].cap -= f;
  e[id ^ 1].cap += f; return f;}
 //returns {maxflow, mincost}
 pair<T, T> solve(int _s, int _t, T goal = inf) {
  s = _s; t = _t;
  flow = 0, cost = 0;
  potential.assign(n, 0);
  if (neg) {
   // Run Bellman-Ford to find starting potential on the starting graph
   // If the starting graph (before pushing flow in the residual graph) is a DAG, 
   // then this can be calculated in O(V + E) using DP:
   // potential(v) = min({potential[u] + cost[u][v]}) for each u -> v and potential[s] = 0
   d.assign(n, inf);
   d[s] = 0;
   bool relax = true;
   for (int i = 0; i < n && relax; i++) {
    relax = false;
    for (int u = 0; u < n; u++) {
     for (int k = 0; k < (int)g[u].size(); k++) {
      int id = g[u][k];
      int v = e[id].v;
      T cap = e[id].cap, w = e[id].cost;
      if (d[v] > d[u] + w && cap > 0) {
       d[v] = d[u] + w;
       relax = true;}}}}
   for(int i = 0; i < n; i++) if(d[i] < inf) potential[i] = d[i];}
  while (flow < goal && dijkstra()) flow += send_flow(t, goal - flow);
  flow_through.assign(mxid + 10, 0);
  for (int u = 0; u < n; u++) {
   for (auto v : g[u]) {
    if (e[v].id >= 0) flow_through[e[v].id] = e[v ^ 1].cap;}}
  return make_pair(flow, cost);}};
int main() {
  int n; cin >> n;
 assert(n <= 10); MCMF F(2 * n);
 for (int i = 0; i < n; i++) {
  for (int j = 0; j < n; j++) {
   int k; cin >> k;
   F.add_edge(i, j + n, 1, k, i * 20 + j);}}
 int s = 2 * n + 1, t = s + 1;
 for (int i = 0; i < n; i++) {
  F.add_edge(s, i, 1, 0);
  F.add_edge(i + n, t, 1, 0);}
 auto ans = F.solve(s, t).second;
 long long w = 0;
 set<int> se;
 for (int i = 0; i < n; i++) {
  for (int j = 0; j < n; j++) {
   int p = i * 20 + j;
   if (F.flow_through[p] > 0) {
    se.insert(j);
    w += F.flow_through[p];}}}
 assert(se.size() == n && w == n);
 cout << ans << '\n';}
\end{minted}
\subsection{Min cost Vertex Cover}
\begin{minted}{c++}
// O(2^(m / 2) * m)
const int N = 42;
int g[N][N], c[N], cost[1 << 20], dp[1 << 20], need[1 << 20];
int32_t main() {
  int n, m; cin >> n >> m;
 int last;
 for (int i = 1; i <= n; i++) {
  int x; cin >> x; --x;
  if (i == 1) g[x][x] = 1; //self loop is okay, it means we must take x in our vertex cover
  if (i == n) g[x][x] = 1;
  if (i > 1) {
   g[last][x] = 1;
   g[x][last] = 1;}
  last = x;}
 // graph containing m nodes
 for (int i = 0; i < m; i++) {
  cin >> c[i];}
 // in this problem, answer for m = 1 is c[0], change this when necessary
 if (m == 1) {
  cout << c[0] << '\n';
  return 0;}
 int one = m / 2, two = m - m / 2;
 for (int mask = 0; mask < (1 << one); mask++) {
  bool good = true;
  for (int i = 0; i < one; i++) {
   for (int j = 0; j < m; j++) {
    if (g[i][j]) {
     if (j < one) {
      if (~mask >> i & 1 and ~mask >> j & 1) {
       good = false;}}
     else {
      if (~mask >> i & 1) {
       need[mask] |= 1 << (j - one);}}}}
   if (mask >> i & 1) {
    cost[mask] += c[i];}}
  if (!good) {
   cost[mask] = 1e9;}}
 for (int mask = 0; mask < (1 << two); mask++) {
  bool good = true;
  for (int i = 0; i < two; i++) {
   for (int j = 0; j < two; j++) {
    if (g[i + one][j + one]) {
     if (~mask >> i & 1 and ~mask >> j & 1) {
      good = false;}}}
   if (mask >> i & 1) {
    dp[mask] += c[i + one];}}
  if (!good) {
   dp[mask] = 1e9;}}
 for (int i = 0; i < two; i++) {
  for (int mask = (1 << two) - 1 ; mask >= 0 ; mask--) {
   if ((mask & (1 << i)) == 0) dp[mask] = min(dp[mask], dp[mask ^ (1 << i)]);}}
 int ans = 1e9;
 for (int mask = 0; mask < (1 << one); mask++) {
  ans = min(ans, cost[mask] + dp[need[mask]]);}
 cout << ans << '\n';
 return 0;}\end{minted}
\subsection{Prufer Code}
\begin{minted}{c++}
const int N = 3e5 + 9;/*
prufer code is a sequence of length n-2 to uniquely determine a labeled tree with n vertices
Each time take the leaf with the lowest number and add the node number the leaf is connected to
the sequence and remove the leaf. Then break the algo after n-2 iterations*/
//0-indexed
int n;vector<int> g[N];
int parent[N], degree[N];void dfs (int v) {
 for (size_t i = 0; i < g[v].size(); ++i) {
  int to = g[v][i];
  if (to != parent[v]) {
   parent[to] = v;
   dfs (to);}}}
vector<int> prufer_code() {
 parent[n - 1] = -1; dfs (n - 1); int ptr = -1;
 for (int i = 0; i < n; ++i) {
  degree[i] = (int) g[i].size();
  if (degree[i] == 1 && ptr == -1) ptr = i;}
 vector<int> result;
 int leaf = ptr;
 for (int iter = 0; iter < n - 2; ++iter) {
  int next = parent[leaf];
  result.push_back (next);
  --degree[next];
  if (degree[next] == 1 && next < ptr) leaf = next;
  else {
   ++ptr;
   while (ptr < n && degree[ptr] != 1) ++ptr;
   leaf = ptr;}}
 return result;}
vector < pair<int, int> > prufer_to_tree(const vector<int> & prufer_code) {
 int n = (int) prufer_code.size() + 2;
 vector<int> degree (n, 1);
 for (int i = 0; i < n - 2; ++i) ++degree[prufer_code[i]]; int ptr = 0;
 while (ptr < n && degree[ptr] != 1) ++ptr;
 int leaf = ptr; vector < pair<int, int> > result;
 for (int i = 0; i < n - 2; ++i) {
  int v = prufer_code[i];
  result.push_back (make_pair (leaf, v));
  --degree[leaf];
  if (--degree[v] == 1 && v < ptr) leaf = v;
  else {
   ++ptr;
   while (ptr < n && degree[ptr] != 1) ++ptr;
   leaf = ptr;}}
 for (int v = 0; v < n - 1; ++v) if (degree[v] == 1) result.push_back (make_pair (v, n - 1));
 return result;}
\end{minted}
\subsection{SCC}
\begin{minted}{c++}
// given a directed graph return the minimum number of edges to be added so that the whole graph become an SCC
bool vis[N];
vector<int> g[N], r[N], G[N], vec; //G is the condensed graph
void dfs1(int u) {
 vis[u] = 1;
 for(auto v: g[u]) if(!vis[v]) dfs1(v);
 vec.push_back(u);}vector<int> comp;
void dfs2(int u) {
 comp.push_back(u);
 vis[u] = 1;
 for(auto v: r[u]) if(!vis[v]) dfs2(v);}int idx[N], in[N], out[N];
int main() {
 int n, m; cin >> n >> m;
 for(int i = 1; i <= m; i++) {
  int u, v; cin >> u >> v;
  g[u].push_back(v); r[v].push_back(u);}
 for(int i = 1; i <= n; i++) if(!vis[i]) dfs1(i);
 reverse(vec.begin(), vec.end());
 memset(vis, 0, sizeof vis);
 int scc = 0;
 for(auto u: vec) {
  if(!vis[u]) {
   comp.clear(); dfs2(u); scc++;
   for(auto x: comp) idx[x]=scc;}}
 for(int u = 1; u <= n; u++) {
  for(auto v: g[u]) {
   if(idx[u] != idx[v]) {
    in[idx[v]]++, out[idx[u]]++;
    G[idx[u]].push_back(idx[v]);}}}
 int needed_in=0, needed_out=0;
 for(int i = 1; i <= scc; i++) {
  if(!in[i]) needed_in++;
  if(!out[i]) needed_out++;}
 int ans = max(needed_in, needed_out);
 if(scc == 1) ans = 0;
 cout << ans << '\n'; return 0;}
\end{minted}
\subsection{SPFA}
\begin{minted}{c++}
//SPFA is a improvement of the Bellman-Ford algorithm
vector< pair<int,int> > adj[MAX]; int d[MAX], cnt[MAX];
bool inqueue[MAX];
bool spfa(int n, int s){
  for(int i=1; i<=n; i++){
    d[i]=inf; inqueue[i]=0; cnt[i]=0; }
  queue<int> q; q.push(s); inqueue[s] = 1; d[s] = 0;
  while(!q.empty()){
  int u = q.front(); q.pop(); inqueue[u] = 0;
  for(auto edge : adj[u]){
  int v = edge.first; int w = edge.second;
  if(d[u] + w < d[v]){
    d[v] = d[u] + w;
    if(!inqueue[v]){
      q.push(v); inqueue[v] = 1; cnt[v]++; if(cnt[v] > n) return 0;/*negative cycle*/} } } }
    return 1;}
\end{minted}
\subsection{Stoer Wagner}
\begin{minted}{c++}
//O(n^3) but faster, 1 indexed
struct StoerWagner {
 int n; long long G[N][N], dis[N];
 int idx[N]; bool vis[N];
 const long long inf = 1e18; StoerWagner() {}
 StoerWagner(int _n) {
 n = _n; memset(G, 0, sizeof G);}
 void add_edge(int u, int v, long long w) { //undirected edge, multiple edges are merged into one edge
 if (u != v) {G[u][v] += w; G[v][u] += w;}}
 long long solve() {
 long long ans = inf; for (int i = 0; i < n; ++ i) idx[i] = i + 1;
 shuffle(idx, idx + n, rnd);
 while (n > 1) {
 int t = 1, s = 0; for (int i = 1; i < n; ++ i) {
 dis[idx[i]] = G[idx[0]][idx[i]];
 if (dis[idx[i]] > dis[idx[t]]) t = i;}
 memset(vis, 0, sizeof vis); vis[idx[0]] = true;
 for (int i = 1; i < n; ++ i) {
 if (i == n - 1) {
  if (ans > dis[idx[t]]) ans = dis[idx[t]]; //idx[s] - idx[t] is in two halves of the mincut
  if (ans == 0) return 0;
  for (int j = 0; j < n; ++ j) {
  G[idx[s]][idx[j]] += G[idx[j]][idx[t]];
  G[idx[j]][idx[s]] += G[idx[j]][idx[t]];}
  idx[t] = idx[-- n];}
 vis[idx[t]] = true;
 s = t; t = -1;
 for (int j = 1; j < n; ++ j) {
  if (!vis[idx[j]]) {
  dis[idx[j]] += G[idx[s]][idx[j]];
  if (t == -1 || dis[idx[t]] < dis[idx[j]]) t = j;}}}}
 return ans;}};
\end{minted}
\subsection{Tree Isomorphism}
\begin{minted}{c++}
struct Tree {
 int n; vector<vector<int>> g;
 Tree() {} Tree(int _n) : n(_n) {g.resize(n + 1);}
 void add_edge(int u, int v) {
 g[u].push_back(v); g[v].push_back(u);}
 vector<int> bfs(int s) {
 queue<int> q; vector<int> d(n + 1, n * 2);
 d[0] = -1; q.push(s); d[s] = 0;
 while(!q.empty()) {
  int u = q.front(); q.pop();
  for(auto v : g[u]) if(d[u] + 1 < d[v]) {
   d[v] = d[u] + 1; q.push(v);}}
 return d;}
 vector<int> get_centers() {
 auto du = bfs(1); int v = max_element(du.begin(), du.end()) - du.begin();
 auto dv = bfs(v);int u = max_element(dv.begin(), dv.end()) - dv.begin();
 du = bfs(u); vector<int> ans;
 for(int i = 1; i <= n; i++) if(du[i] + dv[i] == du[v] && du[i] >= du[v] / 2 && dv[i] >= du[v] / 2) {
  ans.push_back(i);}return ans;}
 mint yo(int u, int pre = 0) {
 vector<mint> nw; mint ans = 0;
 for(auto v : g[u]) if(v != pre) nw.push_back(yo(v, u));
 for(auto x : nw) ans = ans + P.pow(x.value);
 ans = ans * Q + R; return ans;}
 bool iso(Tree &t) {
 auto a = get_centers(); auto b = t.get_centers();
 for(auto x : a) for(auto y : b) if(yo(x) == t.yo(y)) return 1;
 return 0;}};
\end{minted}
\subsection{Virtual Tree}
\begin{minted}{c++}
const int N = 3e5 + 9;vector<int> g[N];
int par[N][20], dep[N], sz[N], st[N], en[N], T;
void dfs(int u, int pre) {
 par[u][0] = pre;
 dep[u] = dep[pre] + 1;
 sz[u] = 1;
 st[u] = ++T;
 for (int i = 1; i <= 18; i++) par[u][i] = par[par[u][i - 1]][i - 1];
 for (auto v : g[u]) {
  if (v == pre) continue; dfs(v, u);
  sz[u] += sz[v];} en[u] = T;}
int lca(int u, int v) {
 if (dep[u] < dep[v]) swap(u, v);
 for (int k = 18; k >= 0; k--) if (dep[par[u][k]] >= dep[v]) u = par[u][k];
 if (u == v) return u;
 for (int k = 18; k >= 0; k--) if (par[u][k] != par[v][k]) u = par[u][k], v = par[v][k];
 return par[u][0];}
int kth(int u, int k) {
 for (int i = 0; i <= 18; i++) if (k & (1 << i)) u = par[u][i];
 return u;}
int dist(int u, int v) {
 int lc = lca(u, v);
 return dep[u] + dep[v] - 2 * dep[lc];}
int isanc(int u, int v) {
 return (st[u] <= st[v]) && (en[v] <= en[u]);}
vector<int> t[N];
// given specific nodes, construct a compressed directed tree with these vertices(if needed some other nodes included)
// returns the nodes of the tree
// nodes.front() is the root
// t[] is the specific tree
vector<int> buildtree(vector<int> v) {
 // sort by entry time
 sort(v.begin(), v.end(), [](int x, int y) {
  return st[x] < st[y];});
 // finding all the ancestors, there are few of them
 int s = v.size();
 for (int i = 0; i < s - 1; i++) {
  int lc = lca(v[i], v[i + 1]);
  v.push_back(lc);}
 // removing duplicated nodes
 sort(v.begin(), v.end());
 v.erase(unique(v.begin(), v.end()), v.end());
 // again sort by entry time
 sort(v.begin(), v.end(), [](int x, int y) {
  return st[x] < st[y];});
 stack<int> st;
 st.push(v[0]);
 for (int i = 1; i < v.size(); i++) {
  while (!isanc(st.top(), v[i])) st.pop();
  t[st.top()].push_back(v[i]);
  st.push(v[i]);}
 return v;}
int ans; int imp[N];
int yo(int u) {
 vector<int> nw;
 for (auto v : t[u]) nw.push_back(yo(v));
 if (imp[u]) {
  for (auto x : nw) if (x) ans++;
  return 1;} else {
  int cnt = 0;
  for (auto x : nw) cnt += x > 0;
  if (cnt > 1) {
   ans++; return 0;}
  return cnt;}}
int32_t main() {
  int i, j, k, n, m, q, u, v; cin >> n;
 for (i = 1; i < n; i++) cin >> u >> v, g[u].push_back(v), g[v].push_back(u);
 dfs(1, 0); cin >> q;
 while (q--) {
  cin >> k; vector<int> v;
  for (i = 0; i < k; i++) cin >> m, v.push_back(m), imp[m] = 1;
  int fl = 1;
  for (auto x : v) if (imp[par[x][0]]) fl = 0;
  ans = 0; vector<int> nodes;
  if (fl) nodes = buildtree(v);
  if (fl) yo(nodes.front());
  if (!fl) ans = -1;
  cout << ans << '\n';
  // clear the tree
  for (auto x : nodes) t[x].clear();
  for (auto x : v) imp[x] = 0;}}
\end{minted}
\section{Math, Number Theory}
\subsection{BerleKamp Massey}
\begin{minted}{c++}
vector<mint> BerlekampMassey(vector<mint> S) {
 int n = (int)S.size(), L = 0, m = 0;
 vector<mint> C(n), B(n), T;
 C[0] = B[0] = 1; mint b = 1;
 for(int i = 0; i < n; i++) {
  ++m; mint d = S[i];
  for(int j = 1; j <= L; j++) d += C[j] * S[i - j];
  if (d == 0) continue;
  T = C; mint coef = d * b.inv();
  for(int j = m; j < n; j++) C[j] -= coef * B[j - m];
  if (2 * L > i) continue;
  L = i + 1 - L; B = T; b = d; m = 0;}
 C.resize(L + 1); C.erase(C.begin());
 for(auto &x: C) x *= -1;
 return C;}
vector<mint> combine (int n, vector<mint> &a, vector<mint> &b, vector<mint> &tr) {
 vector<mint> res(n * 2 + 1, 0);
 for (int i = 0; i < n + 1; i++) {
  for (int j = 0; j < n + 1; j++) res[i + j] += a[i] * b[j];}
 for (int i = 2 * n; i > n; --i) {
  for (int j = 0; j < n; j++) res[i - 1 - j] += res[i] * tr[j];}
 res.resize(n + 1); return res;
};
// transition -> for(i = 0; i < x; i++) f[n] += tr[i] * f[n-i-1]
// S contains initial values, k is 0 indexed
mint LinearRecurrence(vector<mint> &S, vector<mint> &tr, long long k) {
 int n = S.size(); assert(n == (int)tr.size());
 if (n == 0) return 0;
 if (k < n) return S[k];
 vector<mint> pol(n + 1), e(pol);
 pol[0] = e[1] = 1;
 for (++k; k; k /= 2) {
  if (k % 2) pol = combine(n, pol, e, tr); e = combine(n, e, e, tr); }
 mint res = 0;
 for (int i = 0; i < n; i++) res += pol[i + 1] * S[i]; return res;}
int prime[] = {2, 3, 5, 7, 11, 13, 17, 19}, ok[20];
int dp[2000][20];
int yo(int i, int last) {
 if (i == 0) return 1; int &ret = dp[i][last];
 if (ret != -1) return ret;
 ret = 0;
 for (int k = 1; k <= 9; k++) {
  if (!last || ok[last + k]) ret = (ret + yo(i - 1, k)) % mod; }return ret;}
int32_t main() {
 memset(dp, -1, sizeof dp);
 for (int i = 0; i < 8; i++) ok[prime[i]] = 1;
 vector<mint> S; S.push_back(4); mint sum = 4;
 for (int i = 2; i < 100; i++) sum += yo(i, 0), S.push_back(sum);
 auto tr = BerlekampMassey(S);
 S.resize((int)tr.size());
 int q; cin >> q;
 while (q--) {
  int n; cin >> n; --n;
  cout << LinearRecurrence(S, tr, n) << '\n';}
 return 0;}
\end{minted}
\subsection{CRT}
\begin{minted}{c++}
#define __int128 lll
/*A CRT solver which works even when moduli are not pairwise
coprime Call CRT(k) to get {x, N} pair, where x is the 
unique solution modulo N. O(n*Log L) L=LCM(p1,p2,....) 
Assumptions: 1. LCM of all mods will fit into long long.*/
pll CRT(ll n){
  ll r1,r2,p1,p2,x,y,ans,g,mod;  r1=r[0],p1=p[0];
  for(int i=1; i<n; i++){
    r2=r[i],p2=p[i];  g=__gcd(p1,p2);
    if(r1%g != r2%g) return {-1, -1}; //no solution
    EGCD(p1/g, p2/g, x, y);  mod=p1/g*p2;
ans=((lll)r1*(p2/g)%mod*y%mod+(lll)r2*(p1/g)%mod*x%mod)%mod;
    r1=ans;  if(r1<0) r1+=mod;  p1=mod;}  return {r1,mod};  }
\end{minted}
\subsection{Discrete Log}
\begin{minted}{c++}
//returns minimum integer x such that a^x = b (mod m)
// a and m are co-prime, O(sqrt(m))
int discrete_log(int a, int b, int m) {
  int n=(int)sqrt(m+.0)+1, pw=1;
  for(int i=0; i<n; ++i) pw=(1LL*pw*a)%m;
  gp_hash_table<int, int> vals;
  for(int p=1, cur=pw; p<=n; ++p){
    if(!vals[cur]) vals[cur] = p;
    cur=(1LL*cur*pw)%m; }
  int ans = inf;
  for(int q=0, cur=b; q<=n; ++q){
    if(vals.find(cur) != vals.end()){
      ll nw=1LL*vals[cur]*n-q;
      if(nw<ans) ans = nw;  }
    cur = (1LL * cur * a) % m;  }
  if(ans == inf) ans = -1;  return ans; }
ll inverse(ll a, ll m) {
  ll x, y; ll g = e_gcd(a, m, x, y);
  if (g != 1) return -1;return (x % m + m) % m;}
// discrete log but a and m may not be co-prime
int discrete_log_noncoprime(int a, int b, int m) {
  if(m == 1 || b == 1) return 0;
  if(__gcd(a, m) == 1) return discrete_log(a, b, m);
  int g = __gcd(a, m); if(b % g != 0)  return -1;
  int p = inverse(a / g, m / g);
  int nw = discrete_log_noncoprime(a, 1LL * b / g * p % (m / g), m / g);
  if (nw == -1) return -1;  return nw + 1; }
\end{minted}
\subsection{Discrete and Primitive Root}
\begin{minted}{c++}
/*Finds the primitive root modulo p. g is a primitive root 
mod p if and only if for any integer a such that gcd(a,p)
=1, there exists an integer k such that: g^k = a(mod p)*/
int PrimitiveRoot (int p) {
  vector<int> fact;
  int phi = p-1; // if p is prime
  int n = phi;
  for(int i=2; i*i<=n; ++i){
    if(n%i == 0){
      fact.push_back (i);  while(n % i == 0)  n /= i; }}
  if(n>1) fact.push_back (n);
  for(int res=2; res<=p; ++res){
    bool ok = true;
    for(size_t i=0; i<fact.size() && ok; ++i)
      ok &= BigMod (res, phi/fact[i], p) != 1;
    if (ok)  return res;}  return -1;  }
//find all numbers x such that x^k = a (mod n)
void printDiscreteRoot(int k, int a, int n){
  if(a==0){ cout<<"1\n0\n"; return;  }
  int g=PrimitiveRoot(n);
  int phi=n-1; //if n is not a prime calculate phi(n)
  int sq=(int) sqrt(n+0.0)+1;
  vector<pair<int,int> > dec (sq);
  for(int i=1; i<=sq; ++i)
    dec[i-1]={BigMod(g, 1LL*i*sq%phi*k%phi,n), i};
  sort(dec.begin(), dec.end());
  int any_ans = -1;
  for(int i=0; i<sq; ++i) {
    int my=BigMod(g, 1LL*i*k%phi, n)*1LL*a%n;
 auto it=lower_bound(dec.begin(),dec.end(),make_pair(my,0));
    if(it != dec.end() && it->first == my){
      any_ans = it->second * sq - i;  break;  }   }
  if(any_ans==-1){  cout<<"0\n"; return ;   }
  // all possible answers
  int delta = (n-1) / __gcd(k, n-1);
  vector<int> ans;
  for(int cur=any_ans%delta; cur<n-1; cur+=delta)
    ans.push_back(powmod(g,cur,n));
  sort(ans.begin(), ans.end());
  cout<<(int)ans.size()<<"\n";
  for(int answer : ans) cout<<answer<<" ";  }
\end{minted}
\subsection{E gcd}
\begin{minted}{c++}
ll EGCD(ll a, ll b, ll &x, ll &y){
  if(a==0){  x=0; y=1; return b; }
  ll x1,y1;  ll d = EGCD(b%a, a, x1, y1);
  x = y1 - (b/a)*x1;  y = x1;
  return d; }
\end{minted}
\subsection{FFT}
\begin{minted}{c++}
const double PI = acos(-1);
struct base {
  double a, b;
  base(double a = 0, double b = 0) : a(a), b(b) {}
  const base operator + (const base &c) const
    { return base(a + c.a, b + c.b); }
  const base operator - (const base &c) const
    { return base(a - c.a, b - c.b); }
  const base operator * (const base &c) const
    { return base(a * c.a - b * c.b, a * c.b + b * c.a); }};
void fft(vector<base> &p, bool inv = 0) {
  int n = p.size(), i = 0;
  for(int j = 1; j < n - 1; ++j) {
    for(int k = n >> 1; k > (i ^= k); k >>= 1);
    if(j < i) swap(p[i], p[j]);}
  for(int l = 1, m; (m = l << 1) <= n; l <<= 1) {
    double ang = 2 * PI / m;
    base wn = base(cos(ang), (inv ? 1. : -1.) * sin(ang)), w;
    for(int i = 0, j, k; i < n; i += m) {
      for(w = base(1, 0), j = i, k = i + l; j < k; ++j, w = w * wn) {
        base t = w * p[j + l];
        p[j + l] = p[j] - t;
        p[j] = p[j] + t;}}}
  if(inv) for(int i = 0; i < n; ++i) p[i].a /= n, p[i].b /= n;}
vector<long long> multiply(vector<int> &a, vector<int> &b) {
  int n = a.size(), m = b.size(), t = n + m - 1, sz = 1;
  while(sz < t) sz <<= 1;
  vector<base> x(sz), y(sz), z(sz);
  for(int i = 0 ; i < sz; ++i) {
    x[i] = i < (int)a.size() ? base(a[i], 0) : base(0, 0);
    y[i] = i < (int)b.size() ? base(b[i], 0) : base(0, 0);}
  fft(x), fft(y);
  for(int i = 0; i < sz; ++i) z[i] = x[i] * y[i];
  fft(z, 1);
  vector<long long> ret(sz);
  for(int i = 0; i < sz; ++i) ret[i] = (long long) round(z[i].a);
  while((int)ret.size() > 1 && ret.back() == 0) ret.pop_back();
  return ret;}
\end{minted}
\subsection{FWHT}
\begin{minted}{c++}
int POW(long long n, long long k) {
 int ans = 1 % mod; n %= mod; if (n < 0) n += mod;
 while (k) {
  if (k & 1) ans = (long long) ans * n % mod;
  n = (long long) n * n % mod;
  k >>= 1;}return ans;}
const int inv2 = (mod + 1) >> 1;
#define M (1 << 20)
#define OR 0
#define AND 1
#define XOR 2
struct FWHT{
 int P1[M], P2[M];
 void wt(int *a, int n, int flag = XOR) {
  if (n == 0) return;
  int m = n / 2;
  wt(a, m, flag); wt(a + m, m, flag);
  for (int i = 0; i < m; i++){
   int x = a[i], y = a[i + m];
   if (flag == OR) a[i] = x, a[i + m] = (x + y) % mod;
   if (flag == AND) a[i] = (x + y) % mod, a[i + m] = y;
   if (flag == XOR) a[i] = (x + y) % mod, a[i + m] = (x - y + mod) % mod;}}
 void iwt(int* a, int n, int flag = XOR) {
  if (n == 0) return;
  int m = n / 2;
  iwt(a, m, flag); iwt(a + m, m, flag);
  for (int i = 0; i < m; i++){
   int x = a[i], y = a[i + m];
   if (flag == OR) a[i] = x, a[i + m] = (y - x + mod) % mod;
   if (flag == AND) a[i] = (x - y + mod) % mod, a[i + m] = y;
   if (flag == XOR) a[i] = 1LL * (x + y) * inv2 % mod, a[i + m] = 1LL * (x - y + mod) * inv2 % mod; // replace inv2 by >>1 if not required}}
 vector<int> multiply(int n, vector<int> A, vector<int> B, int flag = XOR) {
  assert(__builtin_popcount(n) == 1);
  A.resize(n); B.resize(n);
  for (int i = 0; i < n; i++) P1[i] = A[i];
  for (int i = 0; i < n; i++) P2[i] = B[i];
  wt(P1, n, flag); wt(P2, n, flag);
  for (int i = 0; i < n; i++) P1[i] = 1LL * P1[i] * P2[i] % mod;
  iwt(P1, n, flag);
  return vector<int> (P1, P1 + n);}
 vector<int> pow(int n, vector<int> A, long long k, int flag = XOR) {
  assert(__builtin_popcount(n) == 1);
  A.resize(n);
  for (int i = 0; i < n; i++) P1[i] = A[i];
  wt(P1, n, flag);
  for(int i = 0; i < n; i++) P1[i] = POW(P1[i], k);
  iwt(P1, n, flag);
  return vector<int> (P1, P1 + n);}}t;
\end{minted}
\subsection{Gaussian Elimination}
\begin{minted}{c++}
const double eps = 1e-9;
int Gauss(vector<vector<double>> a, vector<double> &ans) {
 int n = (int)a.size(), m = (int)a[0].size() - 1;
 vector<int> pos(m, -1);
 double det = 1; int rank = 0;
 for(int col = 0, row = 0; col < m && row < n; ++col) {
  int mx = row;
  for(int i = row; i < n; i++) if(fabs(a[i][col]) > fabs(a[mx][col])) mx = i;
  if(fabs(a[mx][col]) < eps) {det = 0; continue;}
  for(int i = col; i <= m; i++) swap(a[row][i], a[mx][i]);
  if (row != mx) det = -det;
  det *= a[row][col];
  pos[col] = row;
  for(int i = 0; i < n; i++) {
   if(i != row && fabs(a[i][col]) > eps) {
    double c = a[i][col] / a[row][col];
    for(int j = col; j <= m; j++) a[i][j] -= a[row][j] * c;}} ++row; ++rank;}
 ans.assign(m, 0);
 for(int i = 0; i < m; i++) {
  if(pos[i] != -1) ans[i] = a[pos[i]][m] / a[pos[i]][i];}
 for(int i = 0; i < n; i++) {
  double sum = 0;
  for(int j = 0; j < m; j++) sum += ans[j] * a[i][j];
  if(fabs(sum - a[i][m]) > eps) return -1; //no solution}
 for(int i = 0; i < m; i++) if(pos[i] == -1) return 2; //infinte solutions
 return 1; //unique solution}
\end{minted}
\subsection{Linear Diophantine Eqn}
\begin{minted}{c++}
void shift_solution(int &x,int &y,int a,int b,int cnt){
  x+=cnt*b; y-=cnt*a; }
void findPosSol(int a,int b,int &x,int &y,int g){
  a/=g; b/=g; int mov;
  if(x<0){mov=(0-x+b-1)/b;shift_solution(x,y,a,b,mov);}
  if(y<0) { mov=(-y)/a; if((-y)%a) mov++;
    shift_solution(x,y,a,b,-mov); } }
bool find_any_solution(int a,int b,int c,int &x0,int &y0,int &g){
  g=egcd(abs(a),abs(b),x0,y0); if(c%g) return 0;
  x0*=c/g; y0*=c/g; if(a<0) x0*=-1; if(b<0) y0*=-1;
  /*findPosSol(a,b,x0,y0,g);*/ return 1; }
int find_all_solutions(int a, int b, int c, int minX, int maxX, int minY, int maxY){
  if(!a && !b){ if(c) return 0LL;
    return (maxX-minX+1LL)*(maxY-minY+1LL); }
  if(!a){ if(c%b) return 0LL; int y=c/b;
    if(minY<=y && y<=maxY) return (maxX-minX+1LL);
    return 0LL; }
  if(!b){ if(c%a) return 0LL; int x=c/a;
    if(minX<=x && x<=maxX) return (maxY-minY+1LL);
    return 0LL; }
  int x, y, g;
  if(!find_any_solution(a,b,c,x,y,g)) return 0;
  a/=g; b/=g;
  int sign_a = a > 0 ? +1 : -1;
  int sign_b = b > 0 ? +1 : -1;
  shift_solution(x, y, a, b,(minX-x)/b);
  if(x<minX) shift_solution(x, y, a, b, sign_b);
  if(x>maxX)  return 0;
  int lx1 = x; shift_solution(x,y,a,b,(maxX-x)/b);
  if(x>maxX) shift_solution(x, y, a, b, -sign_b);
  int rx1 = x; shift_solution(x,y,a,b, -(minY-y)/a);
  if(y<minY) shift_solution(x, y, a, b, -sign_a);
  if(y>maxY)  return 0; int lx2 = x;
  shift_solution(x, y, a, b, -(maxY-y)/a);
  int rx2 = x; if(lx2>rx2) swap(lx2, rx2);
  int lx = max(lx1, lx2), rx = min(rx1, rx2);
  if(lx>rx) return 0; return (rx-lx)/abs(b) + 1;}
/** LD Less or Equal: ax+by <= c **/
ll sumsq(ll n) { return n / 2 * ((n - 1) | 1);}
//\sum_{i = 0}^{n - 1}{(a + d * i) / m}, O(log m)
ll floor_sum(ll a, ll d, ll m, ll n) {
  if(d < 0) { //if d<0, reversing the series
    a = a + (n - 1) * d; d = -d; }
  ll res = d / m * sumsq(n) + a / m * n;
  d %= m; a %= m; if (!d) return res;
  ll to = (n * d + a) / m;
  return res+(n-1)*to-floor_sum(m-1-a,m,d,to);}
// number of integer solutions to ax + by <= c s.t. x, y >= 0
ll lattice_cnt(ll a, ll b, ll c) { // log(a)
  if (a == 0 or b == 0) { // infinite solutions
    return -1; }
  ll terms = 1 + c / b;
  return terms + floor_sum(c, -b, a , terms); }
\end{minted}
\subsection{NOD SOD POD}
\begin{minted}{c++}
/*{p^a} in prime factorization
nod=nod*(a+1)%MOD; 
sod=sod*(BigMod(p,a+1)-1+MOD)%MOD*inverse(p-1,MOD)%MOD;
pod=BigMod(pod,a+1)*BigMod(BigMod(p,(a*(a+1)/2)),tnod)%MOD;
tnod=tnod*(a+1)%(MOD-1);*/
//O(N^1/3),need prime till n^1/3
int CB(int x) {return x*x*x;}
int NOD(int n){
  int cnt=1;
  for(int i=0; CB(prime[i])<=n; i++){
    int c=1;
    while(n&prime[i]==0) { c++; n/=prime[i]; } cnt*=c; }
  if(isPrime[n]) cnt*=2; else if(isPrimeSq[n]) cnt*=3;
  else if(n!=1) cnt*=4;  return cnt; }
//O(sqrt(n)), NOD(1)+..+NOD(n)
ll CNOD(ll n) {
  ll i,cnt=0,sq=sqrt(n);
  for(i=1; i<=sq; i++) cnt+=(n/i)-i;
  cnt*=2; cnt+=sq;  return cnt; }
//O(sqrt(n))
ll CSOD(ll n){
  ll i,j,ans=0LL;
  for(i=1; i*i<=n; i++){
    j=n/i; ans+=(i+j)*(j-i+1)/2LL;
    ans+=i*(j-i); } return ans; }
//O(NLogN)
void SOD(){
  for(i=1; i<MAX; i++) for(j=i; j<MAX; j+=i) sod[j]+=i; }
\end{minted}
\subsection{NTT With Any Prime}
\begin{minted}{c++}
struct base {
 double x, y;
 base() { x = y = 0; }
 base(double x, double y): x(x), y(y) { }};
inline base operator + (base a, base b) { return base(a.x + b.x, a.y + b.y); }
inline base operator - (base a, base b) { return base(a.x - b.x, a.y - b.y); }
inline base operator * (base a, base b) { return base(a.x * b.x - a.y * b.y, a.x * b.y + a.y * b.x); }
inline base conj(base a) { return base(a.x, -a.y); }
int lim = 1;
vector<base> roots = {{0, 0}, {1, 0}};
vector<int> rev = {0, 1};
const double PI = acosl(- 1.0);
void ensure_base(int p) {
 if(p <= lim) return;
 rev.resize(1 << p);
 for(int i = 0; i < (1 << p); i++) rev[i] = (rev[i >> 1] >> 1) + ((i & 1) << (p - 1));
 roots.resize(1 << p);
 while(lim < p) {
  double angle = 2 * PI / (1 << (lim + 1));
  for(int i = 1 << (lim - 1); i < (1 << lim); i++) {
   roots[i << 1] = roots[i];
   double angle_i = angle * (2 * i + 1 - (1 << lim));
   roots[(i << 1) + 1] = base(cos(angle_i), sin(angle_i));}
  lim++;}}
void fft(vector<base> &a, int n = -1){
 if(n == -1) n = a.size();
 assert((n & (n - 1)) == 0);
 int zeros = __builtin_ctz(n);
 ensure_base(zeros);
 int shift = lim - zeros;
 for(int i = 0; i < n; i++) if(i < (rev[i] >> shift)) swap(a[i], a[rev[i] >> shift]);
 for(int k = 1; k < n; k <<= 1) {
  for(int i = 0; i < n; i += 2 * k) {
   for(int j = 0; j < k; j++) {
    base z = a[i + j + k] * roots[j + k];
    a[i + j + k] = a[i + j] - z;
    a[i + j] = a[i + j] + z;}}}}
//eq = 0: 4 FFTs in total
//eq = 1: 3 FFTs in total
vector<int> multiply(vector<int> &a, vector<int> &b, int eq = 0) {
 int need = a.size() + b.size() - 1;
 int p = 0;
 while((1 << p) < need) p++;
 ensure_base(p);
 int sz = 1 << p;
 vector<base> A, B;
 if(sz > (int)A.size()) A.resize(sz);
 for(int i = 0; i < (int)a.size(); i++) {
  int x = (a[i] % mod + mod) % mod;
  A[i] = base(x & ((1 << 15) - 1), x >> 15);}
 fill(A.begin() + a.size(), A.begin() + sz, base{0, 0});
 fft(A, sz);
 if(sz > (int)B.size()) B.resize(sz);
 if(eq) copy(A.begin(), A.begin() + sz, B.begin());
 else {
  for(int i = 0; i < (int)b.size(); i++) {
   int x = (b[i] % mod + mod) % mod;
   B[i] = base(x & ((1 << 15) - 1), x >> 15);}
  fill(B.begin() + b.size(), B.begin() + sz, base{0, 0});
  fft(B, sz);}
 double ratio = 0.25 / sz;
 base r2(0, - 1), r3(ratio, 0), r4(0, - ratio), r5(0, 1);
 for(int i = 0; i <= (sz >> 1); i++) {
  int j = (sz - i) & (sz - 1);
  base a1 = (A[i] + conj(A[j])), a2 = (A[i] - conj(A[j])) * r2;
  base b1 = (B[i] + conj(B[j])) * r3, b2 = (B[i] - conj(B[j])) * r4;
  if(i != j) {
   base c1 = (A[j] + conj(A[i])), c2 = (A[j] - conj(A[i])) * r2;
   base d1 = (B[j] + conj(B[i])) * r3, d2 = (B[j] - conj(B[i])) * r4;
   A[i] = c1 * d1 + c2 * d2 * r5;
   B[i] = c1 * d2 + c2 * d1;}
  A[j] = a1 * b1 + a2 * b2 * r5;
  B[j] = a1 * b2 + a2 * b1;}
 fft(A, sz); fft(B, sz);
 vector<int> res(need);
 for(int i = 0; i < need; i++) {
  long long aa = A[i].x + 0.5;
  long long bb = B[i].x + 0.5;
  long long cc = A[i].y + 0.5;
  res[i] = (aa + ((bb % mod) << 15) + ((cc % mod) << 30))%mod;}return res;}
\end{minted}
\subsection{NTT}
\begin{minted}{c++}
const int mod = 998244353;
inline int mul(int a, int b) {return (a * 1ll * b) % mod;}
inline int add(int a, int b) {a += b; if (a >= mod)a -= mod; return a;}
inline int sub(int a, int b) {a -= b; if (a < 0)a += mod; return a;}
inline int power(int a, int b) {int rt = 1; while (b > 0) {if (b & 1)rt = mul(rt, a); a = mul(a, a); b >>= 1;} return rt;}
inline int inv(int a) {return power(a, mod - 2);}
inline void modadd(int &a, int &b) {a += b; if (a >= mod)a -= mod;}
int base = 1, root = -1, max_base = -1;
vector<int> rev = {0, 1}, roots = {0, 1};
void init() {
 int temp = mod - 1;
 max_base = 0;
 while (temp % 2 == 0) {
  temp >>= 1;
  ++max_base;}
 root = 2;
 while (true) {if (power(root, 1 << max_base) == 1 && power(root, 1 << max_base - 1) != 1) {break;}++root;}}
void ensure_base(int nbase) {
 if (max_base == -1) {init();}
 if (nbase <= base) {return;}
 assert(nbase <= max_base);
 rev.resize(1 << nbase);
 for (int i = 0; i < (1 << nbase); i++) {
  rev[i] = (rev[i >> 1] >> 1) + ((i & 1) << (nbase - 1));}
 roots.resize(1 << nbase);
 while (base < nbase) {
  int z = power(root, 1 << (max_base - 1 - base));
  for (int i = 1 << (base - 1); i < (1 << base); i++) {
   roots[i << 1] = roots[i];
   roots[(i << 1) + 1] = mul(roots[i], z);}
  base++;}}
void fft(vector<int> &a) {
 int n = (int) a.size();
 assert((n & (n - 1)) == 0);
 int zeros = __builtin_ctz(n);
 ensure_base(zeros);
 int shift = base - zeros;
 for (int i = 0; i < n; i++) {
  if (i < (rev[i] >> shift)) {
   swap(a[i], a[rev[i] >> shift]);}}
 for (int k = 1; k < n; k <<= 1) {
  for (int i = 0; i < n; i += 2 * k) {
   for (int j = 0; j < k; j++) {
    int x = a[i + j];
    int y = mul(a[i + j + k], roots[j + k]);
    a[i + j] = x + y - mod;
    if (a[i + j] < 0) a[i + j] += mod;
    a[i + j + k] = x - y + mod;
    if (a[i + j + k] >= mod) a[i + j + k] -= mod;}}}}
vector<int> multiply(vector<int> a, vector<int> b, int eq = 0) {
 int need = (int) (a.size() + b.size() - 1);
 int nbase = 0;
 while ((1 << nbase) < need) nbase++;
 ensure_base(nbase);
 int sz = 1 << nbase;
 a.resize(sz);
 b.resize(sz);
 fft(a);
 if (eq) b = a; else fft(b);
 int inv_sz = inv(sz);
 for (int i = 0; i < sz; i++) {
  a[i] = mul(mul(a[i], b[i]), inv_sz);}
 reverse(a.begin() + 1, a.end());
 fft(a);
 a.resize(need);
 return a;}
\end{minted}
\subsection{Pollard Rho Miller Rabin}
\begin{minted}{c++}
using ll = uint64_t; using u128 = __uint128_t;
bool isComposite(ll n,ll a,ll d,int s){
  ll x=BigMod(a, d, n);
  if(x==1LL || x==n-1) return 0;
  for(int r=1; r<s; r++){
    x = ( (u128)x*x)%n;
    if(x==n-1) return 0;  }  return 1; }
bool MillerRabin(ll n){
  int s=0; ll d=n-1;
  while((d&1) == 0){  d >>= 1; s++; }
  for(ll a : {2,3,5,7,11,13,17,19,23,29,31,37}){
    if(n==a) return 1;
    if(isComposite(n, a, d, s)) return 0;} return 1; }
bool isPrime(ll n){
  if(n==2LL || n==3LL) return 1;
  if(n<5LL || n%2==0 || n%3==0) return 0;
  return MillerRabin(n); }
int pollard_rho(int n){
  if(n%2==0) return 2;
  int x = rand()%n+1; int c = rand()%n+1;
  int y=x, g=1;
  while(g==1){
    x=((x*x)%n+c)%n; y=((y*y)%n+c)%n;
    y=((y*y)%n+c)%n; g= __gcd(abs(x-y),n);} return g; }
vector<int>fact;
//O(N^1/4)
void factorize(int n){ if(n==1) return;
  if(isPrime(n)) { fact.push_back(n); return;}
  int divisor=pollard_rho(n);
  factorize(divisor);factorize(n/divisor); }
\end{minted}
\subsection{Prime Counting Function}
\begin{minted}{c++}
namespace pcf {
 // initialize once by calling init()
 #define MAXN 20000010 // initial sieve limit
 #define MAX_PRIMES 2000010 // max size of the prime array for sieve
 #define PHI_N 100000 #define PHI_K 100
 int len = 0; // total number of primes generated by sieve
 int primes[MAX_PRIMES];
 int pref[MAXN]; // pref[i] --> number of primes <= i
 int dp[PHI_N][PHI_K]; // precal of yo(n,k)
 bitset <MAXN> f;
 void sieve(int n) {
  f[1] = true;
  for (int i = 4; i <= n; i += 2) f[i] = true;
  for (int i = 3; i * i <= n; i += 2) {
   if (!f[i]) { for (int j = i * i; j <= n; j += i << 1) f[j] = 1;}}
  for (int i = 1; i <= n; i++) {
   if (!f[i]) primes[len++] = i; pref[i] = len;}}
 void init() {
  sieve(MAXN - 1);
  // precalculation of phi upto size (PHI_N,PHI_K)
  for (int n = 0; n < PHI_N; n++) dp[n][0] = n;
  for (int k = 1; k < PHI_K; k++) {
   for (int n = 0; n < PHI_N; n++) {
    dp[n][k] = dp[n][k - 1] - dp[n / primes[k - 1]][k - 1];
   }}}
 // returns the number of integers less or equal n which are
 // not divisible by any of the first k primes
 // recurrence --> yo(n, k) = yo(n, k-1) - yo(n / p_k , k-1)
 // for sum of primes yo(n, k) = yo(n, k-1) - p_k * yo(n / p_k , k-1)
 long long yo(long long n, int k) {
  if (n < PHI_N && k < PHI_K) return dp[n][k];
  if (k == 1) return ((++n) >> 1);
  if (primes[k - 1] >= n) return 1;
  return yo(n, k - 1) - yo(n / primes[k - 1], k - 1);}
 // complexity: n^(2/3).(log n^(1/3))
 long long Legendre(long long n) {
  if (n < MAXN) return pref[n];
  int lim = sqrt(n) + 1;
  int k = upper_bound(primes, primes + len, lim) - primes;
  return yo(n, k) + (k - 1);}
 // runs under 0.2s for n = 1e12
 long long Lehmer(long long n) {
  if (n < MAXN) return pref[n];
  long long w, res = 0;
  int b = sqrt(n), c = Lehmer(cbrt(n)), a = Lehmer(sqrt(b)); b = Lehmer(b);
  res = yo(n, a) + ((1LL * (b + a - 2) * (b - a + 1)) >> 1);
  for (int i = a; i < b; i++) {
   w = n / primes[i];
   int lim = Lehmer(sqrt(w)); res -= Lehmer(w);
   if (i <= c) {
    for (int j = i; j < lim; j++) {
     res += j;
     res -= Lehmer(w / primes[j]);
    }}}
  return res;}}
\end{minted}
\subsection{Some Primes}
\begin{minted}{c++}
1e9+9, 1e9+21, 1e9+33, 1e9+87, 1e9+93, 1e9+97
1e18+21, 1e18+33, 1e18+87, 1e18+93, 2e18+11, 2e18+33, 2e18+63, 2e18+87
\end{minted}
\subsection{mobius}
\begin{minted}{c++}
// returns mu[x]  O(logn)
ll mu(ll x){ll cnt=0;while(x>1){ll cur=0,d=spf[x];
    while(x%d==0){x/=d;cur++;if(cur>1) return 0;}
    cnt++;}if(cnt&1) return -1;else return 1;}
// from 1 to n
void mobius(){mu[1] = 1;for(ll i=2;i<N;i++){
    mu[i]--;for(ll j=2*i;j<N;j+=i)mu[j]-=mu[i];}}
\end{minted}
\subsection{nCr MOD p}
\begin{minted}{c++}
//nCr_mod_p_by_lucas_for_small_p._O(LogN)
ll Lucas(ll n, ll m){
  if(n==0 && m==0) return 1ll;
  ll ni = n % MOD; ll mi = m % MOD;
  if(mi>ni) return 0;
  return (Lucas(n/MOD, m/MOD)*nCr(ni, mi))%MOD; }
ll nCr_by_lucas(ll n, ll r){ return Lucas(n, r); }
//nCr_mod_M, M_is_not_prime
const int N = 142858; //need primes <=M (nCr%M)
// returns n! % mod, mod = multiple of p
// O(mod * log(n))
int factmod(ll n, int p, const int mod){
  vector<int> f(mod + 1); f[0] = 1 % mod;
  for(int i = 1; i <= mod; i++){
    if (i % p) f[i] = 1LL * f[i - 1] * i % mod;
    else f[i] = f[i - 1]; }
  int ans = 1 % mod;
  while (n > 1) {
    ans = 1LL * ans * f[n % mod] % mod;
    ans = 1LL*ans*BigMod(f[mod], n/mod, mod)%mod;
    n /= p;
  } return ans; }
ll multiplicity(ll n, int p){
  ll ans = 0;
  while (n) { n /= p; ans += n;}  return ans;}
// C(n, r) modulo p^k
// O(p^k log n)
int ncr(ll n, ll r, int p, int k){
  if(n < r || r < 0) return 0;
  int mod = 1;
  for (int i = 0; i < k; i++) { mod *= p; }
  ll t = multiplicity(n, p) - multiplicity(r, p) - multiplicity(n - r, p);
  if(t >= k) return 0;
  int ans = 1LL * factmod(n, p, mod) * inverse(factmod(r, p, mod), mod) % mod * inverse(factmod(n - r, p, mod), mod) % mod;
  ans = 1LL * ans * BigMod(p, t, mod) % mod;
  return ans;}
pair<ll, ll> CRT(ll a1, ll m1, ll a2, ll m2) {
  ll p, q;
  ll g = e_gcd(m1, m2, p, q);
  if(a1 % g != a2 % g) return make_pair(0, -1);
  ll m = m1 / g * m2;
  p = (p % m + m) % m;  q = (q % m + m) % m;
  return make_pair((p * a2 % m * (m1 / g) % m + q * a1 % m * (m2 / g) % m) %  m, m); }
// O(m log(n) log(m))
int ncr(ll n, ll r, int m) {
  if(n < r || r < 0) return 0;
  pair<ll, ll> ans({0, 1});
  while (m > 1) {
    int p = spf[m], k = 0, cur = 1;
    while (m % p == 0) {
      m /= p; cur *= p; ++k; }
    ans=CRT(ans.first, ans.second, ncr(n,r,p,k),cur);
  } return ans.first; }
\end{minted}
\subsection{phi and sieve(linear)}
\begin{minted}{c++}
int lp[MAX+1],phi[MAX+1];
vector<int> prm;
void pre() {
  phi[1]=1;
  for(int i=2; i<=MAX; ++i){
    if(lp[i] == 0){
      lp[i]=i,phi[i]=i-1; prm.push_back(i);} else {
      if(lp[i]==lp[i/lp[i]]) phi[i]=phi[i/lp[i]]*lp[i];
      else phi[i]=phi[i/lp[i]]* (lp[i]-1);  }
    for(int j=0; j<(int)prm.size()&&prm[j]<=lp[i]&& i*prm[j]<=MAX; ++j) 
      lp[i*prm[j]] = prm[j]; } }
ll Euler_Phi(ll n) {
  ll i,result=n;
  for(i=0;i<prime.size()&&prime[i]*prime[i]<=n;i++){
    if(n%prime[i]==0){
      while(n%prime[i]==0) n/=prime[i];
      result=(result/prime[i])*(prime[i]-1); } }
  if(n>1) result=(result/n)*(n-1); return result; }
\end{minted}
\section{String}
\subsection{Aho Corasick}
\begin{minted}{c++}
#define t_sz 26
struct node{
  int ending,link,par,next[t_sz]; vector<int> idx;
  node(){
    ending=0; par=-1; link=-1; clean(next,-1); } };
struct aho_corasick{
  vector< node > aho;
  aho_corasick(){ aho.emplace_back(); }
  int ID(char ch){  if('a'<=ch && ch<='z') return ch-'a';
    return ch-'A'; }
  void ADD(string &s,int id=0){
    int now,u=0;
    for(auto ch:s){
      now=ID(ch);
      if(aho[u].next[now]==-1) {
        aho[u].next[now]=aho.size(); aho.emplace_back();}
      u=aho[u].next[now]; }
    aho[u].ending++; aho[u].idx.push_back(id); }
  int transition(int u, int i) {
    if(u==-1) return 0;
    if(~aho[u].next[i]) return aho[u].next[i];
    return aho[u].next[i]=transition(aho[u].link, i); }
  void push_links(){
    queue<int> q; q.push(0);
    while(!q.empty()){
      int u=q.front();q.pop();
      for(int i=0; i<t_sz; i++){
        int v=aho[u].next[i]; if(v==-1) continue;
        aho[v].par=u; aho[v].link=transition(aho[u].link,i);
        aho[v].ending+=aho[aho[v].link].ending;
        for(auto &id : aho[aho[v].link].idx)
            aho[v].idx.push_back(id);
        q.push(v); } } }
  int process(string &s, vector<string> &v){
    //v: list of pattern,int n=v.size();vector<int>pos[n];
    int u=0,sum=0;
    for(int i=0; i<s.size(); i++){
      int x=ID(s[i]); u=transition(u, x);
      sum+=aho[u].ending;
      //for(auto id : aho[u].idx) pos[id].push_back(i);} return sum; }
  void clear(){aho.clear();aho.emplace_back();  } };
main() {
  aho_corasick ac;  for (int i = 0; i < n; i++) ac.ADD(v[i], i);
  ac.push_links(); ac.process(text); ac.clear();}
\end{minted}
\subsection{KMP}
\begin{minted}{c++}
vector<int> prefix_function(string s){
  int n = (int)s.length(); vector<int> pi(n);
  for(int i = 1; i < n; i++) {
    int j = pi[i-1];
    while(j>0 && s[i] != s[j]) j = pi[j-1];
    if(s[i] == s[j]) j++;
    pi[i] = j; }  return pi;  }
\end{minted}
\subsection{Manachar}
\begin{minted}{c++}
int d[2][MAX];
void manachar(string &s) {
  int n=(int)s.size();
  for(int t=0; t<2; t++) {
    for(int i=0,l=0,r=-1; i<n; i++){
      int k=(i>r)? 0:min(d[t][l+r-i+(t^1)], r-i+1);
      while(0<=i-k-(t^1) && i+k<n && s[i-k-(t^1)]==s[i+k]) 
        k++;
      d[t][i]=k--;
      if(i+k>r) {  l=i-k-(t^1); r=i+k;}}}}
\end{minted}
\subsection{Palindromic Tree}
\begin{minted}{c++}
const int N = 10004; const int t_sz= 26;
struct Palindromic_Tree{
  int tree[N][t_sz], idx,len[N], link[N], t,n;
  int endHere[N],occ[N],total=0; string s="#";
  int stop[N];//for printing
  Palindromic_Tree(){
    memset(occ, 0, sizeof(occ));
    len[1] = -1, link[1] = 1; len[2] = 0, link[2] = 1;
    endHere[1]=endHere[2]=0; idx = t = 2; }
  void add(int p) {
    while(s[p-len[t]-1] != s[p]) t=link[t];
    int x=link[t], c=s[p]-'a';
    while(s[p-len[x]-1] != s[p]) x=link[x];        
    if(!tree[t][c]) {
      tree[t][c] = ++idx; len[idx] = len[t] + 2;
      link[idx] = len[idx] == 1 ? 2 : tree[x][c];
      endHere[idx]=1+endHere[link[idx]];  }
    t = tree[t][c]; occ[t]++; stop[t]=p; }
  void init(string &ss){
    s+=ss; n=(int)s.size();
    for(int i=1; i<n; i++){ add(i); total+=endHere[t];
    /*palindrom end in pos i=endHere[t]*/  }
    for(int i=idx; i>2; i--){
      occ[link[i]]+=occ[i];
      /*print palindrom 
      int r=stop[i],l=stop[i]-len[i]+1;
      cout<<s.substr(l,len[i])<<endl;
      cout<<"len="<<len[i]<<",cnt="<<occ[i]<<"\n";*/  }  }
  void clear(){
    for(int i=0; i<=idx; i++){
      occ[i]=endHere[i]=len[i]=link[i]=0;
      for(int j=0; j<t_sz; j++) tree[i][j]=0;  }
    len[1] = -1, link[1] = 1;  len[2] = 0, link[2] = 1;
    endHere[1]=endHere[2]=0; idx = t = 2; total=0; s="#";}
  int cntDistinct(){ return idx-2; }
  int cntTotal(){ return total; }  };
//counting
vi suf;//suf[i]= # pali end at i
vi pref;//pref[i]= #pali start from i
p.init(s, suf); p.clear();
reverse(s);p.init(s,pref);reverse(pref);
void init(string &ss, vi &v) {
  //...//palindrom end in pos i=endHere[t]
  v.push_back(endHere[t]); /*....*/ }
//common palindrom?
void add(int p, int sId) {
  /*...*/ occ[t][sId]++; }
ll init(string &s1, string &s2) {
  ll commonPal=0LL; s+=s1; n=s.size();
  add(i, 0) for i = [1,n-1]
  s+=char('a'+26); s+=s2; n2=s.size();
  add(i, 1); for i = [n, n2-1]
  for(int i = idx; i > 2; i--){
    occ[link[i]][0] += occ[i][0];
    occ[link[i]][1] += occ[i][1];
    commonPal+=1LL*occ[i][0]*occ[i][1]; }
  return commonPal; }
\end{minted}
\subsection{Suffix Array}
\begin{minted}{c++}
int sa[N], rank[2 * N], oldrank[2 * N], lcp[N];
int t;
bool compare(int i, int j) { return oldrank[i + t] < oldrank[j + t];}
void build(const string &s) {
 int n = (int)s.size();
 if (n == 0) return;
 int bc[256] = {0};
 for (int i = 0; i < n; ++i) ++bc[s[i]];
 for (int i = 1; i < 256; ++i) bc[i] += bc[i - 1];
 for (int i = 0; i < n; ++i) sa[--bc[s[i]]] = i;
 for (int i = 0; i < n; ++i) rank[i] = bc[s[i]];
 for (t = 1; t < n; t <<= 1) {
  for (int i = 0; i < n; ++i) oldrank[i] = rank[i];
  for (int i = n; i < 2 * n; ++i) oldrank[i] = -1;
  for (int i = 0, j = 1; j < n; i = j++) {
   while (j < n && oldrank[sa[j]] == oldrank[sa[i]]) j++;
   if (j - i == 1) continue;
   int *start = sa + i, *end = sa + j;
   sort(start, end, compare); int num = i; rank[*start] = num;
   for (int *p = start + 1; p < end; ++p) {
    if (oldrank[*p + t] != oldrank[*(p - 1) + t]) num = (int)(p - sa);
    rank[*p] = num;}}}
 int size = 0; lcp[0] = 0;
 for (int i = 0; i < n; i++) {
  if (rank[i] > 0) { int j = sa[rank[i] - 1];
   while (i + size < n && j + size < n && s[i + size] == s[j + size]) ++size;
   lcp[rank[i]] = size; if (size > 0) --size;}}}
\end{minted}
\subsection{Suffix Automata}
\begin{minted}{c++}
struct SuffixAutomata {
 struct node { int len, link, firstpos;
 map<char, int> nxt; };
 int sz, last; vector<node> t;
 SuffixAutomata() {}
 SuffixAutomata(string &s) {
 int n = s.size(); sz = 1; last = 0;
 t.resize(2 * n); t[0].len = 0;
 t[0].link = -1; t[0].firstpos = 0;
 for (auto x: s) extend(x); }
 void extend(char c) {
 int p = last; if (t[p].nxt.count(c)){
  int q = t[p].nxt[c];
  if(t[q].len == t[p].len + 1) {
   last = q; return; }
  int clone = sz++; t[clone] = t[q];
  t[clone].len = t[p].len + 1;
  t[q].link = clone; last = clone;
  while (p != -1 && t[p].nxt[c] == q) {
   t[p].nxt[c] = clone; p = t[p].link; }
  return; }
 int cur = sz++; t[cur].len = t[last].len + 1;
 t[cur].firstpos = t[cur].len; p = last;
 while (p != -1 && !t[p].nxt.count(c)) {
  t[p].nxt[c] = cur; p = t[p].link; }
 if(p == -1) t[cur].link = 0;
 else { int q = t[p].nxt[c];
  if(t[p].len + 1 == t[q].len) t[cur].link = q;
  else { int clone = sz++; t[clone] = t[q];
  t[clone].len = t[p].len + 1;
  while (p != -1 && t[p].nxt[c] == q) {
   t[p].nxt[c] = clone; p = t[p].link; }
  t[q].link=t[cur].link=clone; }} last=cur;}
 ll unique_substrings() { /*O(N)*/ ll tot = 0;
 for(int i = 1; i < sz; i++)
  tot += t[i].len - t[t[i].link].len;
 return tot; }
 ll tot_len_diff_substings() { ll tot = 0;
 for(int i = 1; i < sz; i++) {
  ll shortest = t[t[i].link].len + 1;
  ll longest = t[i].len;
  ll num_strings = longest - shortest + 1;
  ll cur = num_strings * (longest + shortest)/2;
  tot += cur; } return tot; }
 int first_occurrence(string &s) { int u = 0;
 for(auto ch : s) { if(!t[u].nxt.count(ch))
  return -1; u = t[u].nxt[ch]; }
 return t[u].firstpos - s.size() + 1; } };
//lexicographically_kth_distinct
ll dp[MAX]; SuffixAutomata sa;
ll dfs_path_count(int u = 0) {
 if(dp[u]) return dp[u]; dp[u] = 1;
 for(auto [ch, v] : sa.t[u].nxt)
  dp[u] += dfs_path_count(v);
 return dp[u]; }
void kth_lexi(ll k, int u = 0) {
 if(!k) {cout << "\n"; return;} k--;
 for(auto [ch, v] : sa.t[u].nxt) {
 if(dp[v] <= k) k -= dp[v];
 else {cout<<ch; return kth_lexi(k, v); }}}
main() { sa = SuffixAutomata(s);
 dfs_path_count(); kth_lexi(k);}
\end{minted}
\subsection{Z}
\begin{minted}{c++}
vector<int> z_function(string s){
  int n = (int) s.length(); vector<int> z(n); /*z[0]=n;*/
  for(int i=1, l=0, r=0; i<n; ++i) {
    if(i <= r)  z[i] = min(r-i+1, z[i-l]);
    while(i+z[i]<n && s[z[i]]==s[i+z[i]]) ++z[i];
    if(i + z[i] - 1 > r) l=i, r=i+z[i]-1;}  return z;   }
\end{minted}



\section{Formulae}
\footnotesize
\subsection{Combinatorics}

\begin{enumerate}
    \item $\displaystyle \displaystyle \sum \limits_{0\leq k \leq n} {n-k \choose k} = Fib_{n+1}$
    \item $\displaystyle \displaystyle {n \choose k}={n \choose n-k}$
    \item $\displaystyle \displaystyle {n \choose k}+{n \choose k+1}={n+1 \choose k+1}$
    \item $\displaystyle \displaystyle k{n \choose k}=n{n-1 \choose k-1}$
    \item $\displaystyle \displaystyle {n \choose k}=\frac{n}{k}{n-1 \choose k-1}$
    \item $\displaystyle \sum \limits_{i=0}^n{n \choose i}=2^n$
    \item $\displaystyle \sum \limits_{i\geq 0}{n \choose 2i}=2^{n-1}$
    \item $\displaystyle \sum \limits_{i\geq 0}{n \choose 2i+1}=2^{n-1}$
    \item $\displaystyle \sum \limits_{i= 0}^k \left( -1 \right) ^i{n \choose i}=\left( -1 \right) ^k{n-1 \choose k}$
    \item $\displaystyle \sum \limits_{i= 0}^k{n+i \choose i}= \sum \limits_{i= 0}^k{n+i \choose n} = {n+k+1 \choose k}$
    \item $\displaystyle \displaystyle1{n \choose 1}+2{n \choose 2}+3{n \choose 3}+\dots+n{n \choose n}=n2^{n-1}$
    \item $\displaystyle \displaystyle1^2{n \choose 1}+2^2{n \choose 2}+3^2{n \choose 3}+\dots+n^2{n \choose n}=(n+n^2)2^{n-2}$
    \item \textbf{Vandermonde’s Identify:} \\ $\displaystyle \sum \limits_{k=0}^r{m \choose k}{n \choose r-k}={m+n \choose r}$
    \item \textbf{Hockey-Stick Identify:} \\ $\displaystyle \ n,r \in N, n > r, \sum \limits_{i=r}^n{i \choose r}={n+1 \choose r+1}$
    \item $\displaystyle \sum \limits_{i=0}^k{k \choose i}^2={2k \choose k}$
    \item $\displaystyle \sum \limits_{k=0}^n{n \choose k}{n \choose n-k}={2n \choose n}$
    \item $\displaystyle \sum \limits_{k=q}^n{n \choose k}{k \choose q}=2^{n-q}{n \choose q}$
    \item $\displaystyle \sum \limits_{i=0}^nk^i{n \choose i}=(k+1)^n$
    \item $\displaystyle \sum \limits_{i=0}^n{2n \choose i}=2^{2n-1}+\frac{1}{2}{2n \choose n}$
    \item $\displaystyle \sum \limits_{i=1}^n{n \choose i}{n-1 \choose i-1}={2n-1 \choose n-1}$
    \item $\displaystyle \sum \limits_{i=0}^n{2n \choose i}^2=\frac{1}{2} \left( {4n \choose 2n}+{2n \choose n}^2 \right)$
    \item \textbf{Highest Power of $\displaystyle 2$ that divides $^{2n}C_n$:}
    Let $\displaystyle x$ be the number of $\displaystyle 1$s in the binary representation. Then the number of odd terms will be $\displaystyle 2^x$. Let it form a sequence. The $n$-th value in the sequence (starting from $n$ = 0) gives the highest power of $\displaystyle 2$ that divides $^{2n}C_n$.
    \item \textbf{Pascal Triangle}
    \begin{enumerate}
        \item In a row $\displaystyle p$ where $\displaystyle p$ is a prime number, all the terms in that row except the $\displaystyle 1$s are multiples of $\displaystyle p$.
        \item Parity: To count odd terms in row $n$, convert $n$ to binary. Let $\displaystyle x$ be the number of $\displaystyle 1$s in the binary representation. Then the number of odd terms will be $\displaystyle 2^x$.
        \item Every entry in row $\displaystyle 2^n-1, n\geq 0,$ is odd.
    \end{enumerate}
    \item An integer $n\geq 2$ is prime if and only if all the intermediate binomial coefficients \\ \\  $\displaystyle \displaystyle {n \choose 1},{n \choose 2},\dots, {n \choose n-1}$ are divisible by $n$.
    \item \textbf{Kummer’s Theorem:} For given integers $n\geq m\geq 0$ and a prime number $\displaystyle p$, the largest power of $\displaystyle p$ dividing $\displaystyle \displaystyle {n \choose m}$ is equal to the number of carries when $m$ is added to $n$-$m$ in base $\displaystyle p$. For implementation take inspiration from lucas theorem.
    \item Number of different binary sequences of length $n$ such that no two $0$’s are adjacent=$\displaystyle Fib_{n+1}$
    \item \textbf{Combination with repetition:} Let’s say we choose $\displaystyle \displaystyle k$ elements from an $n$-element set, the order doesn’t matter and each element can be chosen more than once. In that case, the number of different combinations is: $\displaystyle \displaystyle {n+k-1 \choose k}$
    \item Number of ways to divide $n$ persons in $\displaystyle \frac{n}{k}$ equal groups i.e. each having size $\displaystyle \displaystyle k$ is
    \[\displaystyle \frac{n!}{k!^{\frac{n}{k}} \left( \frac{n}{k} \right) !}= \prod \limits_{n \geq k}^{n-=k}{n-1 \choose k-1}\]
    \item The number non-negative solution of the equation: $\displaystyle x_1+x_2+x_3+\dots+x_k=n \text{ is}{n+k-1 \choose n}$
    \item Number of ways to choose $n$ ids from $\displaystyle 1$ to b such that every id has distance at least k $\displaystyle \displaystyle =\left( \frac{b-(n-1)(k-1)}{n} \right)$
    \item $\displaystyle \displaystyle \sum \limits_{i=1,3,5,\dots}^{i\leq n} \binom{n}{i} a^{n-i}b^i = \frac{1}{2} ((a+b)^n-(a-b)^n)$
    \item $\displaystyle \displaystyle \sum \limits_{i=0}^{n} \dfrac{\binom{k}{i}}{\binom{n}{i}} = \dfrac{\binom{n+1}{n-k+1}}{\binom{n}{k}}$
    \item \textbf{Derangement:} a permutation of the elements of a set, such that no element appears in its original position. Let $d(n)$ be the number of derangements of the identity permutation fo size $n$.
    \[\begin{array}{@{}l}d(n) = (n - 1)\bigl(d(n - 1) + d(n - 2)\bigr),\\ \text{where } d(0)=1,\ d(1)=0.\end{array}\]
    \item \textbf{Involutions:} permutations such that $\displaystyle p^2 =$ identity permutation. $\displaystyle a_0 = a_1 = 1$ and $\displaystyle a_n=a_{n-1} + (n-1)a_{n-2}$ for $n>1$.
    \item Let $T(n,k)$ be the number of permutations of size $n$ for which all cycles have length $\displaystyle \leq k$.
    \[T(n, k)=
    \begin{cases}
      n! & ; n \le k\\
      n \cdot T(n-1, k) \\ - F(n-1, k) \\ \cdot T(n-k-1, k) & ;n > k \\
    \end{cases}\]
    Here $\displaystyle F(n, k) = n \cdot (n - 1) \cdot \dots \cdot (n - k + 1)$
    \item \textbf{Lucas Theorem}
    \begin{enumerate}
      \item If $\displaystyle p$ is prime, then $\displaystyle \left(\frac{p^{a}}{k}\right) \equiv 0 (\mod\; p)$
      \item For non-negative integers $m$ and $n$ and a prime $\displaystyle p$, the following congruence relation holds:
      \(\displaystyle \left(\frac{m}{n}\right) \equiv \prod_{i=0}^{k} \left(\frac{m_{i}}{n_{i}}\right) (\text{mod}\; p),\)
      where,
      $m=m_{k} p^{k} +m_{k-1} p^{k-1}+\dots+m_{1} p+m_{0}$,
      and
      $n=n_{k}p^{k}+n_{k-1}p^{k-1}+\dots+n_{1}p+n_{0}$
      are the base $\displaystyle p$ expansions of $m$ and $n$ respectively. This uses the convention that $\displaystyle \left(\frac{m}{n}\right)=0$,when $m<n$.
    \end{enumerate}
    \item $\displaystyle \sum_{i=0}^{n}  {n \choose i } \cdot i^{k}$
    \\= $\displaystyle \sum_{i=0}^{n}  {n \choose i } \cdot \sum_{j=0}^{k}{ \left\{ \begin{matrix} k\cr j \end{matrix} \right\} \cdot i^{\underline{j}}}$
    \\= $\displaystyle \sum_{i=0}^{n} {n \choose i } \cdot \sum_{j=0}^{k} \left\{ \begin{matrix} k\cr j \end{matrix} \right\} \cdot  j! {n \choose i }$
    \\= $\displaystyle \sum_{i=0}^{n} \frac{n!}{(n-i)!} \cdot \sum_{j=0}^{k} \left\{ \begin{matrix}k\cr j \end{matrix} \right\} \cdot \frac{1}{(i-j)!}$
    \\= $\displaystyle \sum_{i=0}^{n} \sum_{j=0}^{k} \frac{n!}{(n-i)!} \cdot  \left\{\begin{matrix} k\cr j \end{matrix} \right\} \cdot \frac{1}{(i-j)!}$
    \\= n!$\displaystyle \sum_{i=0}^{n} \sum_{j=0}^{k} \left\{ \begin{matrix} k\cr j \end{matrix} \right\} \cdot \frac{1}{(n-i)!} \cdot \frac{1}{(i-j)!}$
    \\= n!$\displaystyle \sum_{i=0}^{n} \sum_{j=0}^{k} \left\{ \begin{matrix} k\cr j \end{matrix} \right\} \cdot { n-j \choose n-i} \cdot \frac{1}{(n-j)!}$
    \\= n!$\displaystyle \sum_{j=0}^{k} \left\{ \begin{matrix} k\cr j \end{matrix} \right\} \cdot \frac{1}{(n-j)!} \sum_{i=0}^{n} \cdot { n-j \choose n-i}$
    \\= $\displaystyle \sum_{j=0}^{k} \left\{\begin{matrix} k\cr j \end{matrix} \right\} \cdot n^{\underline{j}} \cdot 2^{n-j}$

    Here $n^{\underline{j}}=P(n,j)=\dfrac{n!}{(n-j)!}$ and $\displaystyle \left\{ \begin{matrix} k\cr j \end{matrix} \right\}$ is stirling number of the second kind.

    So, instead of $O(n)$, now you can calculate the original equation in $O(k^2)$ or even in $O(k \log^2 n)$ using NTT.
    \item $\displaystyle \displaystyle \sum \limits_{i=0}^{n-1} {i \choose j} x^i = x^j(1-x)^{-j-1} \\ \cdot \left( 1- x^n \sum \limits_{i=0}^j {n \choose i} x^{j-i} (1-x)^i \right)$
    \item $\displaystyle x_{0}, x_{1}, x_{2}, x_{3}, \dots , x_{n}$
    \\ $\displaystyle x_{0} + x_{1}, x_{1} + x_{2}, x_{2}+x_{3}, \dots x_{n}$
    \\ \dots
    \\ If we continuously do this $n$ times then the polynomial of the first column of the $n$-th row will be
    \[\displaystyle p(n)=\sum_{k=0}^{n}{n \choose k} \cdot x(k)\]
    \item If $\displaystyle P(n)=\sum_{k=0}^{n}{n \choose k} \cdot Q(k)$, then,
    \[Q(n)=\sum_{k=0}^{n}(-1)^{n-k}{n \choose k} \cdot P(k)\]
    \item If $\displaystyle P(n)=\sum_{k=0}^{n}(-1)^{k}{n \choose k} \cdot Q(k)$ , then,
    \[Q(n)=\sum_{k=0}^{n}(-1)^{k}{n \choose k} \cdot P(k)\]
\end{enumerate}

\subsection{Catalan Numbers}

\begin{enumerate}[resume]
    \item $\displaystyle C_n=\frac{1}{n+1}{2n \choose n}$
    \item $\displaystyle C_0=1,C_1=1\text{ and }C_n=\sum \limits_{k=0}^{n-1}C_k C_{n-1-k}$
    \item Number of correct bracket sequence consisting of $n$ opening and $n$ closing brackets.
    \item The number of ways to completely parenthesize $n$+$\displaystyle 1$ factors.
    \item The number of triangulations of a convex polygon with $n$+$\displaystyle 2$ sides (i.e. the number of partitions of polygon into disjoint triangles by using the diagonals).
    \item The number of ways to connect the $\displaystyle 2n$ points on a circle to form $n$ disjoint i.e. non-intersecting chords.
    \item The number of monotonic lattice paths from point $\displaystyle (0,0)$ to point $\displaystyle (n,n)$ in a square lattice of size $n\times n$, which do not pass above the main diagonal (i.e. connecting $\displaystyle (0,0)$ to $\displaystyle (n,n)$).
    \item The number of rooted full binary trees with $n$+$\displaystyle 1$ leaves (vertices are not numbered). A rooted binary tree is full if every vertex has either two children or no children.
    \item Number of permutations of $\displaystyle {1, \dots, n}$ that avoid the pattern $\displaystyle 123$ (or any of the other patterns of length $3$); that is, the number of permutations with no three-term increasing sub-sequence. For $n = 3$, these permutations are $\displaystyle 132,\ 213,\ 231,\ 312$ and $321.$\displaystyle For $n = 4$, they are $\displaystyle 1432,\ 2143,\ 2413,\ 2431,\ 3142,\ 3214,\\ \ 3241,\ 3412,\ 3421,\ 4132,\ 4213,\ 4231,\ 4312$ and $4321.$
    \item \textbf{Balanced Parentheses count with prefix:}
    The count of balanced parentheses sequences consisting of $n+k$ pairs of parentheses where the first $\displaystyle \displaystyle k$ symbols are open brackets. Let the number be $\displaystyle C_n^{(k)}$, then
    \[\displaystyle \displaystyle C_n^{(k)} = \frac{k+1}{n+k+1} \binom{2n+k}{n}\]
\end{enumerate}

\subsection{Narayana numbers}

\begin{enumerate}[resume]
    \item $N(n,k)=\frac{1}{n}\left(\frac{n}{k}\right)\left(\frac{n}{k-1}\right)$
    \item The number of expressions containing $n$ pairs of parentheses, which are correctly matched and which contain $\displaystyle \displaystyle k$ distinct nestings. For instance, $N(4, 2) = 6$ as with four pairs of parentheses six sequences can be created which each contain two times the sub-pattern ‘()’.
\end{enumerate}

\subsection{Stirling numbers of the first kind}

\begin{enumerate}[resume]
    \item The Stirling numbers of the first kind count permutations according to their number of cycles (counting fixed points as cycles of length one).
    \item $S(n,k)$ counts the number of permutations of $n$ elements with $\displaystyle \displaystyle k$ disjoint cycles.
    \item $S(n,k)=(n-1) \cdot S(n-1,k)+S(n-1,k-1),$
    \(where,\; S(0,0)=1,S(n,0)=S(0,n)=0\)
    \item $\displaystyle \displaystyle\sum_{k=0}^{n}S(n,k) = n!$
    \item The unsigned Stirling numbers may also be defined algebraically, as the coefficient of the rising factorial:
    \[\displaystyle x^{\bar{n}} = x(x+1)\dots(x+n-1) = \sum_{k=0}^{n}{ S(n, k) x^k}\]
    \item Lets $[n, k]$ be the stirling number of the first kind, then
    \[\displaystyle \bigl[\!\begin{smallmatrix} n \\ n\ -\ k \end{smallmatrix}\!\bigr] = \sum_{0 \leq i_1 < i_2 < i_k < n}{i_1i_2\dotsi_k.}\]
\end{enumerate}

\subsection{Stirling numbers of the second kind}

\begin{enumerate}[resume]
    \item Stirling number of the second kind is the number of ways to partition a set of n objects into k non-empty subsets.
    \item $S(n,k)=k \cdot S(n-1,k)+S(n-1,k-1)$,
    \(where \; S(0,0)=1,S(n,0)=S(0,n)=0\)
    \item $S(n,2)=2^{n-1}-1$
    \item $S(n,k) \cdot k!$ = number of ways to color $n$ nodes using colors from $\displaystyle 1$ to $\displaystyle \displaystyle k$ such that each color is used at least once.
    \item An $r$-associated Stirling number of the second kind is the number of ways to partition a set of $n$ objects into $\displaystyle \displaystyle k$ subsets, with each subset containing at least $r$ elements. It is denoted by $S_r( n , k )$ and obeys the recurrence relation.
    $\displaystyle \displaystyle S_r(n+1,k) = k S_r(n,k) + \binom{n}{r-1} S_r(n-r+1,k-1)$
    \item Denote the n objects to partition by the integers $\displaystyle 1, 2, \dots, n$. Define the reduced Stirling numbers of the second kind, denoted $S^d(n, k)$, to be the number of ways to partition the integers $\displaystyle 1, 2, \dots, n$ into k nonempty subsets such that all elements in each subset have pairwise distance at least d.
    That is, for any integers i and j in a given subset, it is required that $|i - j| \geq d$. It has been shown that these numbers satisfy,
    \(S^d(n, k) = S(n - d + 1, k - d + 1), n \geq k \geq d\)
\end{enumerate}

\subsection{Bell number}

\begin{enumerate}[resume]
    \item Counts the number of partitions of a set.
    \item $B_{n+1}=\displaystyle\sum_{k=0}^{n}\left(\frac{n}{k}\right) \cdot B_{k}$
    \item $B_{n}=\displaystyle\sum_{k=0}^{n}S(n,k)$ ,where $S(n,k)$ is stirling number of second kind.
\end{enumerate}

\subsection{Math}

\begin{enumerate}[resume]
    \item $\displaystyle ab \mod\ ac=a(b \mod\ c)$
    \item $\displaystyle \sum_{i = 0}^n{i\cdot i!=(n + 1)! - 1}$.
    \item $\displaystyle a^k - b^k = (a - b) \cdot (a^{k - 1}b^0 + a^{k - 2}b^1 + \dots + a^0b^{k - 1})$
    \item $\displaystyle \min(a + b, c) = a + \min(b, c - a)$
    \item $|a - b| + |b - c| + |c - a| = 2 (\max (a, b, c) - \min (a, b, c))$
    \item $\displaystyle a \cdot b \leq c \rightarrow a \leq \left \lfloor \dfrac{c}{b} \right \rfloor$ is correct
    \item $\displaystyle a \cdot b < c \rightarrow a <  \left \lfloor \dfrac{c}{b} \right \rfloor$ is incorrect
    \item $\displaystyle a \cdot b \geq c \rightarrow a \geq  \left \lfloor \dfrac{c}{b} \right \rfloor$ is correct
    \item $\displaystyle a \cdot b > c \rightarrow a >  \left \lfloor \dfrac{c}{b} \right \rfloor$ is correct
    \item For positive integer $n$, and arbitrary real numbers $m,x$,
    $\displaystyle \left \lfloor \dfrac{\left \lfloor x/m \right \rfloor}{n} \right \rfloor = \left \lfloor \dfrac{x}{mn} \right \rfloor$
    $\displaystyle \left \lceil \dfrac{\left \lceil x/m \right \rceil}{n} \right \rceil = \left \lceil \dfrac{x}{mn} \right \rceil$
    \item Lagrange’s identity:
    \[\displaystyle \displaystyle {\displaystyle {\begin{array}{@{}l}\left(\sum \limits_{k=1}^{n}a_{k}^{2}\right)\left(\sum \limits_{k=1}^{n}b_{k}^{2}\right)-\left(\sum\limits_{k=1}^{n}a_{k}b_{k}\right)^{2} \\ =\sum \limits_{i=1}^{n-1}\sum \limits_{j=i+1}^{n}\left(a_{i}b_{j}-a_{j}b_{i}\right)^{2} \\ ={\frac {1}{2}}\sum \limits_{i=1}^{n}\sum \limits_{j=1,j\neq i}^{n}(a_{i}b_{j}-a_{j}b_{i})^{2}\end{array}}}\]
    \item $\displaystyle \sum_{i = 1}^n{ia^i} = \frac{a(n a^{n + 1} - (n + 1) a^n + 1)}{(a - 1)^2}$
    \item \textbf{Vieta’s formulas:}
    Any general polynomial of degree $n$
    \[\displaystyle p(x) = a_nx^n + a_{n - 1}x^{n - 1} + \dots + a_1x + a_0\]
    (with the coefficients being real or complex numbers and $\displaystyle a_n \neq 0$) is known by the fundamental theorem of algebra to have $n$ (not necessarily distinct) complex roots $r_1, r_2,\dots, r_n$.
    \[\begin{cases}
        \text{$r_1 + r_2 + \dots + r_{n - 1} + r_n = - \frac{a_{n - 1}}{a_n}$}
        \\
\begin{array}{@{}l}
(r_1r_2 + r_1r_3 + \dots + r_1r_n) \\
+ (r_2r_3 + r_2r_4 + \dots + r_2r_n) + \dots \\
+ r_{n-1}r_n
= \dfrac{a_{n-2}}{a_n}
\end{array}
        \\\vdots\\
        \text{$r_1r_2\dots r_n = (-1)^n\frac{a_0}{a_n}.$}
        \end{cases}\]
    Vieta’s formulas can equivalently be written as
    \[\displaystyle \sum_{1 \leq i_1 < i_2 < \dots <i_k \leq n} \Bigg(\prod_{j = 1}^k{r_{i_j}}\Bigg) = (-1)^k\frac{a_{n - k}}{a_n},\]
    \item We are given n numbers $\displaystyle a_1,a_2,\dots,a_n$ and our task is to find a value $\displaystyle x$ that minimizes the sum,
    \[|a_1 - x| + |a_2 - x| + \dots + |a_n - x|\]
    optimal $\displaystyle x=$ median of the array.
    if $n$ is even $\displaystyle x =$ $[$left median,right median$]$ i.e. every number in this range will work.
    For minimizing
    \[\displaystyle (a_1 - x)^2 + (a_2 - x)^2 + \dots + (a_n - x)^2\]
    optimal $\displaystyle x = \dfrac{(a_1 + a_2 + \dots + a_n)}{ n}$
    \item Given an array a of n non-negative integers. The task is to find the sum of the product of elements of all the possible subsets. It is equal to the product of $\displaystyle (a_i + 1)$ for all $\displaystyle a_i$
    \item \textbf{Pentagonal number theorem:}
    In mathematics, the pentagonal number theorem states that
    \[
\begin{array}{@{}l}
\prod_{n=1}^{\infty} (1 - x^n)
= \sum_{k=-\infty}^{\infty} (-1)^k x^{\frac{k(3k-1)}{2}} \\
= 1 + \sum_{k=1}^{\infty} (-1)^k
\left(
x^{\frac{k(3k+1)}{2}} + x^{\frac{k(3k-1)}{2}}
\right).
\end{array}
\]
    In other words,
\[
\begin{array}{@{}l}
(1 - x)(1 - x^2)(1 - x^3)\cdots
\\ = 1 - x - x^2 + x^5 + x^7 - x^{12} \\ - x^{15} + x^{22} + x^{26} - \cdots .
\end{array}
\]
    The exponents $\displaystyle 1, 2, 5, 7, 12,\cdots$ on the right hand side are given by the formula $g_k = \frac{k(3k - 1)}{2}$ for $\displaystyle \displaystyle k = 1, -1, 2, -2, 3, \cdots$ and are called (generalized) pentagonal numbers.
    It is useful to find the partition number in $O(n \sqrt{n})$
\end{enumerate}

\subsection{Fibonacci Number}

\begin{enumerate}[resume]
    \item $\displaystyle F_0=0, F_1=1$ and $\displaystyle F_n=F_{n-1}+F_{n-2}$
    \item $\displaystyle F_n=\sum\limits_{k=0}^{\lfloor\frac{n-1}{2}\rfloor}\binom{n-k-1}{k}$
    \item $\displaystyle F_n=\frac{1}{\sqrt{5}}(\frac{1+\sqrt{5}}{2})^n-\frac{1}{\sqrt{5}}(\frac{1-\sqrt{5}}{2})^n$
    \item $\displaystyle \sum\limits_{i=1}^{n}F_i=F_{n+2}-1$
    \item $\displaystyle \sum\limits_{i=0}^{n-1}F_{2i+1}=F_{2n}$
    \item $\displaystyle \sum\limits_{i=1}^{n}F_{2i}=F_{2n+1}-1$
    \item $\displaystyle \sum\limits_{i=1}^{n}F_{i}^{2}=F_nF_{n+1}$
    \item $\displaystyle F_mF_{n+1}-F_{m-1}F_{n}=(-1)^nF_{m-n}$
    $\displaystyle F_{2n}=F_{n+1}^2-F_{n-1}^2=F_n(F_{n+1}+F_{n-1})$
    \item $\displaystyle F_mF_n+F_{m-1}F_{n-1}=F_{m+n-1}$
    $\displaystyle F_mF_{n+1}+F_{m-1}F_{n}=F_{m+n}$
    \item A number is Fibonacci if and only if one or both of $\displaystyle (5 \cdot n^2 + 4)$ or $\displaystyle (5 \cdot n^2 - 4)$ is a perfect square
    \item Every third number of the sequence is even and more generally, every $\displaystyle \displaystyle k^{th}$ number of the sequence is a multiple of $\displaystyle F_k$
    \item $gcd(F_m, F_n)=F_{gcd(m, n)}$
    \item Any three consecutive Fibonacci numbers are pairwise coprime, which means that, for every n, $gcd(F_n, F_{n+1})=gcd(F_n, F_{n+2}), gcd(F_{n+1}, F_{n+2})=1$
    \item If the members of the Fibonacci sequence are taken $mod$ $n$, the resulting sequence is periodic with period at most $6n$.
\end{enumerate}

\subsection{Pythagorean Triples}

\begin{enumerate}[resume]
    \item A Pythagorean triple consists of three positive integers $\displaystyle a, b,$ and $\displaystyle C$, such that $\displaystyle a^{2}+b^{2}=c^{2}$. Such a triple is commonly written $\displaystyle (a, b, c)$
    \item Euclid’s formula is a fundamental formula for generating Pythagorean triples given an arbitrary pair of integers m and n with $m >n >0$. The formula states that the integers
    \[\displaystyle a = m^{2}-n^{2}, b = 2mn,  c = m^{2}+n^{2}\]
    form a Pythagorean triple. The triple generated by Euclid’s formula is primitive if and only if m and n are coprime and not both odd. When both m and n are odd, then a, b, and c will be even, and the triple will not be primitive; however, dividing a, b, and c by 2 will yield a primitive triple when m and n are coprime and both odd.
    \item The following will generate all Pythagorean triples uniquely:
  \[
\begin{array}{@{}l}
a = k\,(m^{2} - n^{2}), \\
b = k\,(2mn),
c = k\,(m^{2} + n^{2}).
\end{array}
\]
    where m, n, and k are positive integers with $m > n$, and with m and n coprime and not both odd.
    \item \textbf{Theorem:}
    The number of Pythagorean triples
    {a,b,n} with $max{a,b,n} = n$ is given by
    \[\displaystyle \dfrac{1}{2}\left( \displaystyle\prod _{ p^{\alpha} || n}\left( 2\alpha +1\right) -1\right)\]
    where the product is over all prime divisors p of the form $4k+1$.
    The notation $\displaystyle p^{\alpha} || n$ stands for the highest exponent $\displaystyle \alpha$ for which $\displaystyle p^{\alpha}$ divides $n$
    \textbf{Example:}
    For $n = 2\cdot3^{2}\cdot5^{3}\cdot7^{4}\cdot11^{5}\cdot13^{6},$ the number of Pythagorean triples with hypotenuse n is $\displaystyle \dfrac{1}{2}\left( 7.13-1\right) =45$.
    To obtain a formula for the number of Pythagorean triples with hypotenuse less than a specific positive integer N, we may add the numbers corresponding to each $n < N$ given by the Theorem. There is no simple way to compute this as a function of N.
\end{enumerate}

\subsection{Sum of Squares Function}

\begin{enumerate}[resume]
    \item The function is defined as
    \[
r_k(n)
=
\left|
\begin{array}{@{}l}
(a_1,a_2,\dots,a_k) \in \mathbf{Z}^k : \\
n = a_1^2 + a_2^2 + \dots + a_k^2
\end{array}
\right|
\]
    \item The number of ways to write a natural number as sum of two squares is given by $r_2(n)$. It is given explicitly by
    $r_{2}\left(n\right) = 4\left(d_{1}\left(n\right) - d_{3}\left(n\right)\right)$
    where d1(n) is the number of divisors of n which are congruent with 1 modulo 4 and d3(n) is the number of divisors of n which are congruent with 3 modulo 4.
    The prime factorization $n = 2^{g}p_{1}^{f_{1}}p_{2}^{f_{2}}. . .q_{1}^{h_{1}}q_{2}^{h_{2}}. . . ,$ where $\displaystyle p_{i}$ are the prime factors of the form $\displaystyle p_{i} \equiv 1$ (mod 4), and $q_{i}$ are the prime factors of the form $q_{i} \equiv 3$ (mod 4) gives another formula $r_{2}\left(n\right) = 4\left(f_{1}+1\right)\left(f_{2}+1\right)\dots,$ if all exponents $h_{1}$,$h_{2}$,\dots are even. If one or more $h_{i}$ are odd, then $r_{2}\left(n\right) = 0.$
    \item The number of ways to represent n as the sum of four squares is eight times the sum of all its divisors which are not divisible by 4, i.e.
    $r_{4}\left( n \right) = 8 \displaystyle\sum _{d| n;4 \nmid d}d$
    $r_{8}\left(n\right) = 16 \displaystyle\sum _{d|n} \left(-1\right)^{n+d}d^{3}$
\end{enumerate}
\subsection{Number Theory}
\begin{enumerate}[resume]
    \item for $i > j$, $\displaystyle \gcd(i, j)$ $=$ $\displaystyle \gcd(i -j, j)$ $\displaystyle \leq (i -j)$
    \item $\displaystyle \sum_{x = 1}^n \left[ d | x^k \right] = \left \lfloor\dfrac{n}{\prod_{i = 0}{p_i^{\left \lceil \dfrac{e_i}{k}\right \rceil}}}\right \rfloor$,
    where $d = \prod_{i = 0}{p_i^{e_i}}$. Here, $[a | b]$ means if $\displaystyle a$ divides $b$ then it is $\displaystyle 1$, otherwise it is $0$.
    \item The number of lattice points on segment $\displaystyle (x_1,y_1)$ to $\displaystyle (x_2,y_2)$ is $\displaystyle \gcd(abs(x_1-x_2),abs(y_1-y_2)) + 1$
    \item $\displaystyle (n-1)! \mod  n = n -1$ if n is prime, 2 if $n = 4$, $0$ otherwise.
    \item A number has odd number of divisors if it is perfect square
    \item The sum of all divisors of a natural number n is odd if and only if $n=2^r\cdot k^2$ where $r$ is non-negative and $\displaystyle \displaystyle k$ is positive integer.
    \item Let $\displaystyle a$ and $b$ be coprime positive integers, and find integers $\displaystyle a{\prime}$ and $b{\prime}$ such that $\displaystyle aa{\prime} \equiv 1 \mod b$ and $bb{\prime} \equiv 1 \mod a$. Then the number of representations of a positive integers (n) as a non negative linear combination of $\displaystyle a$ and $b$ is
    \[\displaystyle \frac{n}{ab}-\Big\{\frac{b{\prime} n}{a}\Big\}-\Big\{\frac{a{\prime} n}{b}\Big\} + 1\]
    Here, $\displaystyle {x}$ denotes the fractional part of $\displaystyle x$.
\[
\begin{array}{@{}l}
\sum \limits_{i=1}^{a}\sum \limits_{j=1}^{b}\sum \limits_{k=1}^{c} d(i \cdot j \cdot k)
\\ =
\displaystyle \sum \limits_{\gcd (i,j)=\gcd (j,k) = \gcd (k,i) =1}
\displaystyle \left \lfloor \dfrac{a}{i}\right \rfloor \left \lfloor \dfrac{b}{j}\right \rfloor \left \lfloor \dfrac{c}{k} \right \rfloor
\end{array}
\]
    Here, $d(x) =$ number of divisors of $\displaystyle x$.
    \item \textbf{Gauss’s generalization of Wilson’s theorem:},
    Gauss proved that,
    \[\displaystyle \prod \limits_{k=1 \atop \gcd(k,m)=1}^{m}\!\!k\ \]
    \[\displaystyle \equiv {\begin{cases}-1{\pmod {m}}&{\text{if }}m=4,\;p^{\alpha },\;2p^{\alpha }\\\;\;\,1{\pmod {m}}&{\text{otherwise}}\end{cases}}\]
    where $\displaystyle p$ represents an odd prime and $\displaystyle \alpha$ a positive integer. The values of $m$ for which the product is $-1$ are precisely the ones where there is a primitive root modulo $m$.
\end{enumerate}
\subsection{Divisor Function}
\begin{enumerate}[resume]
    \item $\displaystyle \sigma_x(n)=\sum\limits_{d\vert n}^{}d^x$
    \item It is multiplicative i.e if $\displaystyle \gcd(a, b)=1\to \sigma_x(ab)=\sigma_x(a)\sigma_x(b)$.
    \item \[\displaystyle \sigma_x(n)=\prod\limits_{i=1}^{\tau}\frac{p_i^{(a_i+1)x}-1}{p_i^x-1}\]
    \item \textbf{Divisor Summatory Function}
    \begin{enumerate}
        \item Let $\displaystyle \sigma_0(k)$ be the number of divisors of $\displaystyle \displaystyle k$.
        \item $D(x)=\sum\limits_{n\leq{x}}{\sigma_0(n)}$
        \item $D(x)=\sum\limits_{k=1}^{x}\lfloor\frac{x}{k}\rfloor=2\sum\limits_{k=1}^{u}\lfloor\frac{x}{k}\rfloor-u^2$, where $u=\sqrt{x}$
        \item $D(n)=$Number of increasing arithmetic progressions where $n+1$ is the second or later term. (i.e. The last term, starting term can be any positive integer $\displaystyle \le n$. For example, $D(3)=5$ and there are 5 such arithmetic progressions: $\displaystyle (1, 2, 3, 4); (2, 3, 4); (1, 4); (2, 4); (3, 4).$
    \end{enumerate}
    \item Let $\displaystyle \sigma_1(k)$ be the sum of divisors of k. Then, $\displaystyle \sum\limits_{k=1}^{n}\sigma_1(k)=\sum\limits_{k=1}^{n}{k \left\lfloor \frac{n}{k} \right\rfloor}$
    \item $\displaystyle \prod\limits_{d\vert n}^{}d={n^{\frac{\sigma_0}{2}}}$ if $n$ is not a perfect square, and $=\sqrt{n} \cdot n^{\frac{\sigma_0-1}{2}}$ if $n $ is a perfect square.
\end{enumerate}
\subsection{Euler’s Totient function}
\begin{enumerate}[resume]
    \item The function is multiplicative.
    This means that if $\displaystyle \gcd(m, n) = 1$, $\displaystyle \phi(m \cdot n) = \phi(m) \cdot \phi(n)$.
    \item $\displaystyle \phi(n) = n\prod_{p |n}(1 - \frac{1}{p} )$
    \item If p is prime and (k \geq 1), then,
    $\displaystyle \phi(p^k) = p^{k - 1}(p - 1) = p^k(1 - \frac{1}{p})$
    \item $J_k(n)$, the Jordan totient function, is the number of $\displaystyle \displaystyle k$-tuples of positive integers all less than or equal to n that form a coprime $\displaystyle (k + 1)$-tuple together with $n$. It is a generalization of Euler’s totient, $\displaystyle \phi(n) = J_1(n)$.
    \(J_k(n) = n^k\prod_{p |n}(1 - \frac{1}{p^k})\)
    \item $\displaystyle \sum_{d|n}J_k(d) = n^k$
    \item $\displaystyle \sum_{d|n}\phi(d) = n$
    \item $\displaystyle \phi(n) = \sum_{d|n}\mu(d)\cdot\frac{n}{d} = n\sum_{d|n}\frac{\mu(d)}{d}$
    \item $\displaystyle \phi(n) = \sum_{d|n}d\cdot\mu(\frac{n}{d})$
    \item $\displaystyle \displaystyle {a}\vert{b} \to \varphi(a)\vert{\varphi(b)}$
    \item $n\vert{\varphi(a^n - 1)}$ for $\displaystyle a, n > 1$
    \item $\displaystyle \varphi(mn)=\varphi(m)\varphi(n)\cdot\frac{d}{\varphi(d)}$ where $d=gcd(m, n)$
    Note the special cases
    \[\varphi(2m)=
    \begin{cases}
      2\varphi(m) & ; if \, m  \, is \, even\\
      \varphi(m) & ; if \, m \, is \, odd\\
    \end{cases}\]
    \[\displaystyle \varphi(n^m)=n^{m-1}\varphi(n)\]
    \item $\displaystyle \varphi(lcm(m, n)) \cdot \varphi(gcd(m, n))=\varphi(m) \cdot \varphi(n)$
    Compare this to the formula $lcm(m, n)\cdot gcd(m,n)=m\cdot n$
    \item $\displaystyle \varphi(n)$ is even for $n\geq3$. Moreover, if if $n$ has $r$ distinct odd prime factors, $\displaystyle 2^r \vert \varphi(n)$
    \item $\displaystyle \sum\limits_{d\vert{n}}^{} \frac{\mu^2(d)}{\varphi(d)}=\frac{n}{\varphi(n)}$
    \item $\displaystyle \sum\limits_{1\leq k \leq n, \gcd(k, n)=1}k=\frac{1}{2}n\varphi(n)$ for $n > 1$
    \item $\displaystyle \frac{\varphi(n)}{n}=\frac{\varphi(rad(n))}{rad(n)}$ where $rad(n)=$
    $\displaystyle \prod\limits_{p\vert n, p \, prime} p$
    \item $\displaystyle \phi(m) \geq \log_2{m}$
    \item $\displaystyle \phi(\phi(m))\leq\frac{m}{2}$
    \item When $\displaystyle x \geq \log_2m$, then
    \[n^x\mod m=n^{\phi(m) + x\mod \phi(m)}\mod m\]
    \item $\displaystyle \sum\limits_{1\leq k\leq n, \gcd(k, n)=1}\gcd(k-1, n)=\varphi(n)d(n)$ where $d(n)$ is number of divisors. Same equation for $\displaystyle \gcd(a \cdot k-1, n)$ where $\displaystyle a$ and $n$ are coprime.
    \item For every $n$ there is at least one other integer $m\ne n$ such that $\displaystyle \varphi(m)=\varphi(n).$
    \item $\displaystyle \sum\limits_{i=1}^{n} {\varphi(i) \cdot \lfloor\frac{n}{i}\rfloor=\frac{n*(n+1)}{2}}$
    \item $\displaystyle \sum\limits_{i=1, i \% 2 \ne 0 }^{n} \varphi(i) \cdot \lfloor\frac{n}{i}\rfloor=\sum\limits_{k\geq 1}[\frac{n}{2^k}]^2$.
    Note that $[\, ]$ is used here to denote round operator not floor or ceil
    \item \[\displaystyle \sum\limits_{i=1}^{n}\sum\limits_{j=1}^{n}ij[\gcd(i, j)=1]=\sum\limits_{i=1}^{n}\varphi(i)i^2\]
    \item Average of coprimes of $n$ which are less than $n$ is $\displaystyle \dfrac{n}{2}$.
\end{enumerate}
\subsection{Mobius Function and Inversion}
\begin{enumerate}[resume]
    \item It is a multiplicative function.
    \item \[\sum_{d|n}\mu(d) =
    \begin{cases}
      1 & ; n=1\\
      0 & ; n > 0
    \end{cases}\]
    \item $\displaystyle \displaystyle \sum\limits_{n=1}^N {\mu}^2(n) = \displaystyle \sum\limits_{n=1}^{\sqrt{N}} \mu(k) \cdot \left \lfloor \dfrac{N}{k^2} \right\rfloor$
    This is also the number of square-free numbers $\displaystyle \le n$
    \item \textbf{Mobius inversion theorem:} The classic version states that if g and f are arithmetic functions satisfying
    $g(n) = \displaystyle \sum_{d|n}f(d) $ for every integer $n \ge 1$ then $g(n) = \displaystyle \sum_{d|n}\mu(d)g\left(\frac{n}{d}\right)$ for every integer $n \ge 1$
    \item If $\displaystyle F(n) = \displaystyle \prod_{d|n} f(d)$, then $\displaystyle F(n) = \displaystyle \prod_{d|n} F\left(\frac{n}{d}\right)^{\mu(d)}$
    \item $\displaystyle \displaystyle \sum_{d|n}\mu(d)\phi(d) = \displaystyle \prod_{j=1}^K (2 - P_j)$
    where $\displaystyle p_j$ is the primes factorization of $d$
    \item If $\displaystyle F(n)$ is multiplicative, $\displaystyle F \not\equiv 0$, then $\displaystyle \displaystyle \sum_{d|n} \mu(d) f(d) = \displaystyle \prod_{i=1} (1 - f(P_i)) \cdot$
    where $\displaystyle p_i$ are primes of $n$.
\end{enumerate}
\subsection{GCD and LCM}
\begin{enumerate}[resume]
    \item $\displaystyle \gcd(a, 0) = a$
    \item $\displaystyle \gcd(a, b) = \gcd(b, a \mod b)$
    \item Every common divisor of $\displaystyle a$ and $b$ is a divisor of $\displaystyle \gcd(a,b)$.
    \item if $m$ is any integer, then $\displaystyle \gcd(a + m {\cdot} b, b) = \gcd(a, b)$
    \item The gcd is a multiplicative function in the following sense: if $\displaystyle a_1$ and $\displaystyle a_2$ are relatively prime, then $\displaystyle \gcd(a_1 \cdot a_2, b) = \gcd(a_1, b) \cdot \gcd(a_2,b )$.
    \item $\displaystyle \gcd(a, b){\cdot} \operatorname{lcm}(a, b) = |a{\cdot}b|$
    \item $\displaystyle \gcd(a, \operatorname{lcm}(b, c)) = \operatorname{lcm}(\gcd(a, b), \gcd(a, c))$.
    \item $\displaystyle \operatorname{lcm}(a, \gcd(b, c)) = \gcd(\operatorname{lcm}(a, b), \operatorname{lcm}(a, c))$.
    \item For non-negative integers $\displaystyle a$ and $b$, where $\displaystyle a$ and $b$ are not both zero, $\displaystyle \gcd({n^a} - 1, {n^b} - 1) = n^{\gcd(a,b)} - 1$
    \item $\displaystyle \gcd(a, b) = \displaystyle \sum_{k|a \, \text{and} \, k|b} {\phi(k)}$
    \item $\displaystyle \displaystyle \sum_{i=1}^n [\gcd(i, n) = k] = { \phi{\left(\frac{n}{k}\right)}}$
    \item $\displaystyle \displaystyle \sum_{k=1}^n \gcd(k, n) = \displaystyle \sum_{d|n} d \cdot {\phi{\left(\frac{n}{d}\right)}}$
    \item $\displaystyle \displaystyle \sum_{k=1}^n x^{\gcd(k,n)} = \displaystyle \sum_{d|n} x^d \cdot {\phi{\left(\frac{n}{d}\right)}}$
    \item $\displaystyle \displaystyle \sum_{k=1}^n \frac{1}{\gcd(k, n)} = \displaystyle \sum_{d|n} \frac{1}{d} \cdot {\phi{\left(\frac{n}{d}\right)}} = \frac{1}{n} \displaystyle \sum_{d|n} d \cdot \phi(d)$
    \item $\displaystyle \displaystyle \sum_{k=1}^n \frac{k}{\gcd(k, n)} = \frac{n}{2} \cdot \displaystyle \sum_{d|n} \frac{1}{d} \cdot {\phi{\left(\frac{n}{d}\right)}} = \frac{n}{2} \cdot \frac{1}{n} \cdot \displaystyle \sum_{d|n} d \cdot \phi(d)$
    \item $\displaystyle \displaystyle \sum_{k=1}^n \frac{n}{\gcd(k, n)} = 2 * \displaystyle \sum_{k=1}^n \frac{k}{\gcd(k, n)} - 1$, for $n > 1$
    \item $\displaystyle \displaystyle \sum_{i=1}^n \sum_{j=1}^n [\gcd(i, j) = 1] = \displaystyle \sum_{d=1}^n \mu(d) \lfloor {\frac{n}{d} \rfloor}^2$
    \item $\displaystyle \displaystyle \sum_{i=1}^n \displaystyle\sum_{j=1}^n \gcd(i, j) = \displaystyle \sum_{d=1}^n \phi(d) \lfloor {\frac{n}{d} \rfloor}^2$
    \item $\displaystyle \sum_{i=1}^n \sum_{j=1}^n i \cdot j[\gcd(i, j) = 1] = \sum_{i=1}^n \phi(i)i^2$
    \item $\displaystyle F(n) = \displaystyle \sum_{i=1}^n \displaystyle \sum_{j=1}^n \operatorname{lcm}(i, j) = \displaystyle \sum_{l=1}^n {\left(\frac{\left( 1 + \lfloor{\frac{n}{l} \rfloor} \right) \left( \lfloor{\frac{n}{l} \rfloor} \right)} {2} \right)}^2 \displaystyle \sum_{d|l} \mu(d)ld$
    \item $\displaystyle \gcd(\operatorname{lcm}(a, b), \operatorname{lcm}(b, c), \operatorname{lcm}(a, c)) = \operatorname{lcm}(\gcd(a, b), \gcd(b, c), \gcd(a, c))$
    \item $\displaystyle \gcd(A_L, A_{L+1}, \dots, A_R) = \gcd(A_L, A_{L+1} - A_L, \dots, A_R - A_{R-1})$.
    \item Given n, If $SUM = LCM(1,n) + LCM(2,n) + \dots + LCM(n,n)$
    then SUM = $\displaystyle \dfrac{n}{2}( \displaystyle\sum _{ d| n}\left( \phi \left( d\right) \times d\right) +1$
  \end{enumerate}
  \subsection{Legendre Symbol}
  \begin{enumerate}[resume]
    \item $\displaystyle \left(\frac{a}{p}\right)\equiv a^{\frac{p-1}{2}} \pmod{p} \text{ and } \left(\frac{a}{p}\right)\in\{-1,0,1\}.$
    \item The Legendre symbol is periodic in its first (or top) argument: if $\displaystyle a \equiv b \pmod{p}$, then \\
    $\displaystyle \left(\frac{a}{p}\right)=\left(\frac{b}{p}\right).$
    \item The Legendre symbol is a completely multiplicative function of its top argument: \\
    $\displaystyle \left(\frac{ab}{p} \right) = \left(\frac{a}{p} \right) \left(\frac{b}{p} \right)$
    \item The Fibonacci numbers $\displaystyle 1,1,2,3,5,8,13\dots$ are defined by the recurrence $\displaystyle F_1 = F_2 = 1,F_{n+1} = F_n + F_{n-1}.$ If $\displaystyle p$ is a prime number then \\
    $\displaystyle F_{p-\left(\frac{p}{5}\right)} \equiv 0 {\pmod{p}}, \;\;F_p \equiv \left(\frac{p}{5}\right) {\pmod{p}}.$
    \item Continuing from previous point, \\
    $\displaystyle \left(\frac{p}{5}\right)= \text{infinite concatenation of the sequence}$ \\
    $\left(1,-1,-1,1,0\right) {\text{from }}p \geq 1$.
    \item If $n = k^2$ is perfect square then $\displaystyle \left(\frac{n}{p}\right)=1$ for every odd prime except $\displaystyle \left(\frac{n}{k}\right)=0$ if k is an odd prime.
    \end{enumerate}
  \subsection{Miscellaneous}
  \begin{enumerate}[resume]
    \item $\displaystyle \displaystyle a+b = a \oplus b +2(a \& b)$.
    \item $\displaystyle \displaystyle a + b = a \mid b + a \& b$
    \item $\displaystyle \displaystyle a \oplus b = a\mid b - a \& b$
    \item $\displaystyle \displaystyle k_{th}$ bit is set in $\displaystyle x$ iff $\displaystyle x \mod 2^{k - 1}\geq 2^k$. It comes handy when you need to look at the bits of the numbers which are pair sums or subset sums etc.
    \item $\displaystyle \displaystyle k_{th}$ bit is set in $\displaystyle x$ iff $\displaystyle x \mod 2^{k - 1} - x \mod 2^k \neq\ 0$ ($ = 2^k$ to be exact). It comes handy when you need to look at the bits of the numbers which are pair sums or subset sums etc.
    \item $n \mod 2^i = n \& (2^i - 1)$
    \item $\displaystyle 1\oplus 2 \oplus 3 \oplus \cdots \oplus (4k - 1) = 0$ for any $\displaystyle \displaystyle k \ge 0$
    \item \textbf{Erdos Gallai Theorem:}
    The degree sequence of an undirected graph is the non-increasing sequence of its vertex degrees
    A sequence of non-negative integers $d_1 \geq d_2 \geq \cdots \geq d_n$ can be represented as the degree sequence of finite simple graph on $n$ vertices if and only if $d_1 + d_2 + \cdots + d_n$ is even and
    \[\displaystyle \sum_{i = 1}^k {d_i \leq k(k - 1)} + \sum_{i = k + 1}^n \min(d_i, k)\]
    holds for every $\displaystyle \displaystyle k$ in $\displaystyle 1 \leq k \leq n$.
    \item f_n = \sum_{i=0}^{n} \binom{n}{i} g_i \Rightarrow g_n = \sum_{i=0}^{n} (-1)^{n-i} \binom{n}{i} f_i
\end{enumerate}
\subsection{Traingle}
\begin{itemize}[leftmargin=*]
    \item \textbf{Sine Rule:} $\frac{a}{\sin A} = \frac{b}{\sin B} = \frac{c}{\sin C} = 2R $
    \item \textbf{Cosine Rule:} $c^2 = a^2 + b^2 - 2ab \cos C$
    \item \textbf{Area:} $\Delta = \sqrt{s(s-a)(s-b)(s-c)} \\ = \frac{1}{2}bc \sin A$
    \item \textbf{Circumradius ($R$):} $R = \frac{abc}{4\Delta} = \frac{a}{2\sin A}$
    \item \textbf{Inradius ($r$):} $r = \frac{\Delta}{s}$
    \item \textbf{Ex-radii ($r_a, r_b, r_c$):} $r_a = \frac{\Delta}{s-a}$
    \item \textbf{Centroid ($G$):} $\left( \frac{\sum x}{3}, \frac{\sum y}{3} \right)$
    \item \textbf{Incenter ($I$):} $\left( \frac{ax_1+bx_2+cx_3}{a+b+c}, \frac{ay_1+by_2+cy_3}{a+b+c} \right)$
    \item \textbf{Excenter ($I_A$):} Replace $a$ with $-a$ in Incenter formula.
    \item \textbf{Euler Line:} $O, G, H$ are collinear; $HG = 2GO$.
    \item \textbf{Median ($m_a$):} $4m_a^2 = 2b^2 + 2c^2 - a^2$ (Apollonius)
    \item \textbf{Angle Bisector ($l_a$):} $l_a = \frac{2bc}{b+c}\cos(A/2)$
    \item \textbf{Stewart's Theorem:} Cevian(a line from A to a) length $d$ dividing $a$ into $m, n$:
    \[ b^2m + c^2n = a(d^2 + mn) \]
    \item \textbf{Centroid ($G$):} Intersection of \textit{medians} (vertex to midpoint).
    \item \textbf{Incenter ($I$):} Intersection of \textit{angle bisectors}.
    \item \textbf{Circumcenter ($O$):} Intersection of \textit{perpendicular bisectors} of the sides.
    \item \textbf{Orthocenter ($H$):} Intersection of \textit{altitudes} (vertex perpendicular to opposite side).
\end{itemize}
\subsection{Circle}
\begin{itemize}[leftmargin=*]
    \item \textbf{General Eq:} $x^2+y^2+2gx+2fy+c=0$. Center $(-g, -f)$, $r=\sqrt{g^2+f^2-c}$.
    \item \textbf{Tangent at $P(x_1, y_1)$:} $xx_1 + yy_1 + g(x+x_1) + f(y+y_1) + c = 0$
    \item \textbf{Slope of tangent from $P(x, y)$:} $y + f = m(x + g) \pm r\sqrt{1 + m^2}$
    \item \textbf{Tangency Condition ($y=mx+c$):} $c^2 = r^2(1+m^2)$
    \item \textbf{Chord Length:} $L = 2\sqrt{r^2 - d^2}$ ($d = $ dist to center)
    \item \textbf{Power of Point $P$:} $d^2 - r^2$ ($d = $ dist $P$ to center) (Useful for determining if a point is inside or outside)
    \item \textbf{Secant Theorem:} $PA \cdot PB = PC \cdot PD$ (AB and CD are 2 lines drawn from point P over the circle)
    \item \textbf{Sector Area:} $\frac{1}{2}r^2\theta$ ($\theta$ in rads)
    \item \textbf{Segment Area (Cut by chord):} $\frac{1}{2}r^2(\theta - \sin\theta)$
\end{itemize}
\subsection{Points \& 2D Line}
\begin{itemize}[leftmargin=*]
    \item \textbf{Dist. point $(x_1, y_1)$ to line $ax+by+c=0$:} 
    \[ d = \frac{|ax_1+by_1+c|}{\sqrt{a^2+b^2}} \]
    \item \textbf{Dist. between parallel lines:} $d = \frac{|c_1-c_2|}{\sqrt{a^2+b^2}}$
    \item \textbf{Shoelace Area:} $\frac{1}{2} | \sum (x_i y_{i+1} - x_{i+1} y_i) |$
    \item \textbf{foot of Perpendicular $(h,k)$}:
    \[ \frac{h-x_1}{a} = \frac{k-y_1}{b} = -\frac{ax_1+by_1+c}{a^2+b^2} \]
    \item \textbf{Reflection $(h,k)$:} Replace $- (...)$ with $-2 (...)$ above.
    \item \textbf{Rotation of Axes (CCW $\theta$):}
    $x = X\cos\theta - Y\sin\theta, \quad y = X\sin\theta + Y\cos\theta$
\end{itemize}
\subsection{Vector}
\begin{itemize}[leftmargin=*]
    \item \textbf{Direction Cosines:} $l=\cos\alpha, m=\cos\beta, n=\cos\gamma$. $l^2+m^2+n^2=1$.
    \item \textbf{Dot Product:} $\mathbf{a} \cdot \mathbf{b} = |\mathbf{a}||\mathbf{b}|\cos\theta$
    \item \textbf{Projection of $\mathbf{a}$ on $\mathbf{b}$:} $(\mathbf{a} \cdot \hat{b})\hat{b}$
    \item \textbf{Cross Product:} Area $\triangle = \frac{1}{2}|\mathbf{a} \times \mathbf{b}|$.
    \item \textbf{Scalar Triple (Vol):} $[\mathbf{a}\mathbf{b}\mathbf{c}] = \mathbf{a} \cdot (\mathbf{b} \times \mathbf{c})$. Tet Vol $= \frac{1}{6}|[\dots]|$.
    \item \textbf{Vector Triple:} $\mathbf{a} \times (\mathbf{b} \times \mathbf{c}) = (\mathbf{a}\cdot\mathbf{c})\mathbf{b} - (\mathbf{a}\cdot\mathbf{b})\mathbf{c}$
\end{itemize}
\subsection{3D Line}
Line 1: $\mathbf{r} = \mathbf{a}_1 + \lambda\mathbf{b}_1$. Line 2: $\mathbf{r} = \mathbf{a}_2 + \mu\mathbf{b}_2$.
\begin{itemize}[leftmargin=*]
    \item \textbf{Angle Between Lines:} $\cos \theta = \frac{|\mathbf{b}_1 \cdot \mathbf{b}_2|}{|\mathbf{b}_1||\mathbf{b}_2|}$
    \item \textbf{Condition for Intersection:} Scalar Triple Product $[\mathbf{a}_2-\mathbf{a}_1, \mathbf{b}_1, \mathbf{b}_2] = 0$.
    \item \textbf{Intersection Point:} Solve $\mathbf{a}_1 + \lambda\mathbf{b}_1 = \mathbf{a}_2 + \mu\mathbf{b}_2$ for $\lambda, \mu$.
    \item \textbf{Shortest Dist (Skew):} $D = \left| \frac{(\mathbf{a}_2 - \mathbf{a}_1) \cdot (\mathbf{b}_1 \times \mathbf{b}_2)}{|\mathbf{b}_1 \times \mathbf{b}_2|} \right|$
    \item \textbf{Dist Between Parallel Lines:} $D = \frac{|(\mathbf{a}_2 - \mathbf{a}_1) \times \mathbf{b}|}{|\mathbf{b}|}$
    \item \textbf{Dist Point $P$ to Line:} $D = \frac{|(\mathbf{p} - \mathbf{a}) \times \mathbf{b}|}{|\mathbf{b}|}$
    \item \textbf{Foot of Perp $F$ from $P$:} $\mathbf{f} = \mathbf{a} + \left( \frac{(\mathbf{p}-\mathbf{a})\cdot \mathbf{b}}{|\mathbf{b}|^2} \right)\mathbf{b}$
    \item \textbf{Image $P'$ in Line:} $\mathbf{p'} = 2\mathbf{f} - \mathbf{p}$.
\end{itemize}
\subsection{Plane}
Through $\mathbf{a}$, normal $\mathbf{n}$. Eq: $(\mathbf{r} - \mathbf{a}) \cdot \mathbf{n} = 0$.
\begin{itemize}[leftmargin=*]
    \item \textbf{Cartesian:} $Ax + By + Cz + D = 0$. Normal $(A, B, C)$.
    \item \textbf{Dist Point $P$ to Plane:} $D = \frac{|Ax_1 + By_1 + Cz_1 + D|}{\sqrt{A^2+B^2+C^2}}$
    \item \textbf{Angle Between Planes:} $\cos \theta = \frac{|\mathbf{n}_1 \cdot \mathbf{n}_2|}{|\mathbf{n}_1||\mathbf{n}_2|}$
    \item \textbf{Angle Between Line \& Plane:} $\sin \phi = \frac{|\mathbf{b} \cdot \mathbf{n}|}{|\mathbf{b}||\mathbf{n}|}$
    \item \textbf{Intersection (Line of two planes):} Direction $\mathbf{v} = \mathbf{n}_1 \times \mathbf{n}_2$.
    \item \textbf{Family of Planes:} $P_1 + \lambda P_2 = 0$ (Planes through intersection line).
    \item \textbf{Foot of Perp from $P$:} $\frac{x-x_1}{A} = \frac{y-y_1}{B} = \frac{z-z_1}{C} = \frac{-(Ax_1+By_1+Cz_1+D)}{A^2+B^2+C^2}$
    \item \textbf{Image of Point in Plane:} Use $-2$ instead of $-1$ above.
    \item \textbf{Plane Bisectors:} $\frac{A_1x+B_1y+C_1z+D_1}{\sqrt{A_1^2+B_1^2+C_1^2}} = \pm \frac{A_2x+B_2y+C_2z+D_2}{\sqrt{A_2^2+B_2^2+C_2^2}}$
\end{itemize}
\subsection{Sphere}
\begin{itemize}[leftmargin=*]
    \item \textbf{Standard:} $(x-u)^2 + (y-v)^2 + (z-w)^2 = r^2$.
    \item \textbf{Diameter Form:} $(x-x_1)(x-x_2) + (y-y_1)(y-y_2) + (z-z_1)(z-z_2) = 0$.
    \item \textbf{Tangent Plane at $P$:} $(\mathbf{r}-\mathbf{c}) \cdot (\mathbf{p}-\mathbf{c}) = r^2$.
    \item \textbf{Circle of Intersection:} Radius $R = \sqrt{r^2 - d^2}$.
    \item \textbf{Spherical Cosine Rule:} $\cos a = \cos b \cos c + \sin b \sin c \cos A$.
\end{itemize}
\subsection{Solid}
\begin{itemize}[leftmargin=*]
    \item \textbf{Cone:} $V=\frac{1}{3}\pi r^2 h$, $SA = \pi r(r+l)$
    \item \textbf{Sphere:} $V=\frac{4}{3}\pi r^3$, $SA = 4\pi r^2$
    \item \textbf{Spherical Cap:} $SA = 2\pi r h$. $V = \frac{1}{3}\pi h^2(3r-h)$.
    \item \textbf{Frustum (Cone):} $V = \frac{1}{3}\pi h (R^2 + r^2 + Rr)$. $LSA = \pi(R+r)l$.
    \item \textbf{Pyramid:} $V = \frac{1}{3} (\text{Base Area}) \times h$. $LSA = \frac{1}{2} (\text{Perimeter}) \times l$.
    \item \textbf{Regular Tetrahedron:} Side $a$. $H = a\sqrt{2/3}$. $V = \frac{a^3}{6\sqrt{2}}$. Face Angle $\cos\theta = 1/3$.
\end{itemize}
\subsection{Trigonometry}
\begin{itemize}[leftmargin=*]
    \item $\sin(A\pm B) = \sin A\cos B \pm \cos A\sin B$
    \item $\cos(A\pm B) = \cos A\cos B \mp \sin A\sin B$
    \item $\cos 2A = \cos^2 A - \sin^2 A = 2\cos^2 A - 1$
    \item $\tan^{-1}x \pm \tan^{-1}y = \tan^{-1}\left( \frac{x \pm y}{1 \mp xy} \right)$
    \textit{(Note: For $+$, if $xy > 1$, add $\pi$ to the result.)}
    \item $2\tan^{-1}x = \tan^{-1}\left(\frac{2x}{1-x^2}\right) = \sin^{-1}\left(\frac{2x}{1+x^2}\right) = \cos^{-1}\left(\frac{1-x^2}{1+x^2}\right)$
    \item $\tan^{-1}x + \cot^{-1}x = \frac{\pi}{2}$
\end{itemize}
\end{multicols}\end{document}
